% concatenated from arxiv.tex
\documentclass[11pt,reqno]{amsart}

\usepackage{fullpage}
\usepackage{amsmath,amssymb,amsthm}
\usepackage{amsaddr}
\usepackage{amsfonts}
\usepackage[numbers,compress]{natbib}
\usepackage{mathtools}
\usepackage{mathrsfs}
\usepackage{enumitem}
\usepackage{bm}
\usepackage[hidelinks,bookmarks=false]{hyperref}
\usepackage{url}

\usepackage{subcaption}
\captionsetup{skip=4pt}
\captionsetup[subfigure]{skip=2pt}

\usepackage[table,xcdraw]{xcolor}
\usepackage{multirow}

\usepackage{booktabs,tabularx,threeparttable,makecell}
\usepackage{capt-of}

\usepackage{mathtools}
\mathtoolsset{showonlyrefs=true}

\renewcommand{\baselinestretch}{1.2}



\newtheorem{theorem}{Theorem}
\newtheorem{corollary}[theorem]{Corollary}
\newtheorem{lemma}[theorem]{Lemma}
\newtheorem{proposition}[theorem]{Proposition}
\newtheorem{external}[theorem]{External result}
\newtheorem{assumption}{Assumption}
\theoremstyle{definition}
\newtheorem{definition}{Definition}
\newtheorem{remark}{Remark}
\newtheorem{example}{Example}

\newcommand{\R}{\mathbb{R}}
\newcommand{\N}{\mathbb{N}}
\newcommand{\KL}{\mathrm{KL}}
\renewcommand{\d}{\mathrm{d}}
\renewcommand{\hat}{\widehat}
\renewcommand{\tilde}{\widetilde}
\newcommand{\Ep}{\mathbb{E}}
\newcommand{\mJ}{\mathcal{J}}
\newcommand{\mM}{\mathcal{M}}
\newcommand{\mP}{\mathcal{P}}
\newcommand{\mone}{\textbf{1}}
\newcommand{\norm}[1]{\lVert #1 \rVert}
\DeclareMathOperator*{\argmin}{arg min}
\DeclareMathOperator*{\argmax}{arg max}
\newcommand{\supp}{\mathrm{supp}}
\renewcommand{\epsilon}{\varepsilon}

\title{{\Large Minimax Optimal Estimation of Transport-Growth Pairs \\ in Unbalanced Optimal Transport}}



\usepackage{color}
\newcommand{\rc}{\color{red}}
\newcommand{\mc}{\color{magenta}}
\newcommand{\bc}{\color{blue}}
\newcommand{\blue}[1]{{\color{blue}#1}}
\DeclareMathOperator{\clip}{clip}



\allowdisplaybreaks
\sloppy
\raggedbottom

\let\origmaketitle\maketitle
\def\maketitle{
  \begingroup
  \def\uppercasenonmath##1{} % this disables uppercasing title
  \let\MakeUppercase\relax % this disables uppercasing authors
  \origmaketitle
  \endgroup
}

\author[M. Imaizumi and A. Coauthor]{Donlapark Ponnoprat$^{1}$, Noboru Isobe$^{2}$ \and 
Masaaki Imaizumi$^{2,3,4}$
}
\address{
$^{1}$ Chiang Mai University, Chiang Mai, Thailand \\
$^{2}$ RIKEN Center for Advanced Intelligence Project, Tokyo, Japan \\
$^{3}$ Graduate School of Arts and Science, The University of Tokyo, Tokyo, Japan \\
$^{4}$ Graduate School of Science, Kyoto University, Kyoto, Japan
}

\email{donlapark.p@cmu.ac.th}
\email{noboru.isobe@riken.jp}
\email{imaizumi@g.ecc.u-tokyo.ac.jp}


\begin{document}

\maketitle



\begin{abstract}
Unbalanced optimal transport (UOT) extends classical optimal transport to measures with different total masses, but statistical guarantees for Monge-type estimation remain limited. We study unbalanced transport with quadratic cost and Kullback-Leibler marginal penalties and argue that the natural population target is not a map alone, but a transport-growth pair. Consequently, we develop two estimators for the transport-growth pairs under several setups: an optimal transport plan-based estimator for a general case, and a kernel-based estimator for a case with smooth densities. We also show that an error of the estimator achieves the minimax optimal rate by deriving a matching lower bound of the minimax risk. Our main technical contribution is a value-based stability reduction that converts perturbations of the UOT objective into transport and growth risks through a UOT gap condition. These results provide a statistical foundation for Monge-type estimation in unbalanced optimal transport.
\end{abstract}


\section{Introduction}

\subsection{Background}

Optimal transport (OT) provides a principled geometric language for comparing probability distributions and has become a standard tool in statistics and machine learning. In statistical applications, however, the underlying population measures are rarely observed directly; instead, one must infer transport objects from finite samples. For balanced OT, this question is now supported by a substantial statistical theory. Smooth transport maps admit minimax analysis via semi-dual curvature and growth arguments \cite{hutter2021minimax}, while plug-in, barycentric-projection, and entropic estimators have been analyzed in discrete, semi-discrete, and smooth regimes \cite{deb2021barycentric,Pooladian2021,pooladian2023semidiscrete,manole2024plugin}. More recently, the balanced theory has expanded to general function-space analyses and sharper stability reductions for plug-in estimators \cite{divol2025general,balakrishnan2025stability}.

Many datasets are not naturally balanced: the total masses may differ, unmatched observations may be present, and mass creation or destruction may be an intrinsic feature of the phenomenon under study. Unbalanced optimal transport (UOT) addresses this issue by relaxing the hard marginal constraints of OT and penalizing deviations of the plan marginals from the reference measures. This viewpoint underlies the modern entropy-transport framework and includes important models such as logarithmic entropy-transport and the Wasserstein-Fisher-Rao/Hellinger-Kantorovich geometry \cite{ChizatEtAl2018,LieroMielkeSavare2018,savare2024relaxation,gallouet2025regularity}. It has also generated a large computational and applied literature, including Sinkhorn-type solvers and scalable parameterizations \cite{pham2020sinkhorn,sejourne2022faster,Gazdieva2024,yang2018scalable}, with applications ranging from single-cell dynamics to growth modeling and shape analysis \cite{schiebinger2019,sha2023cell,dimarino2020,bauer2022SRNF}.

Despite this rapid development, the statistical theory of UOT remains less developed, especially for Monge-type objects. For the UOT model studied below, \cite{vacher2022stability,vacher2023semi} derived semi-dual formulations and established global Bregman-type stability for the corresponding objectives, yielding the first fast statistical rates for UOT semi-dual potentials. More recently, statistical properties of unbalanced Kantorovich-Rubinstein quantities have been analyzed for finitely supported measures and for spatio-temporal point-process models \cite{HundrieserEtAl2025,struleva2025sharp}. However, compared with the balanced case, a plug-in or near-minimax theory for Monge-type UOT maps is still largely missing.

A key difficulty is structural. In balanced OT, the specific form of the transport objective allows us to bound the transport error by exploiting the semi-dual form of the transport problem \cite{hutter2021minimax,manole2024plugin,balakrishnan2025stability}. UOT, by contrast, has a different objective where the hard marginal constraints are replaced with divergence penalties; this modification introduces a \emph{growth map} that specified how the masses will contract or expand after the transport.
Consequently, a rigorous statistical theory for UOT must account for the errors of transport map and the growth map estimation.

\subsection{Our contribution}

In this paper, we develop a statistical theory for Monge-type estimation in unbalanced optimal transport with KL marginal penalties. In particular, we develop two estimators for the transport-growth pair, then study their estimation error. 

Concretely, our contributions are as follows.
\begin{itemize}
    \item \textbf{Estimators for the transport-growth pair: }We formulate the Monge-type statistical target for UOT as the transport-growth pair, then develop two estimators: a plan-based estimator obtained from a discrete UOT plan, and a smooth plug-in estimator based on regularized marginal estimates. 
    \item \textbf{Prove optimality: } We prove that the estimators achieve the minimax optimal convergence rate, by deriving both an upper bound of their estimation error and a corresponding lower bound of the minimax risk. This result improves upon the existing evaluation, as summarized in Table \ref{tab:main-rate-comparison-highd}.
    \item \textbf{Proof technique by stability bound:} To analyze this target, we establish a stability bound and a first-order expansion of the UOT objective under perturbations of the marginals, which isolates the UOT-specific analytic step from the choice of marginal estimator.
    By employing a stability-based proof, we circumvent the limitations of proofs using semi-dual potentials and enable the evaluation of optimality.
\end{itemize}

\begin{minipage}{\textwidth}
\footnotesize
\centering
\captionof{table}[High-dimensional comparison of statistical rates]{Comparison of convergence rates of the estimation error. 
\(d\ge 5\) is the dimension of samples, \(N = \min\{n,m\}\) where $n$ and $m$ are samples sizes from the measures, and \(\widetilde O(\cdot)\) hides polylogarithmic factors, and $\alpha$ is smoothness of densities of the measure.
For \cite{vacher2022stability}, 
their regularity argument yields
\((\alpha+2)\)-smooth semi-dual potentials with \(\alpha+2<d/2\).
}
\label{tab:main-rate-comparison-highd}
\begin{tabular}{
@{}
>{\raggedright\arraybackslash}p{2.0cm}
>{\centering\arraybackslash}p{2.55cm}
>{\centering\arraybackslash}p{1.7cm}
>{\centering\arraybackslash}p{1.5cm}
>{\centering\arraybackslash}p{3.0cm}
>{\centering\arraybackslash}p{1.1cm}
@{}
}
\toprule
Reference & Setup & \makecell[c]{transport\\estimation} & \makecell[c]{growth\\ estimation} & Convergence rate & \makecell[c]{Lower\\bound} \\
\midrule \midrule

\makecell[l]{ \cite{hutter2021minimax} \\ \cite{manole2024plugin}}
&
\makecell[c]{ balanced \\ (w/ smoothness)}
&
\(\checkmark\)
&
N/A
&
\(\widetilde O \left(N^{-2\alpha/(2\alpha-2+d)}\right)\)
&
\(\checkmark\)
\\ \hline

\addlinespace[2pt]

\makecell[l]{\cite{vacher2022stability}}
&
\makecell[c]{unbalanced \\ (w/ smoothness)}
&
\(\checkmark\)
&
&
\(O \left(N^{-(\alpha+2)/d}\right)\)
&
\\

\addlinespace[2pt]

\makecell[l]{\cite{vacher2023semi}}
&
unbalanced
&
\(\checkmark\)
&
&
\(O \left(N^{-2/d}\right)\)
&
\\

\addlinespace[2pt]

\makecell[l]{\textbf{This paper}\\ \textbf{(plan-based)}}
&
unbalanced
&
\(\checkmark\)
&
\(\checkmark\)
&
\(\widetilde O(N^{-2/d})\)
&
\\

\addlinespace[2pt]

\makecell[l]{\textbf{This paper}\\ \textbf{(kernel-based)}}
&
\makecell[c]{unbalanced \\ (w/ smoothness)}
&
\(\checkmark\)
&
\(\checkmark\)
&
\(\widetilde O \left(N^{-2\alpha/(2\alpha-2+d)}\right)\)
&
\(\checkmark\)
\\

\bottomrule
\end{tabular}
\end{minipage}


\subsection{Notation}
Throughout, $\Omega\subset\R^d$ denotes the ambient domain. We write $C(\Omega)$ for the space of continuous real-valued functions on $\Omega$, $C_b(\Omega)$ for the bounded continuous functions, and $\mathrm{Diff}(\Omega)$ for the class of diffeomorphisms from $\Omega$ onto itself whenever differentiability is imposed. The set of finite positive Radon measures on $\Omega$ is denoted by $\mM_+(\Omega)$, and $\mP(\Omega)$ denotes the Borel probability measures on $\Omega$. For a measurable map $T:\Omega\to\Omega$ and a measure $\mu\in\mM_+(\Omega)$, $T_{\#}\mu$ denotes the pushforward of $\mu$. For $\gamma\in\mM_+(\Omega\times\Omega)$, we write $\gamma_0$ and $\gamma_1$ for its first and second marginals, respectively. If $\eta,\sigma\in\mM_+(\Omega)$ have the same total mass, then $\Pi(\eta,\sigma)$ denotes the set of couplings between them. We use $\|\cdot\|$ for the Euclidean norm, $\|\cdot\|_{\mathrm{op}}$ for the operator norm, $\mathbf{1}\{\cdot\}$ for the indicator function, and $D_\KL$ for the KL divergence. The Legendre-Fenchel transform of a convex function $F$ is denoted by $F^*$. 
For real numbers $a<b$, the clipping function is $\clip_{[a,b]}(t)\coloneqq\min\{b,\max\{a,t\}\}$.

\section{Unbalanced Optimal Transport with Quadratic Cost} \label{sec:setup}

We briefly review the unbalanced optimal transport in the \textit{Gaussian-Hellinger} case. For a complete treatment of the general case, see \cite{LieroMielkeSavare2018,savare2024relaxation,gallouet2025regularity}.
Let $\Omega\subset\R^d$ be a bounded convex domain and $\mu,\nu\in\mM_+(\Omega)$ are finite positive Radon measures. The \emph{c-transform} of a function $\varphi:\Omega \to \R$ is defined as $\varphi^c(x) =\inf_y \{ \|x-y\|^2/2 - \varphi(y) \}$. We say that $\varphi$ is \emph{c-concave} if $\varphi_0=\psi^c$ for some $\psi: \Omega \to \R$.



\subsection{Transport formulations}

\paragraph{Monge problem}: We introduce a Monge-like form of the unbalanced optimal transport setup.
A \textit{transport-growth pair} consists of a measurable map $T:\Omega\to\Omega$ and a measurable weight $\lambda:\Omega\to[0,\infty)$, and acts on $\mu$ through the weighted pushforward
\begin{align}
    (T,\lambda)_{\#}\mu \coloneqq T_{\#}(\lambda^2\mu).
\end{align}
The map $\lambda$ is referred to as the \emph{growth factor}, since it rescales mass after the transport.
In the {Gaussian-Hellinger} case, we have the following cone cost by 
$$
    C((x,r),(y,s))^2
    \coloneqq
    r^2 + s^2 - 2rs e^{-\|x-y\|^2/4}
    ,\qquad x,y\in\Omega,\ r,s\ge 0.
$$
Using this cone cost, we obtain the 
\emph{unbalanced Monge problem} for the Gaussian-Hellinger case:
\begin{align}
\label{eq:gh_monge}
    \mathrm{UM}(\mu,\nu)
    \coloneqq
    \inf_{(T,\lambda): (T,\lambda)_{\#}\mu=\nu}
    \int_\Omega
    \bigl(
        1 + \lambda(x)^2 - 2\lambda(x) e^{-\|x-T(x)\|^2/4}
    \bigr) d\mu(x).
\end{align}
In contrast to balanced OT, where only a transport map $T_0$ appears, UOT also requires a growth map $\lambda_0$ to reconcile the discrepancy between the total masses of $\mu$ and $\nu$.


\paragraph{Kantorovich problem}: We consider another problem for the unbalanced optimal transport by optimizing over plans $\gamma\in\mM_+(\Omega\times\Omega)$ instead of deterministic transport-growth pairs for the Monge problem. Specifically, the Kantorovich problem associated to \eqref{eq:gh_monge} is given by
\begin{align}
\label{eq:gh_kantorovich}
    \mathrm{UOT}(\mu,\nu)
    \coloneqq
    \inf_{\gamma'\in\mM_+(\Omega\times\Omega)}
    \left\{
        \int_{\Omega\times\Omega} \frac{\|x-y\|^2}{2} d\gamma'(x,y)
        + D_\KL(\gamma_0'\mid\mu)
        + D_\KL(\gamma_1'\mid\nu)
    \right\}.
\end{align}
The corresponding dual problem reads:
\begin{align}
\label{eq:gh_dual}
    \mathrm{UOT}(\mu,\nu)
    =
    \sup_{\substack{(\varphi,\psi)\in C_b(\Omega)^2\\ \varphi(x)+\psi(y)\le \|x-y\|^2/2}}
    \left\{
        \int_\Omega \bigl(1-e^{-\varphi(x)}\bigr) d\mu(x)
        + \int_\Omega \bigl(1-e^{-\psi(y)}\bigr) d\nu(y)
    \right\}.
\end{align}


Let $(\varphi_0,\psi_0)$ denote an optimal dual pair for \eqref{eq:gh_dual}, and we refer to them as \emph{potentials}.  Existence of the solutions of \eqref{eq:gh_monge} and \eqref{eq:gh_dual} are guaranteed by the following theorem:
\begin{theorem}[{\cite{gallouet2025regularity}}]
\label{prop:monge-potential}
Let $\Omega\subset\R^d$ be bounded and convex, let $\mu=\rho_0dx$ with $\rho_0>0$ a.e., and let $\nu\ll dx$ be supported in $\Omega$. Then there exists a $\mu$-a.e. unique $c$-concave $\varphi_0$ such that $(\varphi_0,\psi_0)$ solves the dual problem \eqref{eq:gh_dual} with $\psi_0 = \varphi^c_0$, and $\varphi_0$ uniquely determines a solution $(T_0,\lambda_0)$ of the Monge problem \eqref{eq:gh_monge} via:
\begin{align}\label{eq:unique}
    T_0(x)=x-\nabla \varphi_0(x),
    \qquad
    \lambda_0(x)=\exp\left(-\varphi_0(x)+\frac1{4}\|\nabla \varphi_0(x)\|^2\right).
\end{align}
\end{theorem}
\paragraph{Useful form:} Following Theorem~\ref{prop:monge-potential}, we define \emph{active source marginal} $\gamma_0=e^{-\varphi_0}\mu$ and \emph{active-source factor} $a_0(x)\coloneqq e^{-\varphi_0(x)/2}$.
It follows from the complementary slackness that $\varphi_0(x) + \psi_0(T_0(x)) = \| x - T_0(x) \|^2/2$. We thus have the following gradient-free formulation of $T_0$ and $\lambda_0$:
\begin{align} \label{eq:convenient}
    T_0(x)
    &\in
    \argmin_{y\in\Omega}
    \left\{
        \frac{\|x-y\|^2}{2}-\psi_0(y)
    \right\},
    \qquad
    \lambda_0(x)
    =a_0(x)^2
    \exp\left(\frac14\|x-T_0(x)\|^2\right) .
\end{align}







\subsection{Statistical estimation problem of the Monge map}

Given finite positive measures $\mu,\nu \in \mM_+(\Omega)$, we consider total masses
$    M_\mu \coloneqq \mu(\Omega)$ and $    M_\nu \coloneqq \nu(\Omega)$
and define the normalized probability measures $
    \bar\mu \coloneqq {\mu} / {M_\mu},
    \bar\nu \coloneqq {\nu} / {M_\nu}$.
Given sample sizes $n,m\in\N$, we observe mutually independent samples $X_1,\dots,X_n \sim \bar\mu$ and $Y_1,\dots,Y_m \sim \bar\nu$.
Because the data are sampled from the normalized laws $\bar\mu$ and $\bar\nu$, the total masses $M_\mu$ and $M_\nu$ are not identifiable from the samples alone. We therefore assume either that these masses are known, or that additional estimators $\hat M_\mu$ and $\hat M_\nu$ are available from external information.
We then define atomic weights
    $\hat\mu_i \coloneqq \hat\mu_n(\{X_i\}) = \hat M_\mu/n, \hat\nu_j \coloneqq \hat\nu_m(\{Y_j\}) = \hat M_\nu/m$ and the weighted empirical measures
\begin{align}
    \hat\mu_n
    &\coloneqq
    \frac{\hat M_\mu}{n}\sum_{i=1}^n \delta_{X_i},\qquad 
    \hat\nu_m
    \coloneqq
    \frac{\hat M_\nu}{m}\sum_{j=1}^m \delta_{Y_j}.
\end{align} 

    
Our objective is to estimate the population transport-growth pair $(T_0,\lambda_0)$ associated with $(\mu,\nu)$ from the observed samples and the mass estimators. Our proofs also estimate the auxiliary active-source factor $a_0=e^{-\varphi_0/2}$, since it is the quantity directly encoded by the source marginal $\gamma_0=a_0^2\mu$ of the Kantorovich plan.

\begin{remark}


The construction of $\hat M_\mu$ and $\hat M_\nu$ depends on the observation scheme. 
If external total-mass measurements are available, they can be plugged in directly. In Poisson or more general point-process models, the observed counts naturally carry mass information; see Appendix~\ref{sec:appendix_ppp} and the recent UOT analyses of \cite{HundrieserEtAl2025,struleva2025sharp}. 
\end{remark}

\section{Estimator design}

We propose two estimators: (i) a plan-based estimator and (ii) a kernel plugin estimator; the former can handle high dimensional data, while the latter is adaptive to the smoothness of the densities.

We introduce some notations: for a nonnegative matrix $G=(G_{ij})\in\R_+^{n\times m}$, define the discrete marginals $(G_0)_i \coloneqq \sum_{j=1}^m G_{ij}$ and $(G_1)_j \coloneqq \sum_{i=1}^n G_{ij}$.
With $F(r)=r\log r-r+1$, the corresponding discrete KL penalty between mass vectors is
$
    D_\KL(G_0\mid\hat\mu_n)
    \coloneqq
    \sum_{i=1}^n \hat\mu_i  F\bigl((G_0)_i / \hat\mu_i\bigr).
$

\subsection{Plan-based estimator} \label{sec:plan}



This estimator is constructed by solving the Kantorovich problem under the empirical measures $\hat{\mu}_n$ and $\hat{\nu}_m$ and using its optimal transport plan obtained as a matrix, based on the following steps.

\paragraph{(i) Transport plan estimation:}
Let the discrete cost matrix be $C_{ij} \coloneqq \|X_i - Y_j \|^2/2$. We define an estimator $\hat\gamma$ of the unbalanced optimal transport plan $\gamma$. For the Gaussian-Hellinger Kantorovich problem \eqref{eq:gh_kantorovich}, one may take any solution of
\begin{equation}
\hat\gamma\in\argmin_{\gamma'\in\R_+^{n\times m}}
\left\{
\sum_{i=1}^n\sum_{j=1}^m C_{ij}\gamma_{ij}'
+D_\KL(\gamma_0'\mid\hat\mu_n)+D_\KL(\gamma_1'\mid\hat\nu_m)
\right\}.
\end{equation}

\paragraph{(ii) Estimate discrete transport:}
Define the row masses $\hat r_i \coloneqq \sum_{j=1}^m \hat\gamma_{ij}$, $i=1,\dots,n$. We then define a point estimator of the transported component by the Fr\'echet projection: if $\hat r_i > 0$, 
\begin{equation}\label{eq:disc-frechet}
\hat T_i = \argmin_{y\in\Omega}\ \sum_{j=1}^m \hat\gamma_{ij} c(y,Y_j) = \frac{1}{\hat r_i}\sum_{j=1}^m \hat\gamma_{ij} Y_j \in \R^d
\end{equation}
If $\hat r_i = 0$, then the above objective is identically zero and $\hat T_i$ is not identifiable from the plan. In this case, we set $\hat T_i \coloneqq X_i$ by convention. 
Motivated by the active-marginal identity $\gamma_0=a_0^2\mu$ from Section~\ref{sec:setup}, we define an estimator $\hat a_i \coloneqq \sqrt{{\hat r_i}/{\hat\mu_i}}$ for the active-source factor $a_0(X_i)$.
To estimate the growth factor $\lambda_0(X_i)$, we fix constants $0<w_-<w_+$ containing the range of $w_0=e^{-\varphi_0}$ and set
\begin{align}
\label{eq:disc-gh-growth}
    \hat\lambda_i
    \coloneqq
    \clip_{[w_-,w_+]}(\hat a_i^2)
    \exp\left(\frac14\|X_i-\hat T_i\|^2\right).
\end{align}


\paragraph{(iii) Estimate Monge/growth map}

Since $(\hat T_i,\hat a_i,\hat\lambda_i)$ are defined only at the sample points $X_1,\dots,X_n$, we extend them to all of $\Omega$.
Here, we employ an approach of the nearest-neighbour method. Let $(V_i)_{i=1}^n$ be the Voronoi partition induced by $X_1,\dots,X_n$:
    $V_i \coloneqq \{x\in\Omega:\ \|x-X_i\|\le \|x-X_k\|\ \forall k\neq i\}$.
Setting $w_i(x) \coloneqq \mathbf{1}\{x\in V_i\}$ yields the piecewise-constant estimators
\begin{equation}
\label{eq:disc-1nn-map}
    \hat T^{\mathrm{1NN}}(x)
    =
    \sum_{i=1}^n \mathbf{1}\{x\in V_i\} \hat T_i,
    ~~
    \hat a^{\mathrm{1NN}}(x)
    =
    \sum_{i=1}^n \mathbf{1}\{x\in V_i\}\hat a_i,
    ~~
    \hat\lambda^{\mathrm{1NN}}(x)
    =
    \sum_{i=1}^n \mathbf{1}\{x\in V_i\}\hat\lambda_i,
\end{equation}
where $\hat a^{\mathrm{1NN}}$ and $\hat\lambda^{\mathrm{1NN}}$ estimate the active-source factor $a_0$ and the growth map $\lambda_0$, respectively.

In addition to the nearest-neighbor method, a Nadaraya-Watson-type estimator  can also be considered; see Definition~\ref{def:nw_extension_app} in Appendix~\ref{sec:appendix_nw_extension}.




\subsection{Kernel-based estimator} \label{sec:kernel}

We next present a kernel-based estimator for the transport-growth pair on the hypercube $\Omega=[0,1]^d$.
This estimator is based on density estimates of the measures $\mu,\nu$ and has the advantage of adapting to the smoothness of these densities.



\paragraph{(i) Prepare kernels}:
We introduce a new kernel function based on a cosine basis for estimating density functions on $[0,1]^d$. This method makes statistical use \cite{efromovich2010orthogonal,tsybakov2009nonparametric} of the eigenfunctions of the Neumann Laplacian \cite{strang1999dct}. This kernel is useful to relax the constraint that a density function must lie on a torus, a requirement in the balanced case \cite{manole2024plugin}.



We define the kernel function. In preparation, we define a function $\eta(s)  = e^{-1/s} \mathbf{1}\{s > 0\}$ for $s \in \R$, and $\tau(t)
    \coloneqq
    {\eta(2-t)} / ({\eta(2-t)+\eta(t-1)})$ for $t\ge 0$.
Then, $\tau\in C^\infty([0,\infty))$, $0\le \tau\le 1$, $\tau(t)=1$ for $t\in[0,1]$, and $\tau(t)=0$ for $t\ge 2$.
With $c_0 \coloneqq 1,
    c_{\ell} \coloneqq \sqrt{2}\ \ (\ell\ge 1)$, 
and for $L\ge 1$, define the one-dimensional Neumann kernel
\begin{align}
\label{eq:uot_cube_kernel_1d}
    \kappa_L(u,v)
    \coloneqq
    1+\sum_{\ell=1}^{\infty}
    \tau \left(\frac{\pi^2 \ell^2}{L^2}\right)c_\ell^2
    \cos(\pi \ell u)\cos(\pi \ell v),
    \qquad
    u,v\in[0,1].
\end{align}
Because $\tau$ has compact support, the sum in \eqref{eq:uot_cube_kernel_1d} is finite for each $L$. We then define the boundary-adapted separable kernels
\begin{align}
\label{eq:uot_cube_kernel}
    K_L(x,y)
    &\coloneqq
    \prod_{r=1}^d \kappa_L(x_r,y_r),
    \qquad x,y\in[0,1]^d.
\end{align}
Since each one-dimensional factor contains only $O(L)$ nonzero cosine modes, one evaluation of $K_L(x,y)$ costs $O(dL)$ arithmetic operations. Moreover, we have $\int_{[0,1]^d} K_L(x,y) dy = 1$ for $ x\in[0,1]^d$.
This is the cube analogue of a smooth spectral cutoff, but expressed in a form that is computationally tractable in moderate and high dimension: the cosine basis supplies the boundary adaptation, while the coordinatewise multiplier supplies exact separability.


\paragraph{(ii) Estimate measures via densities}:
We estimate the normalized densities first and then attach either the true masses (for the oracle equal-mass objects used in the analysis) or the estimated masses (for the actual fitted measures).
In particular, we define the preliminary kernel density estimators
\begin{align}
    \tilde p_n^{\mathrm{ker}}(\cdot)
    &\coloneqq
    \frac{1}{n}\sum_{i=1}^n K_{L_n}(x,X_i),
    \qquad
    \tilde q_m^{\mathrm{ker}}(\cdot)
    \coloneqq
    \frac{1}{m}\sum_{j=1}^m K_{L_m}(y,Y_j).
\end{align}
We then define estimator for the measures $\mu,\nu$ with the estimated total mass as
\begin{align}
    \hat \mu_n^{\mathrm{ker}}(x)
    &\coloneqq
    \frac{\hat{M}_\mu (\tilde p_n^{\mathrm{ker}}(x))_+ dx}
    {\int_{[0,1]^d}(\tilde p_n^{\mathrm{ker}}(u))_+ du},
    \qquad
    \hat \nu_m^{\mathrm{ker}}(y)
    \coloneqq
    \frac{ \hat{M}_\nu(\tilde q_m^{\mathrm{ker}}(y))_+ dy}
    {\int_{[0,1]^d}(\tilde q_m^{\mathrm{ker}}(v))_+ dv}.
\end{align}



\paragraph{(iii) Estimate maps}:
Let $(\hat\varphi_{nm}^{\mathrm{ker}},\hat\psi_{nm}^{\mathrm{ker}})$ be an optimal dual pair for
$\mathrm{UOT}(\hat\mu_n^{\mathrm{ker}},\hat\nu_m^{\mathrm{ker}})$. We define the estimators of the active-source factor as $\hat a_{nm}^{\mathrm{ker}}(x)
    \coloneqq
    e^{-\hat\varphi_{nm}^{\mathrm{ker}}(x)/2}$.
    Then, we define estimators for the transport map and the growth map by
\begin{align} \label{eq:kernel_estimator}
    \hat T_{nm}^{\mathrm{ker}}(x)
    &\in
    \argmin_{y\in[0,1]^d}
    \left\{
        \|x-y\|^2/2 - \hat\psi_{nm}^{\mathrm{ker}}(y)
    \right\},\\
    \hat\lambda_{nm}^{\mathrm{ker}}(x)
    &\coloneqq
     \clip_{[w_-,w_+]}(\hat a_{nm}^{\mathrm{ker}}(x)^2)
    \exp\left(\frac14\|x-\hat T_{nm}^{\mathrm{ker}}(x)\|^2\right).
\end{align}


\section{Minimax optimal rate}


\subsection{Basic assumptions}


\begin{assumption}\label{assm:curvature}
There exists a constant $\kappa \in (0, 1)$ such that 
\[
    \left(1 - \frac{1}{\kappa}\right)I \preceq \nabla^2 \varphi_0(x) \preceq (1 - \kappa)I, \qquad x \in \Omega.
\]
\end{assumption}
The bounds in Assumption~\ref{assm:curvature} are equivalent to $\kappa I \preceq \nabla^2 f_0(x) \preceq \kappa^{-1} I$ for the Brenier potential $f_0(x) \coloneqq \frac{1}{2}\|x\|^2 - \varphi_0(x)$ commonly used in convergence analysis of balanced OT map estimators \cite{hutter2021minimax,manole2024plugin,divol2025general}. The following consequence will be used regularly in our proofs:
\begin{lemma}
\label{lem:uot_gap_sufficient}
Suppose that Assumption~\ref{assm:curvature} holds. Then the dual potentials $(\varphi_0,\psi_0)$ satisfy:
\begin{align}
    \frac{1}{2}\|x-y\|^2-\varphi_0(x)-\psi_0(y)
    \ge \frac{\kappa}{2}\|y-T_0(x)\|^2,
    \qquad x,y\in\Omega.
\end{align}
\end{lemma}



\begin{assumption}
\label{assm:supp_global}
The domain $\Omega \subset \mathbb{R}^d$ is compact, convex, and satisfies the interior cone condition: there exist $\epsilon_0, \delta_0 > 0$ such that for all $x \in \Omega$ and $\epsilon \in (0, \epsilon_0)$,
$
    \mathrm{Vol}(B(x,\epsilon) \cap \Omega) \geq \delta_0 \mathrm{Vol}(B(x,\epsilon)).
$
\end{assumption}



\begin{assumption}
\label{assm:positivity}
The positive finite measures $\mu$ and $\nu$ have total masses $M_\mu$ and $M_\nu$. Their normalized laws $\bar\mu=\mu/M_\mu$ and $\bar\nu=\nu/M_\nu$ admit Lebesgue densities $p$ and $q$ bounded away from zero and infinity on $\Omega$: $0<\beta_{\min}\le p(x),q(x)\le\beta_{\max}<\infty$. Equivalently, the finite-measure densities are $\rho_\mu=M_\mu p$ and $\rho_\nu=M_\nu q$.
\end{assumption}
The two assumptions are standard in the statistical literature on transport map estimation: similar assumptions on the densities, on the transport potential or on the map itself appear in the balanced case \cite{hutter2021minimax,manole2024plugin,balakrishnan2025stability} and the unbalanced case \cite{gallouet2025regularity}.




\begin{assumption}\label{assm:masses}
There exist positive real sequences $a_n, b_m \to 0$ as $n,m \to \infty$ and a constant $c\in(0,1]$ such that
$
    \hat M_\mu \to M_\mu,
    \hat M_\nu \to M_\nu,
$ and
$\hat M_\mu \ge c M_\mu$, $\hat M_\nu \ge c M_\nu$, almost surely,
and $\Ep[|\hat M_\mu-M_\mu|] \le a_n$, $\Ep[|\hat M_\nu-M_\nu|] \le b_m$ hold for all sufficiently large $n,m$.
\end{assumption}
This assumption is intentionally modular rather than standard in the UOT setup. 
The condition itself is mild and it is satisfied, for example, when the masses are known, when separate total-mass measurements are available, or in point-process/counting models where the counts carry mass information \cite{HundrieserEtAl2025,struleva2025sharp}. For that reason we state it abstractly instead of tying the main theory to a single data-acquisition mechanism.



\subsection{Plan-based estimator}

We study the estimation error of the $1$NN estimator \eqref{eq:disc-1nn-map}. The transport loss is measured under the \emph{active source measure} $\gamma_0 = a_0^2 \mu$, which is the natural source marginal in the Kantorovich problem; the growth loss is measured for $\lambda_0$ under the empirical source measure $\hat\mu_n$.
In preparation, we define the following value:
\begin{align}
    \mathfrak R_n^{\mathrm{emp}}(d)
    \coloneqq
    \begin{cases}
        n^{-1/2}, & d \le 3,\\
        (\log n)n^{-1/2}, & d = 4,\\
        n^{-2/d}, & d \ge 5,
    \end{cases}
\end{align}
which corresponds to the convergence rate of empirical distribution in the Wasserstein distance $W_2$ \cite{weed2019sharp}.
Then, we obtain the following rate: 

\begin{theorem}[Error rate of plan-based estimator]
\label{thm:main_plan_based_rates}
Assume that Assumptions~\ref{assm:curvature}, \ref{assm:supp_global}, \ref{assm:positivity} and \ref{assm:masses} hold.
Let $(\hat T^{\mathrm{1NN}},\hat a^{\mathrm{1NN}},\hat\lambda^{\mathrm{1NN}})$ be the $1$NN estimators defined in Section~\ref{sec:plan}, where $\hat\lambda^{\mathrm{1NN}}$ is the clipped growth estimator in \eqref{eq:disc-gh-growth}. Let $a_n,b_m$ be the sequences from Assumption~\ref{assm:masses}.
Then there exists a constant $C>0$ such that, for all sufficiently large $n,m$,
\begin{align}
\label{eq:section4_plan_map_rate_abstract}
    &\max\left\{\Ep\Bigl[
        \int_\Omega \|\hat T^{\mathrm{1NN}}(x)-T_0(x)\|^2  d\mu(x)
    \Bigr], \Ep\Bigl[
        \int_\Omega |\hat\lambda^{\mathrm{1NN}}(x)-\lambda_0(x)|^2  d\mu(x)
    \Bigr]\right\} \\
    &\qquad \le
    C\log n \Bigl(
        M_\mu \mathfrak R_n^{\mathrm{emp}}(d)
        + M_\nu \mathfrak R_m^{\mathrm{emp}}(d)
        + a_n+b_m
    \Bigr).
\end{align}
\end{theorem}

This result shows that the estimation error of an unbalanced mapping by 1NN, when excluding the effects of mass estimation $a_n$ and $b_m$, is equal to the rate of convergence of the empirical distribution in the $W_2$. 
Furthermore, since $a_n$ and $b_m$ achieve the parametric rate in many cases, they do not hamper the convergence rate of the empirical distribution.






\subsection{Kernel-based estimator on the hypercube}

We study the estimation error of the kernel estimator \eqref{eq:kernel_estimator} on $[0,1]^d$. Here, we additionally introduce an assumption on the smoothness of the normalized densities of $\mu,\nu$.
\begin{definition}[Neumann-compatible H\"older class]
\label{def:neumann_holder_class}
Let $\mathbb{T}_2^d \coloneqq (\mathbb{R}/2\mathbb{Z})^d$. For $f:[0,1]^d \to \mathbb{R}$, define its even $2$-periodic reflection $\mathcal{E}f: \mathbb{T}_2^d \to \mathbb{R}$ by
$
    (\mathcal{E}f)(x_1,\dots,x_d)
    \coloneqq
    f\bigl(\vartheta(x_1),\dots,\vartheta(x_d)\bigr)
$, where $\vartheta(t)
    \coloneqq
    \min_{m\in\mathbb{Z}} |t-2m| \in [0,1]
$.
For $s>0$ and $M>0$, we write
\begin{align}
    \mathcal{C}_N^s([0,1]^d;M)
    \coloneqq
    \left\{
        f:[0,1]^d\to\mathbb{R}:
        \|\mathcal{E}f\|_{C^s(\mathbb{T}_2^d)} \le M
    \right\},
\end{align}
which consists of functions whose coordinatewise even reflection is $C^s$-smooth on the doubled torus. 
\end{definition}
\begin{assumption}[Smooth density]
\label{assm:cube_kernel_smoothness}
Let $\mathcal{C}_N^s([0,1]^d;M)$ be the Neumann-compatible H\"older class.
For some $\alpha>1$ and $M>0$, the normalized densities satisfy $p,q \in \mathcal{C}_N^{\alpha-1}([0,1]^d;M)$.
\end{assumption}
In preparation, we define a benchmark rate as follows, then obtain the result on the estimation error:
\begin{align}
    \mathfrak R_n^{\mathrm{ker}}(\alpha)
    \coloneqq
    \begin{cases}
        n^{-1}, & d=1,\\
        (\log n)n^{-1}, & d=2,\\
        n^{-{2\alpha} / ({2(\alpha-1)+d})}, & d\ge 3
    \end{cases}.
\end{align}



\begin{theorem}[Error rate of kernel estimator]
\label{thm:main_cube_kernel_rates}
Assume that $\Omega=[0,1]^d$, and Assumptions~\ref{assm:curvature}, \ref{assm:positivity} and \ref{assm:masses}, and \ref{assm:cube_kernel_smoothness} hold.
Assume that $L_n \asymp n^{1/(d+2(\alpha-1))}$ and $L_m \asymp m^{1/(d+2(\alpha-1))}$, and let $a_n,b_m$ be the sequences from Assumption~\ref{assm:masses}.
Then there exists a constant $C>0$, depending only on $d,M,\beta_{\min},\beta_{\max}$, the constants in Assumption~\ref{assm:curvature}, and the cutoff $\tau$, such that
\begin{align}
    &\max\left\{\Ep\left[
        \int_{[0,1]^d}
        \|\hat T_{nm}^{\mathrm{ker}}(x)-T_0(x)\|^2 d\mu(x)
    \right],
    \Ep\left[
        \int_{[0,1]^d}
        |\hat\lambda_{nm}^{\mathrm{ker}}(x)-\lambda_0(x)|^2 d\mu(x)
    \right] \right\} \\
    &\qquad \le
    C\Bigl(
        M_\mu \mathfrak R_n^{\mathrm{ker}}(\alpha)
        +
        M_\nu \mathfrak R_m^{\mathrm{ker}}(\alpha)
        + a_n+b_m
    \Bigr).\label{eq:section4_cube_kernel_map_rate}
\end{align}
\end{theorem}

This result provides the following insights: (i) Under the smoothness assumption for the density function, the kernel estimator achieves a faster rate, exceeding the rate $O(n^{-2/d})$ of the empirical estimator in Theorem \ref{thm:main_plan_based_rates}. This mitigates the curse of dimensionality in the rate of Theorem \ref{thm:main_plan_based_rates} through the smoothness. (ii) The kernel estimator enables the estimation of mappings on a hypercube $[0,1]^d$, which contrasts with the estimator in \cite{manole2024plugin}, which performed estimation on a torus. (iii) This rate is minimax optimal up to a logarithmic factor, as we will see in the next section.  

\subsection{Proof outline for the upper bounds: Stability approach}
\label{subsec:proof-outline}




We present the case of the plan-based estimator below, while the approach is similar for the kernel-based estimator. Below, $a \lesssim b$ means $a \le Cb$ for some constant $C>0$.

Let $\hat\gamma$ be the optimizer of $\mathrm{UOT}(\hat\mu_n,\hat\nu_m)$. Define the fitted row and column masses $\hat r_i \coloneqq \sum_{j=1}^m \hat\gamma_{ij}$ and $\hat s_j \coloneqq \sum_{i=1}^n \hat\gamma_{ij}$, the oracle active masses $\hat r^\star_i \coloneqq e^{-\varphi_0(X_i)}\hat\mu_i$ and $\hat s^\star_j \coloneqq e^{-\psi_0(Y_j)}\hat\nu_j$, the barycentric projections $\hat T_i \coloneqq \frac{1}{\hat r_i}\sum_{j=1}^m \hat\gamma_{ij}Y_j$, and the barycentric error $\hat\Delta_{nm}^{\mathrm{bar}} \coloneqq \sum_{i=1}^n \hat r_i \|\hat T_i-T_0(X_i)\|^2$. A key ingredient of the proof is the following bound, which relates the risk under the population active measure $\gamma_0 \coloneqq e^{-\varphi_0}\mu$ to the discrete UOT solution and the Voronoi geometry (Theorem~\ref{thm:uot_complete_1nn_estimated_mass}):
\begin{equation}
\int_\Omega \|\hat T^{\mathrm{1NN}}(x)-T_0(x)\|^2 d\gamma_0(x) \lesssim n M_n \hat\Delta_{nm}^{\mathrm{bar}} + {D_\KL}(\hat r\mid \hat r^\star) + R_n^2,
\end{equation}
Here, $M_n$ and $R_n$ denote the maximum Voronoi cell mass and radius. By the Vapnik-Chervonenkis inequality (Lemma~\ref{lem:uot_voronoi_mass_radius}), we have $M_n \lesssim \log n / n$ and $\Ep[R_n^2] \lesssim (\log n/n)^{2/d}$ with high probability.

To bound $\hat\Delta_{nm}^{\mathrm{bar}}$ and ${D_\KL}(\hat r\mid \hat r^\star)$, we consider the empirical excess $E_n \coloneqq \mathrm{UOT}(\hat\mu_n,\hat\nu_m) - \int (1-e^{-\varphi_0(x)}) d\hat\mu_n - \int (1-e^{-\psi_0(y)}) d\hat\nu_m$. We show that it can be expressed as follows (Proposition~\ref{prop:uot_barycentric_and_1nn_estimated_mass}):
\begin{align}
\label{eq:excess_identity_summary}
E_n = \sum_{i=1}^n\sum_{j=1}^m \hat\gamma_{ij} \Bigl( \tfrac{1}{2}\| X_i - Y_j\|^2-\varphi_0(X_i)-\psi_0(Y_j) \Bigr) + {D_\KL}(\hat r \mid \hat r^\star) + {D_\KL}(\hat s \mid \hat s^\star).
\end{align}
By dual feasibility, all terms on the right-hand side are nonnegative. Hence, ${D_\KL}(\hat r\mid \hat r^\star) \le E_n$. By Jensen's inequality and Lemma~\ref{lem:uot_gap_sufficient}, we have $\hat\Delta_{nm}^{\mathrm{bar}} \le \sum_{i=1}^n\sum_{j=1}^m \hat\gamma_{ij}\|Y_j - T_0(X_i)\|^2 \le \frac{2}{\kappa} E_n$.

Taking the expectation, we then employ our stability bound (Proposition~\ref{prop:estimated_mass_stability}), which states:
\begin{equation}
\begin{aligned}
\Ep[E_n]  = \Ep[\mathrm{UOT}(\hat\mu_n,\hat\nu_m)] - \mathrm{UOT}(\mu,\nu) &\lesssim \Ep[W_2^2(\hat\mu_n,\mu)] + \Ep[W_2^2(\hat\nu_m,\nu)] + a_n + b_m,
\end{aligned}
\end{equation}
where $a_n$ and $b_m$ are the sequences from Assumption~\ref{assm:masses}. Using any existing bound on the Wasserstein distances of the empirical measures, e.g., \cite{fournier2015rate}, we achieve the final bound \eqref{eq:section4_cube_kernel_map_rate}.

To bound the growth risk of $\hat\lambda^{\mathrm{1NN}}$, we relate the empirical growth map error to the active-source factor error and the transport map error using the following bound (Lemma~\ref{lem:discrete_gh_growth_transfer}):
\begin{equation}
 |\hat\lambda_i-\lambda_0(X_i)|^2 \le |\hat a_i-a_0(X_i)|^2+\|\hat T_i-T_0(X_i)\|^2,
\end{equation}
where $\hat a_i \coloneqq \sqrt{\hat r_i/\hat\mu_i}$ is a plug-in estimate of the active-source factor. Multiplying both sides by $\hat\mu_i$ and summing over $i$, we show in Lemma~\ref{lem:uot_complete_1nn_active_fitted} that
$
\sum_{i=1}^n \hat\mu_i |\hat a_i-a_0(X_i)|^2 \le {D_\KL}(\hat r\mid \hat r^\star)
$. The rest of the proof proceeds analogously to that of the transport map.


\subsection{Minimax lower bound on the hypercube}
\label{sec:minimax_lower_bound}

To match Theorem~\ref{thm:main_cube_kernel_rates} with a lower bound in the same smooth hypercube regime, we restrict to $\Omega=[0,1]^d$, assume that the masses $M_\mu,M_\nu$ are known, and the sample sizes are equal: $m=n$. 
We write $\mathbb P_{\mu,\nu}^n$ and $\mathbb E_{\mu,\nu}^n$ for the joint law and expectation, respectively.

Fix $\alpha>1$ and parameters $M,B<\infty$ and $\Lambda>1$. Let $\mathcal U_\alpha(M,B,\Lambda)$ denote the class of pairs $(\mu,\nu)\in\mM_+([0,1]^d)^2$ such that Assumptions~\ref{assm:curvature}, \ref{assm:positivity}, and~\ref{assm:cube_kernel_smoothness} hold, with smoothness radius $M$ in Assumption~\ref{assm:cube_kernel_smoothness}; the unique $c$-concave solution $\varphi_0$ of the dual problem \eqref{eq:gh_dual} and the associated solution $(T_0,\lambda_0)$ of the Monge problem \eqref{eq:gh_monge} obtained via Theorem~\ref{prop:monge-potential} satisfy $\|T_0\|_{C^\alpha([0,1]^d)}+\|\lambda_0\|_{C^\alpha([0,1]^d)}\le B$ and $\Lambda^{-1}\le \lambda_0(x)\le\Lambda$ for all $x\in[0,1]^d$.

\begin{theorem}[Minimax lower bound in the smooth hypercube regime]
\label{thm:lb-ubot}
For every $\alpha>1$, there exist constants $M_0,B_0<\infty$ and $\Lambda_0>1$ such that, for every $M\ge M_0$, $B\ge B_0$, and $\Lambda\ge\Lambda_0$, there is a constant $c_0>0$, depending only on $d,\alpha,M,B,\Lambda$ and the constants in Assumptions~\ref{assm:curvature} and \ref{assm:positivity}, such that for all sufficiently large $n$,
\begin{align}
    \inf_{\hat T}
    \sup_{(\mu,\nu)\in\mathcal U_\alpha(M,B,\Lambda)}
    \mathbb E_{\mu,\nu}^n\left[
        \int_{[0,1]^d}\|\hat T(x)-T_0(x)\|^2 d\mu(x)
    \right]
    &\ge
    c_0\left(
        n^{ -2\alpha/(2\alpha-2+d)}
        \vee
        n^{-1}
    \right),\\
    \inf_{\hat\lambda}
    \sup_{(\mu,\nu)\in\mathcal U_\alpha(M,B,\Lambda)}
    \mathbb E_{\mu,\nu}^n\left[
        \int_{[0,1]^d}|\hat\lambda(x)-\lambda_0(x)|^2 d\mu(x)
    \right]
    &\ge
    c_0\left(
        n^{ -2\alpha/(2\alpha-2+d)}
        \vee
        n^{-1}
    \right),
\end{align}
where the infima are taken over all measurable estimators of $(X_1,\dots,X_n,Y_1,\dots,Y_n,M_\mu,M_\nu)$.
\end{theorem}

The construction in the appendix lies inside the same regime as Theorem~\ref{thm:main_cube_kernel_rates}: the potential $\varphi_0$ is strongly convex, the associated transport map $T_0$ and growth map $\lambda_0$ have uniform $C^\alpha$ bounds, the density functions are bounded above and away from zero. Hence, in view of Theorem~\ref{thm:main_cube_kernel_rates}, our lower bound indicates that our kernel-based estimator is minimax optimal up to the logarithmic factor.





\section{Experiments}

\subsection{Simulation study}

\begin{figure}[t]
    \centering
    \begin{subfigure}{0.35\textwidth}
        \centering
        \includegraphics[width=\textwidth]{uot_benchmark_summary.pdf}
        \caption{}
        \label{fig:simulation}
    \end{subfigure}
    \hfill 
    \begin{subfigure}{0.64\textwidth}
        \centering
        \includegraphics[width=\textwidth]{uot_3d.pdf}
        \caption{}
        \label{fig:3d}
    \end{subfigure}
    \caption{(a) MSE of the four UOT estimators. Each plot shows the average over 10 seeds with one standard error. Top: The MSEs of estimating $T_0$ and $\lambda_0$ vs $n$. Bottom: Learning rates for $T_0$ and $\lambda_0$ vs. $d$. (b) Top: Incomplete 3D shapes. Middle: Complete 3D shapes predicted by the plan-based 1NN estimator. Bottom: Complete 3D shapes predicted by the plan-based kernel estimator.}
\end{figure}

We sample data from source and target measures $\mu,\nu \in \mM_+([0,1]^d)$ whose densities are $1$-Hölder smooth, with $\mu([0,1]^d)=1$ and $\nu([0,1]^d)=2.5$. We benchmark several estimators of the oracle pair $(T_0, \lambda_0)$, namely plan-based 1NN (PB-1NN), plan-based kernel (PB-Kernel), plug-in kernel (PI-Kernel), and plug-in wavelet (PI-Wavelet cf. Appendix~\ref{sec:wavelet}). We also include the SSUOT estimator proposed by~\cite{vacher2023semi} as a competing baseline.  Additional details are provided in Appendix~\ref{sec:sim_detail}. 

Figure~\ref{fig:simulation} shows MSEs and empirical learning rates of the five methods. Specifically, the two learning rate plots (the bottom row) indicate that PB-1NN, PB-Kernel and SSUOT's learning rates for $T_0$ and $\lambda_0$ are faster than both theoretical upper bound of $n^{-0.5}$ in Theorem~\ref{thm:main_plan_based_rates}, while the learning rates of our plug-in estimators closely match the theoretical upper bound in Theorem~\ref{thm:main_cube_kernel_rates} (the dashed line). Moreover, the two MSE plots (the top row) indicate that our plug-in kernel and wavelet estimators are more accurate and yield faster learning rates than the plan-based estimators and SSUOT.

\subsection{Application to 3D shape completion}

To showcase our plan-based estimators for high-dimensional tasks, we apply them to an unpaired 3D shape completion task using the Completion3D dataset~\cite{Chang2015,Yuan2018}. Specifically, we use the set of incomplete 3D point clouds of chairs as the source dataset and their complete versions as the target dataset. To test robustness, we introduce 30 incomplete cars into the source dataset as outliers. The point clouds are first transformed into a $16\times 16\times 16$ grid and subsequently projected via PCA into 1024 dimensions, yielding a source-target dataset pair with $n=5780$, $m=5750$, and $d=1024$. In this task, both source and target measures have unit mass. Nonetheless, we can leverage UOT to selectively reject outliers through the growth map assignment.

We fit PB-1NN and PB-Kernel to this dataset pair and use them to predict the complete shapes of four chairs and four cars in a held-out test set. As Figure~\ref{fig:3d} illustrates, both estimators successfully recover the shapes of the test chairs. However, their behavior diverges on the outlier test cars: while PB-1NN attempts to reconstruct chairs from the cars, PB-Kernel predicts the mean in the latent PCA space and yields values of $\hat\lambda$ that are close to zero. This indicates PB-Kernel implicitly discards the outliers, demonstrating that it can perform unpaired 3D shape completion in the presence of outliers.


\section{Conclusion}
We studied statistical estimation in unbalanced optimal transport with quadratic cost, where the target estimand consists of the transport map and the growth map. We developed the plan-based estimator and the kernel-based estimator, then derive convergence rates of their estimation errors. We also showed the optimality of the kernel-based estimator by deriving the lower bound. Main technical contribution is a reduction of the estimation to the stability bound, and proving the bound itself. Overall, our results provide a statistical foundation for estimation in UOT and clarify the distinct roles of active marginals, growth estimation, and first-order bias in this problem. 
A limitation is that our setup does not cover entropic regularization, and this is an interesting future work.




\newpage

\appendix

\section{Related work}

\paragraph{Foundations of unbalanced transport.}
Modern UOT is largely built on entropy-transport formulations, which relax the marginal constraints by convex divergences and recover important geometries such as logarithmic entropy-transport and Hellinger-Kantorovich/Wasserstein-Fisher-Rao \cite{ChizatEtAl2018,LieroMielkeSavare2018}. Recent work has clarified the Monge viewpoint, primal-dual optimality conditions, and regularity theory for these models \cite{savare2024relaxation,gallouet2025regularity}. These results provide the structural background for our Monge-type analysis.

\paragraph{Statistical map estimation in balanced OT.}
In the balanced setting, OT map estimation is now supported by a fairly rich statistical theory. Minimax-optimal rates for smooth maps were established via semi-dual curvature arguments by \cite{hutter2021minimax}. On the plug-in side, barycentric-projection estimators based on empirical or smoothed couplings were analyzed by \cite{deb2021barycentric}, while entropic estimators were developed by \cite{Pooladian2021} and shown to be particularly effective in semi-discrete and discontinuous settings by \cite{pooladian2023semidiscrete}. Sharp smooth plug-in guarantees were obtained by \cite{manole2024plugin}, and the stability-based perspective has recently been broadened both to more general function classes and to sharper reductions from map estimation to distribution estimation \cite{divol2025general,balakrishnan2025stability,ponnoprat2025minimax}. Our work can be viewed as an unbalanced counterpart to this stability/plugin line, but with a different target: the transport-growth pair induced by the active marginals of the UOT plan.

\paragraph{Statistical theory for UOT}
The closest prior works are those of \cite{vacher2022stability,vacher2023semi}, which derive semi-dual formulations and global Bregman-type stability bounds for quadratic UOT, leading to the first fast rates for UOT semi-dual potentials. Our contribution differs in both target and technique: rather than estimating a potential under a semi-dual metric, we study the population Monge-type transport-growth pair $(T_0,\lambda_0)$ and relate its risk directly to perturbations of the UOT value and to Wasserstein/$L^1$ errors of fitted measures. {Rigorous statistical results for other UOT models remain scarce.} In particular, \cite{HundrieserEtAl2025} study unbalanced Kantorovich-Rubinstein distances, plans, and barycenters on finite spaces, and \cite{struleva2025sharp} analyze sharp rates for empirical unbalanced KR quantities in spatio-temporal point-process models. These works address different costs, observation models, and statistical targets from the Monge-type estimation problem considered here.

\paragraph{Algorithms and applications.}
A large parallel literature develops scalable algorithms and applications for UOT; see \cite{sejourne2023theorynumerics} for a broad overview. On the computational side, generalized Sinkhorn methods and their complexity or acceleration for unbalanced problems were studied by \cite{pham2020sinkhorn,sejourne2022faster}, and lightweight or neural parameterizations have been proposed for large-scale settings \cite{Gazdieva2024,yang2018scalable,choi2024}. On the application side, UOT and WFR-type models have been used in single-cell trajectory inference and population dynamics \cite{schiebinger2019,sha2023cell}, tumor-growth and reaction-diffusion models \cite{dimarino2020}, gradient-flow learning \cite{yan2024}, and geometric shape analysis \cite{bauer2022SRNF}. These works strongly motivate statistical guarantees for UOT, but they do not provide a nonparametric estimation theory for Monge-type UOT maps.

\section{Proof outline for kernel-based estimator} \label{sec:proof_outline_appendix}



The proof of Theorem~\ref{thm:main_cube_kernel_rates} follows the same abstract stability principle as the proof of Theorem \ref{thm:main_plan_based_rates}, but with the empirical measures $\hat{\mu}_n, \hat{\nu}_m$ replaced by smooth fitted measures $\hat\mu_n^{\mathrm{ker}}$ and $\hat\nu_m^{\mathrm{ker}}$.

\paragraph{(i) Stability argument}: Lemma~\ref{lem:uot_plugin_exact_excess} shows that if one solves UOT between fitted measures $(\hat\mu_n^{\mathrm{ker}},\hat\nu_m^{\mathrm{ker}})$, then the empirical UOT cost decomposes into the same three pieces as in the plan-based case: a map-mismatch term and {two $D_\KL$ penalties} against the oracle active marginals $e^{-\varphi_0}\hat\mu_n^{\mathrm{ker}}$ and $e^{-\psi_0}\hat\nu_m^{\mathrm{ker}}$. Consequently, the transport map risk and the growth-factor risk are reduced to controlling $W_2^2(\hat\mu_n^{\mathrm{ker}},\mu)$, $W_2^2(\hat\nu_m^{\mathrm{ker}},\nu)$, and the corresponding $L^1$ errors.




\paragraph{(ii) Analysis of the kernel}: 
The estimator-specific step is the analysis of the designed kernel function. Proposition~\ref{prop:cube_kernel_density_rates} proves the one-sample bounds for the fitted kernel measures. Informally, for the resolution choice $L_n\asymp n^{1/(d+2(\alpha-1))}$ one obtains
 $   \Ep\bigl[W_2^2(\hat\mu_n^{\mathrm{ker}},\mu)\bigr]
    +
    \Ep\bigl[W_2^2(\hat\nu_m^{\mathrm{ker}},\nu)\bigr]
    \lesssim
    \mathfrak R_n^{\mathrm{ker}}(\alpha)
    +
    \mathfrak R_m^{\mathrm{ker}}(\alpha)
    +
    a_n+b_m,
$
together with matching $L^1$ bounds. The nontrivial point here is the boundary: the estimator cannot use the Fourier kernel directly because the cube densities are not periodic. Instead, we work in the cosine basis, which corresponds to smoothing the even reflection on the doubled torus and preserves the boundary. The coordinatewise multiplier is chosen so that the kernel factorizes exactly into one-dimensional sums, avoiding an $O(L^d)$ summation over multi-indices.
Finally, once the estimated measures satisfy the required $W_2^2$ and $L^1$ rates, combining them to the stability lemmas yields the statement. 




\section{Additional constructions and remarks}
\subsection{Nadaraya-Watson extension of the discrete plan-based estimator}
\label{sec:appendix_nw_extension}

\begin{definition}[Nadaraya-Watson extension]
\label{def:nw_extension_app}
Let $(\hat T_i,\hat a_i,\hat\lambda_i)_{i=1}^n$ be the discrete quantities defined in \eqref{eq:disc-frechet}, \eqref{eq:disc-gh-growth} and their around. Fix a nonnegative kernel $K:\mathbb R^d\to[0,\infty)$ and a bandwidth $h>0$. For $x\in\Omega$, define the normalized weights
\begin{align*}
    w_{i,h}(x)
    \coloneqq
    \frac{K \left((x-X_i)/h\right)}{\sum_{k=1}^n K \left((x-X_k)/h\right)},
\end{align*}
whenever the denominator is positive. The associated Nadaraya--Watson extension is then given by
\begin{align*}
    \hat T_h^{\mathrm{NW}}(x)
    &\coloneqq
    \sum_{i=1}^n w_{i,h}(x)\hat T_i, \quad 
    \hat a_h^{\mathrm{NW}}(x)
    \coloneqq
    \sum_{i=1}^n w_{i,h}(x)\hat a_i, \quad 
    \hat\lambda_h^{\mathrm{NW}}(x)
    \coloneqq
    \sum_{i=1}^n w_{i,h}(x)\hat\lambda_i.
\end{align*}
The second display smooths the auxiliary active-source factor, whereas the third smooths the already corrected Gaussian--Hellinger growth estimator.
\end{definition}
This is the standard kernel-smoothing analogue of the $1$NN extension. As in the main text, the smoothed fields need not preserve the exact empirical marginal constraint, but they provide continuous alternatives when the kernel and the bandwidth are chosen suitably.




\subsection{Obtaining samples through Poisson point processes}\label{sec:appendix_ppp}

We consider that the observations follows two independent Poisson point processes as
\begin{align}
    \mathcal{X} = \{X_i\}_{i=1}^{N_\mu} \sim \mathrm{PPP}(\mu), \qquad \mathcal{Y} = \{Y_j\}_{j=1}^{N_\nu} \sim \mathrm{PPP}(\nu).
\end{align}
Here, the sample sizes satisfy $N_\mu \sim \mathrm{Poisson}(M_\mu)$ and $N_\nu \sim \mathrm{Poisson}(M_\nu)$.
Conditional on $N_\mu$, $X_i$ independently and identically follows $(M_\mu)^{-1}\mu$.
Then, we define an empirical measure
\begin{align}
    \hat{\mu} \coloneqq \sum_{i=1}^{ N_\mu} \delta_{X_i} \qquad \hat{\nu} \coloneqq \sum_{j=1}^{N_\nu} \delta_{Y_j}.
\end{align}
In this setup, we have $\Ep[\hat{\mu}] = \mu$. This setup with slight modification is studied by \cite{struleva2025sharp}.



\subsection{Proof of Lemma~\ref{lem:uot_gap_sufficient}} \label{sec:uot_gap}

We proof Lemma~\ref{lem:uot_gap_sufficient} that allows us to convert the constraint on the Hessian of the UOT potential into a lower bound for the gap in the constraint $\varphi_0(x)+\psi_0(y) \le \frac{1}{2}\|x-y\|^2$. 
\begin{lemma}[{Lemma~\ref{lem:uot_gap_sufficient}, restated}]
Suppose that Assumption~\ref{assm:curvature} holds. Then the dual potentials satisfy the following uniform lower bound on the gap:
\begin{align}
    \frac{1}{2}\|x-y\|^2-\varphi_0(x)-\psi_0(y)
    \ge \frac{\kappa}{2}\|y-T_0(x)\|^2,
    \qquad x,y\in\Omega.
\end{align}
\end{lemma}
\begin{proof}
Define the associated Brenier potential $f_0(x) \coloneqq \frac{1}{2}\|x\|^2 - \varphi_0(x)$. The lower bound in Assumption~\ref{assm:curvature}, namely $\nabla^2 \varphi_0(x) \succeq (1 - \kappa^{-1})I$, implies that the Hessian of the Brenier potential obeys the uniform upper bound:
\begin{align}
    \nabla^2 f_0(x) = I - \nabla^2 \varphi_0(x) \preceq \frac{1}{\kappa} I.
\end{align}
 Let $f_0^*(y) \coloneqq \tfrac{1}{2}\|y\|^2 - \psi_0(y)$ be the Legendre-Fenchel conjugate of $f_0(x)$. Since the Hessian matrix of $f^*_0$ is the inverse of the Hessian matrix of $f_0$, the hypothesis $\nabla^2 f_0(x) \preceq \kappa^{-1} I$ implies
\begin{align}
   \nabla^2 f_0^*(y)  = I - \nabla^2 \psi_0(y) \succeq \kappa I.
\end{align}
 For each fixed $x\in\Omega$, define
\begin{align}
    g_x(y) \coloneqq c(x,y)-\psi_0(y)
    = \tfrac{1}{2}\|x-y\|^2 - \psi_0(y).
\end{align}
Its Hessian is
$
    \nabla_y^2 g_x(y) = I - \nabla^2\psi_0(y) \succeq \kappa I,
$
so $g_x$ is $\kappa$-strongly convex in $y$.
By the $c$-transform relation,
\begin{align}
    \varphi_0(x) = \inf_{y\in\Omega} \bigl\{ c(x,y)-\psi_0(y) \bigr\}
    = g_x(T_0(x)),
\end{align}
hence $T_0(x)$ minimizes $g_x$.
Strong convexity therefore implies
\begin{align}
    g_x(y)-g_x(T_0(x))
    \ge \frac{\kappa}{2}\|y-T_0(x)\|^2,
\end{align}
which is exactly the desired inequality.  
\end{proof}



\section{Stability Bound}

We recall two extensions of the empirical measures $\hat{\mu}_n$ and $\hat{\nu}_m$.
Their normalized empirical probability measures are defined as
\begin{align}
    \bar\mu_n \coloneqq {\hat\mu_n/\hat M_\mu} = {1/n}\sum_{i=1}^n \delta_{X_i}, \qquad \bar\nu_m \coloneqq {\hat\nu_m/\hat M_\nu} = {1/m}\sum_{j=1}^m \delta_{Y_j}.
\end{align}
We also define the weighted empirical measure with the total mass $M_\mu$ as
\begin{align}
    \tilde\mu_n \coloneqq {M_\mu/n}\sum_{i=1}^n \delta_{X_i} = M_\mu\bar\mu_n, \qquad \tilde\nu_m \coloneqq {M_\nu/m}\sum_{j=1}^m \delta_{Y_j} = M_\nu\bar\nu_m.
\end{align}
Note that these measure has the same total mass of $\mu$ and $\nu$, i.e., $\tilde\mu_n(\Omega) = \mu(\Omega) = M_\mu$ and $\tilde\nu_m(\Omega) = \nu(\Omega) = M_\nu$ hold.

Throughout the proof, we denote $c(x,y) = \tfrac{1}{2} \| x - y \|^2$.




\subsection{Stability bound with known masses}

% We temporary define
% \begin{align}
%     \mathrm{UOT}(\hat\mu, \hat\nu) &\coloneqq \min_{\gamma\in\R_+^{n\times m}}
% \left\{
% \sum_{i=1}^n\sum_{j=1}^m C_{ij}\gamma_{ij}
% +{D_\KL}(\gamma_0|\hat\mu)+{D_\KL}(\gamma_1|\hat\nu)
% \right\}.
% \end{align}

We are now ready to state the two-sample stability bound for the UOT problem  with $\tilde\mu_n$ and $\tilde\nu_m$, in the case that the masses $M_\mu$ and $M_\nu$ are known.
\begin{proposition}[Two-sample stability of $\mathrm{UOT}$]
\label{prop:two_sample_stability}
Let $\mu, \nu, \tilde\mu_n, \tilde\nu_m \in \mathcal{M}_+(\Omega)$. {Consider the UOT model introduced in Section~\ref{sec:setup}.} Let $\varphi_0, \psi_0$ be the optimal dual potentials for $(\mu,\nu)$. Assume that the problem setting satisfies Assumptions \ref{assm:curvature}-\ref{assm:positivity}. Then there exists a constant $C_\Lambda$ such that:
\begin{align}
\label{eq:uot_two_sample}
&\mathrm{UOT}(\tilde\mu_n, \tilde\nu_m) - \mathrm{UOT}(\mu, \nu) \\
&\leq \int \zeta_0 d(\tilde\mu_n - \mu) + \int \xi_0 d(\tilde\nu_m - \nu) + C_\Lambda \left( M_\mu W_2^2(\bar\mu_n, \bar\mu) + M_\nu W_2^2(\bar\nu_m, \bar\nu) \right).
\end{align}
\end{proposition}

\begin{proof}[Proof of Proposition \ref{prop:two_sample_stability}]
Let $\gamma$ be the optimal continuous coupling for $(\mu,\nu)$. Since $\mu$ and $\nu$ are both absolutely continuous (Assumption~\ref{assm:positivity}), Brenier's theorem ensures the existence of optimal transport maps $T_\mu: \Omega \to \Omega$ and $T_\nu: \Omega \to \Omega$ such that $(T_\mu)_\# \mu = \tilde\mu_n$, $(T_\nu)_\# \nu = \tilde\nu_m$, $\int \|T_\mu(x)-x\|^2 d\mu(x) = W_2^2(\tilde\mu_n, \mu)=M_\mu W_2^2(\bar\mu_n,\bar\mu)$, and $\int \|T_\nu(y)-y\|^2 d\nu(y) = W_2^2(\tilde\nu_m, \nu)=M_\nu W_2^2(\bar\nu_m,\bar\nu)$. Define the candidate coupling $\hat \gamma = (T_\mu, T_\nu)_\# \gamma$.
{By the data-processing inequality, $D_\KL$ contracts under deterministic maps:}
$$ {D_\KL}(\hat\gamma_0 \mid \tilde\mu_n) = {D_\KL}((T_\mu)_\# \gamma_0 \mid (T_\mu)_\# \mu) \leq {D_\KL}(\gamma_0 \mid \mu). $$
Similarly, ${D_\KL}(\hat\gamma_1 \mid \tilde\nu_m) \leq {D_\KL}(\gamma_1 \mid \nu)$. Thus, the UOT objective for the candidate $\hat \gamma$ is bounded:
$$ \mathrm{UOT}(\tilde\mu_n, \tilde\nu_m) \leq \int c d\hat\gamma + {D_\KL}(\hat\gamma_0 \mid \tilde\mu_n) + {D_\KL}(\hat\gamma_1 \mid \tilde\nu_m) \leq \int c d\hat\gamma + {D_\KL}(\gamma_0 \mid \mu) + {D_\KL}(\gamma_1 \mid \nu). $$
Subtracting this inequality by $\mathrm{UOT}(\mu,\nu) = \int c d\gamma + {D_\KL}(\gamma_0 \mid \mu) + {D_\KL}(\gamma_1 \mid \nu)$, we get:
$$ \mathrm{UOT}(\tilde\mu_n, \tilde\nu_m) - \mathrm{UOT}(\mu,\nu) \leq \int c d\hat\gamma - \int c d\gamma. $$
By Lemma~\ref{lem:uot_cost_expansion}, we have
\begin{equation}
\label{eq:cost_diff_exact_0}
\begin{aligned}
\int c d\hat\gamma - \int c d\gamma 
&\le \int \|T_\mu(x)-x\|^2 d\gamma_0(x) + \int \|T_\nu(y)-y\|^2 d\gamma_1(y) \\
&\quad + \int \langle T_\mu(x)-x, x-y \rangle d\gamma(x,y) \\
&\quad + \int \langle T_\nu(y)-y, y-x \rangle d\gamma(x,y).
\end{aligned}
\end{equation}
First, we bound the two terms that involve $T_\mu(x)$.
By Assumption \ref{assm:curvature}, the dual potentials are bounded, so $M_\varphi \coloneqq \sup_{x\in\Omega} e^{-\varphi_0(x)} < \infty$ and $M_\psi \coloneqq \sup_{y\in\Omega} e^{-\psi_0(y)} < \infty$. Consequently,
\begin{align}
    \int \|T_\mu(x)-x\|^2 d\gamma_0(x)
    &=
    \int \|T_\mu(x)-x\|^2 e^{-\varphi_0(x)} d\mu(x)
    \\
    &\le
    M_\varphi \int \|T_\mu(x)-x\|^2 d\mu(x)
    =
    M_\varphi M_\mu W_2^2(\bar\mu_n, \bar\mu).
\end{align}
For the other term, we use the first-order optimality relation
\begin{align*}
    x-y = \nabla \varphi_0(x)
    \qquad\text{for $\gamma$-a.e.\ }(x,y),
\end{align*}
which follows from $c(x,y)=\tfrac{1}{2}\|x-y\|^2$ and the complementary slackness for the
UOT dual: for $\gamma$-a.e.\ $(x,y)$, $\varphi_0(x)+\psi_0(y)=c(x,y)$, while
$\varphi_0(x')+\psi_0(y)\le c(x',y)$ for all $x'\in\Omega$. Differentiating in $x$
at points where $\varphi_0$ is differentiable (which holds $\mu$-a.e., hence
$\gamma_0$-a.e.\ since $\gamma_0\ll\mu$, by Theorem~\ref{prop:monge-potential}
under Assumption~\ref{assm:curvature}) yields $\nabla c(x,y)=\nabla\varphi_0(x)$,
i.e., $x-y=\nabla\varphi_0(x)$.
Therefore,
\begin{align*}
    \int \langle T_\mu(x)-x, x-y \rangle d\gamma(x,y)
    &=
    \int \langle T_\mu(x)-x, \nabla \varphi_0(x) \rangle d\gamma_0(x)
    \\
    &=
    \int \langle T_\mu(x)-x, \nabla \varphi_0(x) \rangle e^{-\varphi_0(x)} d\mu(x)
    \\
    &=
    \int \langle T_\mu(x)-x, \nabla \zeta_0(x) \rangle d\mu(x),
\end{align*}
because $\nabla \zeta_0(x)=e^{-\varphi_0(x)}\nabla \varphi_0(x)$.
Since both $\nabla \zeta_0$ and $\nabla \xi_0$ are Lipschitz on $\Omega$ by Lemma~\ref{lem:smoothness_effective}, by taking
$\Lambda \coloneqq \max\bigl(\mathrm{Lip}(\nabla\zeta_0),\mathrm{Lip}(\nabla\xi_0)\bigr)>0$,
we have the following inequality for all $x,y\in\Omega$:
\begin{align*}
    \bigl|\zeta_0(y)-\zeta_0(x)-\langle \nabla\zeta_0(x),y-x\rangle\bigr|
    \le
    \frac{\Lambda}{2}\|y-x\|^2.
\end{align*}
Applying this with $y=T_\mu(x)$ yields
\begin{align*}
    \langle \nabla \zeta_0(x), T_\mu(x)-x \rangle
    \le
    \zeta_0(T_\mu(x))-\zeta_0(x)+\frac{\Lambda}{2}\|T_\mu(x)-x\|^2.
\end{align*}
Integrating with respect to $\mu$ and using $(T_\mu)_\#\mu=\tilde\mu_n$, we obtain
\begin{align*}
    \int \langle T_\mu(x)-x, \nabla \zeta_0(x) \rangle d\mu(x)
    &\le
    \int (\zeta_0(T_\mu(x))-\zeta_0(x)) d\mu(x)
    + \frac{\Lambda}{2}\int \|T_\mu(x)-x\|^2 d\mu(x)
    \\
    &=
    \int \zeta_0 d(\tilde\mu_n-\mu)
    + \frac{\Lambda}{2}M_\mu W_2^2(\bar\mu_n,\bar\mu).
\end{align*}
The terms in \eqref{eq:cost_diff_exact_0} that involve $T_\nu(y)$ are handled in the same way, using $y-x=\nabla\psi_0(y)$ for $\gamma$-a.e.\ $(x,y)$ and noting that $\nabla\xi_0$ is also $\Lambda$-Lipschitz by our choice of $\Lambda$.
Combining both bounds, we may take
\begin{align*}
    C_\Lambda \coloneqq \max(M_\varphi,M_\psi)+\frac{\Lambda}{2}.
\end{align*}
This proves \eqref{eq:uot_two_sample}.
\end{proof}


\subsection{Extension to estimated masses}

We now remove the assumption that the empirical measures have exactly the same total masses as the target measures. We have $\hat\mu_n = \alpha_n \tilde\mu_n$ and $ \hat\nu_m = \beta_m \tilde\nu_m,$
where $    \alpha_n \coloneqq {\hat{M}_\mu}/  {M_\mu}$ and $ \beta_m \coloneqq {\hat{M}_\nu} / {M_\nu}$. Our bound relies on the following assumption that the masses are estimable:


{In preparation, we study the fixed $D_\KL$-penalized formulation.}
Since $\hat\mu_n$ and $\tilde\mu_n$ have the same support and differ only by a global multiplicative factor, the change in the objective can be controlled directly through the {$D_\KL$ terms}.

\begin{lemma}[$D_\KL$ under measure scaling]
\label{lem:kl_scaling_reference}
Let $\eta,\mu \in \mathcal{M}_+(\Omega)$ with $\eta \ll \mu$, write $M_\mu \coloneqq \mu(\Omega)$ and let $\alpha>0$. Then
\begin{align}
    {D_\KL}(\eta\mid \alpha\mu) = {D_\KL}(\eta\mid \mu) - \eta(\Omega)\log \alpha + (\alpha-1)M_\mu.
\end{align}
\end{lemma}


\begin{proof}[Proof of Lemma \ref{lem:kl_scaling_reference}]
Write $\eta = f\mu$. Then $\eta = (f/\alpha)(\alpha\mu)$, and therefore
\begin{align*}
    {D_\KL}(\eta\mid \alpha\mu)
    &= \int \left( \frac{f}{\alpha}\log\frac{f}{\alpha} - \frac{f}{\alpha} + 1 \right)  \d(\alpha\mu) \\
    &= \int \bigl( f\log f - f + 1 \bigr)  \d\mu - (\log \alpha)\int f \d\mu + (\alpha-1)M_\mu \\
    &= {D_\KL}(\eta\mid \mu) - \eta(\Omega)\log \alpha + (\alpha-1)M_\mu.
\end{align*}
\end{proof}

\begin{lemma}[Bound on $\gamma(\Omega\times\Omega)$]
\label{lem:uot_plan_mass_bound}
Let $\eta,\sigma \in \mathcal{M}_+(\Omega)$ with $\eta(\Omega)=A>0$ and $\sigma(\Omega)=B>0$, and let $\gamma$ be any optimal plan for $\mathrm{UOT}(\eta,\sigma)$ with quadratic cost. Then
\begin{align}
    \gamma(\Omega\times\Omega) \le e\sqrt{AB}.
\end{align}
\end{lemma}


\begin{proof}[Proof of Lemma \ref{lem:uot_plan_mass_bound}]
Let $s \coloneqq \gamma(\Omega\times\Omega) = \gamma_0(\Omega)=\gamma_1(\Omega)$. Since the zero plan is feasible, we have
\begin{align}
    \mathrm{UOT}(\eta,\sigma) \le \eta(\Omega) + \sigma(\Omega) = A+B.
\end{align}
Because the transport cost is nonnegative, it follows that
\begin{align}
    {D_\KL}(\gamma_0\mid \eta) + {D_\KL}(\gamma_1\mid \sigma) \le A+B.
\end{align}
Now let $f = \d\gamma_0/\d\eta$. Since $t\mapsto t\log t - t + 1$ is convex, Jensen's inequality yields
\begin{align}
    {D_\KL}(\gamma_0\mid \eta)
    &= \int \bigl( f\log f - f + 1 \bigr) \d\eta \\
    &\ge A\left( \frac{s}{A}\log\frac{s}{A} - \frac{s}{A} + 1 \right)
     = s\log\frac{s}{A} - s + A.
\end{align}
Similarly,
\begin{align}
    {D_\KL}(\gamma_1\mid \sigma) \ge s\log\frac{s}{B} - s + B.
\end{align}
Summing the last two inequalities and using the upper bound $A+B$ gives
\begin{align}
    s\log\frac{s^2}{AB} - 2s + A+B \le A+B,
\end{align}
that is,
\begin{align}
    s\left( \log\frac{s^2}{AB} - 2 \right) \le 0.
\end{align}
If $s=0$ there is nothing to prove. Otherwise $\log(s^2/(AB))\le 2$, hence $s \le e\sqrt{AB}$.
\end{proof}

\begin{proposition}[Stability of $\mathrm{UOT}$ under estimated masses]
\label{prop:estimated_mass_stability}
Assume that Assumptions~\ref{assm:curvature}, \ref{assm:supp_global}, \ref{assm:positivity}, and \ref{assm:masses} hold.
Let $C_\Lambda$ be the stability constant from Proposition~\ref{prop:two_sample_stability}.
Then, almost surely,
\begin{equation}
\label{eq:estimated_mass_stability}
\begin{aligned}
\mathrm{UOT}(\hat\mu_n,\hat\nu_m) - \mathrm{UOT}(\mu,\nu)
&\le \int \zeta_0  \d(\tilde\mu_n-\mu) + \int \xi_0  \d(\tilde\nu_m-\nu) \\
&\quad + C_\Lambda\Bigl( M_\mu W_2^2(\bar\mu_n,\bar\mu) + M_\nu W_2^2(\bar\nu_m,\bar\nu) \Bigr) \\
&\quad + \Bigl( M_\mu|\alpha_n-1| + M_\nu|\beta_m-1| + e\sqrt{M_\mu M_\nu}\bigl(|\log \alpha_n| + |\log \beta_m|\bigr) \Bigr).
\end{aligned}
\end{equation}
Consequently, there exists a constant $C_{M_\mu,M_\nu}>0$ such that, for all sufficiently large $n,m$, the following bound holds:
\begin{equation}
\label{eq:estimated_mass_stability_linearized}
\begin{aligned}
\Ep\bigl[\mathrm{UOT}(\hat\mu_n,\hat\nu_m)\bigr] - \mathrm{UOT}(\mu,\nu)
&\le C_\Lambda\Bigl( \Ep[M_\mu W_2^2(\bar\mu_n,\bar\mu)] + \Ep[M_\nu W_2^2(\bar\nu_m,\bar\nu)] \Bigr) \\
&\quad + C_{M_\mu,M_\nu}\bigl( a_n + b_m \bigr),
\end{aligned}
\end{equation}
where $a_n,b_m$ are the rates from Assumption~\ref{assm:masses}.
In particular, if the right-hand side of Proposition~\ref{prop:two_sample_stability} applied to $(\tilde\mu_n,\tilde\nu_m)$ converges to $0$ and Assumption \ref{assm:masses} holds, then, we obtain the following convergence in probability:
\begin{align}
    \mathrm{UOT}(\hat\mu_n,\hat\nu_m) - \mathrm{UOT}(\mu,\nu) \to 0.
\end{align}
\end{proposition}

\begin{proof}[Proof of Proposition \ref{prop:estimated_mass_stability}]
Let $\tilde\gamma_{n,m}$ be an optimal plan for $\mathrm{UOT}(\tilde\mu_n,\tilde\nu_m)$, and define 
    $s_{n,m} \coloneqq \tilde\gamma_{n,m}(\Omega\times\Omega)$.
Since we have $\hat\mu_n = \alpha_n\tilde\mu_n$ and $\hat\nu_m = \beta_m\tilde\nu_m$, the same plan $\tilde\gamma_{n,m}$ is admissible for $\mathrm{UOT}(\hat\mu_n,\hat\nu_m)$. Hence, we obtain
\begin{align*}
    \mathrm{UOT}(\hat\mu_n,\hat\nu_m)
    &\le \int c \d\tilde\gamma_{n,m}
    + {D_\KL}((\tilde\gamma_{n,m})_0\mid \hat\mu_n)
    + {D_\KL}((\tilde\gamma_{n,m})_1\mid \hat\nu_m).
\end{align*}
Applying Lemma~\ref{lem:kl_scaling_reference} with $\eta=(\tilde\gamma_{n,m})_0$ and $\mu=\tilde\mu_n$, and then with $\eta=(\tilde\gamma_{n,m})_1$ and $\mu=\tilde\nu_m$, yields
\begin{align*}
    {D_\KL}((\tilde\gamma_{n,m})_0\mid \hat\mu_n)
    &= {D_\KL}((\tilde\gamma_{n,m})_0\mid \tilde\mu_n) - s_{n,m}\log \alpha_n + (\alpha_n-1)M_\mu, \\
    {D_\KL}((\tilde\gamma_{n,m})_1\mid \hat\nu_m)
    &= {D_\KL}((\tilde\gamma_{n,m})_1\mid \tilde\nu_m) - s_{n,m}\log \beta_m + (\beta_m-1)M_\nu.
\end{align*}
Therefore, we obtain
\begin{align}
\begin{aligned}
    \mathrm{UOT}(\hat\mu_n,\hat\nu_m) - \mathrm{UOT}(\tilde\mu_n,\tilde\nu_m)
    \le \Bigl( -s_{n,m}(\log \alpha_n + \log \beta_m) + (\alpha_n-1)M_\mu + (\beta_m-1)M_\nu \Bigr).
\end{aligned}
\end{align}
By Lemma~\ref{lem:uot_plan_mass_bound}, we also obtain
\begin{align}
    s_{n,m} \le e\sqrt{M_\mu M_\nu}.
\end{align}
We then obtain
\begin{align}
\label{eq:empirical_mass_error_term}
\begin{aligned}
    &\mathrm{UOT}(\hat\mu_n,\hat\nu_m) - \mathrm{UOT}(\tilde\mu_n,\tilde\nu_m) \\
    &\le \Bigl( M_\mu|\alpha_n-1| + M_\nu|\beta_m-1| + e\sqrt{M_\mu M_\nu}\bigl(|\log \alpha_n| + |\log \beta_m|\bigr) \Bigr).
\end{aligned}
\end{align}
Next, Proposition~\ref{prop:two_sample_stability} applies to $(\tilde\mu_n,\tilde\nu_m)$ because $\tilde\mu_n(\Omega)=M_\mu$ and $\tilde\nu_m(\Omega)=M_\nu$. Hence
\begin{align}
\label{eq:balanced_part_estimated_mass}
\begin{aligned}
    &\mathrm{UOT}(\tilde\mu_n,\tilde\nu_m) - \mathrm{UOT}(\mu,\nu) \\
    &\le \int \zeta_0  \d(\tilde\mu_n-\mu) + \int \xi_0  \d(\tilde\nu_m-\nu)  + C_\Lambda\Bigl( M_\mu W_2^2(\bar\mu_n,\bar\mu) + M_\nu W_2^2(\bar\nu_m,\bar\nu) \Bigr).
\end{aligned}
\end{align}
Adding \eqref{eq:empirical_mass_error_term} and \eqref{eq:balanced_part_estimated_mass} proves  \eqref{eq:estimated_mass_stability}.

We now derive \eqref{eq:estimated_mass_stability_linearized}.
Taking expectations in \eqref{eq:estimated_mass_stability}, the linear terms
$\Ep[\int\zeta_0  d(\tilde\mu_n-\mu)]$ and $\Ep[\int\xi_0  d(\tilde\nu_m-\nu)]$
vanish by unbiasedness of $\tilde\mu_n=M_\mu\bar\mu_n$ and $\tilde\nu_m=M_\nu\bar\nu_m$.
By the a.s.\ lower bound $\hat M_\mu\ge cM_\mu$ in Assumption~\ref{assm:masses}, $\alpha_n\ge c$ a.s., so the mean-value theorem applied to $\log$ on $[c,\infty)$ gives
\(
|\log\alpha_n| \le |\alpha_n-1|/c
\)
a.s., and analogously $|\log\beta_m|\le|\beta_m-1|/c$ a.s. Hence, using the identities $M_\mu|\alpha_n-1|=|\hat M_\mu-M_\mu|$ and $M_\nu|\beta_m-1|=|\hat M_\nu-M_\nu|$,
\begin{align*}
    \Ep \bigl[ M_\mu|\alpha_n-1| + e\sqrt{M_\mu M_\nu} |\log\alpha_n| \bigr]
    &\le \Ep[|\hat M_\mu-M_\mu|] + \frac{e\sqrt{M_\mu M_\nu}}{c M_\mu} \Ep[|\hat M_\mu-M_\mu|] \\
    &\le \left(1 + c^{-1}e\sqrt{M_\nu/M_\mu}\right) a_n,
\end{align*}
and analogously $\Ep \bigl[ M_\nu|\beta_m-1| + e\sqrt{M_\mu M_\nu} |\log\beta_m| \bigr]
    \le \left(1 + c^{-1}e\sqrt{M_\mu/M_\nu}\right) b_m$. Setting
$C_{M_\mu,M_\nu}\coloneqq 1+c^{-1}e\sqrt{M_\mu/M_\nu}+c^{-1}e\sqrt{M_\nu/M_\mu}$
and absorbing the constant yields \eqref{eq:estimated_mass_stability_linearized}.
The final convergence claim is immediate from \eqref{eq:estimated_mass_stability} together with $\alpha_n,\beta_m\to 1$ and $|\log\alpha_n|,|\log\beta_m|\to 0$ a.s.
\end{proof}

\subsection{Supportive Results}


In preparation, we introduce the following result to characterize the Monge map.
\begin{proposition}[Monge structure of $T_0$]
\label{prop:population_target_monge}
Assume that the population UOT optimizer is deterministic, namely
\begin{align}
    \gamma = (\mathrm{id},T_0)_\# \gamma_0.
\end{align}
Then $\gamma$ is an optimal coupling for the \emph{balanced} Kantorovich problem between the active marginals $\gamma_0$ and $\gamma_1$, i.e.
\begin{align}
    \int_{\Omega^2} c(x,y)   d\gamma(x,y)
    =
    \inf_{\pi\in \Pi(\gamma_0,\gamma_1)}
    \int_{\Omega^2} c(x,y)   d\pi(x,y).
\end{align}
Consequently, $T_0$ solves the Monge problem
\begin{align}
\label{eq:population_monge_problem_active}
    \inf_{T:  T_\# \gamma_0 = \gamma_1}
    \int_\Omega c(x,T(x))   d\gamma_0(x).
\end{align}
If, in addition, $c(x,y)=\tfrac{1}{2}\|x-y\|^2$ and $\gamma_0$ is absolutely continuous with respect to Lebesgue measure, then this Monge solution is unique $\gamma_0$-almost everywhere.
\end{proposition}

\begin{proof}[Proof of Proposition \ref{prop:population_target_monge}]
Fix any $\pi \in \Pi(\gamma_0,\gamma_1)$. {Since $\pi$ has the same marginals as $\gamma$, the $D_\KL$ terms in the UOT objective coincide:}
\begin{align}
    {D_\KL}(\pi_0\mid \mu) = {D_\KL}(\gamma_0\mid \mu),
    \qquad
    {D_\KL}(\pi_1\mid \nu) = {D_\KL}(\gamma_1\mid \nu).
\end{align}
Because $\gamma$ minimizes the UOT objective, we therefore obtain
\begin{align}
    \int_{\Omega^2} c(x,y)   d\gamma(x,y)
    \le
    \int_{\Omega^2} c(x,y)   d\pi(x,y)
    \qquad
    \forall \pi \in \Pi(\gamma_0,\gamma_1),
\end{align}
which proves the balanced Kantorovich optimality. Since $\gamma=(\mathrm{id},T_0)_\#\gamma_0$, we have $T_{0\#}\gamma_0=\gamma_1$ and
\begin{align}
    \int_{\Omega^2} c(x,y)   d\gamma(x,y)
    =
    \int_\Omega c(x,T_0(x))   d\gamma_0(x),
\end{align}
so $T_0$ attains the infimum in \eqref{eq:population_monge_problem_active}.
For uniqueness, let $S:\Omega\to\Omega$ be any other minimizer in \eqref{eq:population_monge_problem_active}.
Then $(\mathrm{id},S)_\#\gamma_0$ is also an optimal coupling in $\Pi(\gamma_0,\gamma_1)$.
Because $\gamma_0$ is absolutely continuous and the cost is $\tfrac{1}{2}\|x-y\|^2$, the balanced transport problem between $\gamma_0$ and $\gamma_1$ has a unique optimal plan, and that plan is induced by a map.
Hence
\begin{align*}
    (\mathrm{id},S)_\#\gamma_0
    =
    \gamma
    =
    (\mathrm{id},T_0)_\#\gamma_0,
\end{align*}
which implies $S=T_0$ $\gamma_0$-almost everywhere.
\end{proof}



We will need the following two lemmas in order to prove the stability bound:
\begin{lemma}[Lipschitzness of $\nabla\zeta_0$ and $\nabla\xi_0$]
\label{lem:smoothness_effective}
Assume Assumptions~\ref{assm:curvature} and ~\ref{assm:supp_global}, and define
\begin{align*}
    \zeta_0(x) \coloneqq -\bigl(e^{-\varphi_0(x)}-1\bigr),
    \qquad
    \xi_0(y) \coloneqq -\bigl(e^{-\psi_0(y)}-1\bigr).
\end{align*}
Then $\nabla \zeta_0$ and $\nabla \xi_0$ are Lipschitz on $\Omega$.
\end{lemma}

\begin{proof}[Proof of Lemma \ref{lem:smoothness_effective}]
We first prove this for $\zeta_0$. The gradient is $\nabla \zeta_0(x) = e^{-\varphi_0(x)} \nabla \varphi_0(x)$. By the product rule,
\begin{align*}
    \nabla^2 \zeta_0(x) = - e^{-\varphi_0(x)} (\nabla \varphi_0(x))(\nabla \varphi_0(x))^\top + e^{-\varphi_0(x)} \nabla^2 \varphi_0(x).
\end{align*}
By Assumption~\ref{assm:curvature} and~\ref{assm:supp_global}, $\varphi_0$ and $\nabla \varphi_0(x)$ are bounded on the compact set $\Omega$, and $\|\nabla^2 \varphi_0(x)\|_{\mathrm{op}} \le \Lambda \coloneqq \max(1-\kappa, \kappa^{-1}-1)$. Thus $\nabla \zeta_0$ is $\Lambda$-Lipschitz.

For $\xi_0$, we have the dual relation $(I - \nabla^2 \psi_0) = (I - \nabla^2 \varphi_0)^{-1}$ $\mu$-a.e., which implies $(1-\kappa^{-1})I \preceq \nabla^2 \psi_0 \preceq (1-\kappa)I$. The rest of the proof remains identical to that of $\zeta_0$.
\end{proof}


\begin{lemma}[Bound on $\int c d\hat\gamma-\int c d\gamma$]
\label{lem:uot_cost_expansion}
Let $\gamma \in \mathcal{M}_+(\Omega \times \Omega)$ have marginals $\gamma_0$ and $\gamma_1$, and let $T_\mu, T_\nu: \Omega \to \Omega$ be measurable maps.
Define the pushforward coupling $\hat\gamma = (T_\mu, T_\nu)_\# \gamma$. Then
\begin{equation}
\label{eq:cost_diff_exact}
\begin{aligned}
\int c d\hat\gamma - \int c d\gamma 
&\le \int \|T_\mu(x)-x\|^2 d\gamma_0(x) + \int \|T_\nu(y)-y\|^2 d\gamma_1(y) \\
&\quad + \int \langle T_\mu(x)-x, x-y \rangle d\gamma(x,y) \\
&\quad + \int \langle T_\nu(y)-y, y-x \rangle d\gamma(x,y).
\end{aligned}
\end{equation}
\end{lemma}


\begin{proof}[Proof of Lemma \ref{lem:uot_cost_expansion}]
By definition of the pushforward,
\begin{align*}
    \int c(x,y) d\hat\gamma(x,y)
    =
    \frac{1}{2}\int \|T_\mu(x) - T_\nu(y)\|^2 d\gamma(x,y).
\end{align*}
Consequently,
\begin{align*}
\tfrac{1}{2}\|T_\mu(x) - T_\nu(y)\|^2
&=
\tfrac{1}{2}\|(T_\mu(x)-x) - (T_\nu(y)-y) + (x-y)\|^2 \\
&=
\tfrac{1}{2}\|x-y\|^2 + \tfrac{1}{2}\|T_\mu(x)-x\|^2 + \tfrac{1}{2}\|T_\nu(y)-y\|^2 \\
&\quad + \langle x-y, T_\mu(x)-x \rangle - \langle x-y, T_\nu(y)-y \rangle \\
&\quad - \langle T_\mu(x)-x, T_\nu(y)-y \rangle \\
&\le
\tfrac{1}{2}\|x-y\|^2 + \|T_\mu(x)-x\|^2 + \|T_\nu(y)-y\|^2 \\
&\quad + \langle T_\mu(x)-x, x-y \rangle + \langle T_\nu(y)-y, y-x \rangle.
\end{align*}
Integrating this inequality with respect to $d\gamma(x,y)$ and using the identities
\begin{align*}
    \int \|T_\mu(x)-x\|^2 d\gamma(x,y)
    &=
    \int \|T_\mu(x)-x\|^2 d\gamma_0(x),\\
    \int \|T_\nu(y)-y\|^2 d\gamma(x,y)
    &=
    \int \|T_\nu(y)-y\|^2 d\gamma_1(y),
\end{align*}
yields \eqref{eq:cost_diff_exact}.
\end{proof}


\section{Proof of Theorem \ref{thm:main_plan_based_rates}} \label{sec:proof_plan_based_rates}

We now formulate a first empirical two-sample analogue of \cite[Proposition~13]{manole2024plugin} for the {$D_\KL$-penalized} unbalanced problem.
The main point is that, in the unbalanced case, the natural excess quantity is no longer purely a transport mismatch term.
Besides the deviation of $Y_j$ from the population Monge map values $T_0(X_i)$ from Proposition~\ref{prop:population_target_monge}, one must also keep track of how the fitted row and column masses differ from the \emph{active} population marginals transported by the optimal Monge pair $(T_0,\lambda_0)$.


Let $\tilde\gamma=(\tilde\gamma_{ij})$ be any optimizer of $\mathrm{UOT}(\tilde\mu_n,\tilde\nu_m)$, and write
\begin{align}
    \tilde r_i \coloneqq \sum_{j=1}^m \tilde\gamma_{ij},
    \qquad
    \tilde s_j \coloneqq \sum_{i=1}^n \tilde\gamma_{ij}.
\end{align}
We also denote the atomic masses of the oracle empirical measures by
    $\tilde\mu_i \coloneqq {M_\mu} / {n}$ and $ \tilde\nu_j \coloneqq {M_\nu} / {m}$.
Finally, recall that for the population optimizer $(\varphi_0,\psi_0)$ we have
    $\gamma_0 = e^{-\varphi_0}\mu$ and $
    \gamma_1 = e^{-\psi_0}\nu$,
and $\zeta_0(x)= -(e^{-\varphi_0(x)}-1)$, $\xi_0(y)= -(e^{-\psi_0(y)}-1)$. Throughout the proof, we denote $c(x,y) = \tfrac{1}{2} \| x - y \|^2$.



\begin{proposition}[Empirical UOT excess identity]
\label{prop:uot_exact_empirical_excess}
Define the oracle empirical active marginals
$r_i^\star \coloneqq e^{-\varphi_0(X_i)}\tilde\mu_i$ and $s_j^\star \coloneqq e^{-\psi_0(Y_j)}\tilde\nu_j$.
Then, we have
\begin{equation}
\label{eq:uot_exact_empirical_excess}
\begin{aligned}
&\mathrm{UOT}(\tilde\mu_n,\tilde\nu_m)
- \int \zeta_0  d\tilde\mu_n
- \int \xi_0  d\tilde\nu_m \\
&\qquad=
\sum_{i=1}^n\sum_{j=1}^m
\tilde\gamma_{ij}
\Bigl(
c(X_i,Y_j)-\varphi_0(X_i)-\psi_0(Y_j)
\Bigr) 
+ {D_\KL}(\tilde r \mid r^\star)
+ {D_\KL}(\tilde s \mid s^\star).
\end{aligned}
\end{equation}
In particular, the right-hand side is nonnegative.
Moreover, we obtain
\begin{equation}
\label{eq:uot_exact_empirical_excess_expectation}
\begin{aligned}
&\Ep\Biggl[
\sum_{i=1}^n\sum_{j=1}^m
\tilde\gamma_{ij}
\Bigl(
c(X_i,Y_j)-\varphi_0(X_i)-\psi_0(Y_j)
\Bigr)
+ {D_\KL}(\tilde r \mid r^\star)
+ {D_\KL}(\tilde s \mid s^\star)
\Biggr] \\
&\qquad=
\Ep\Bigl[
\mathrm{UOT}(\tilde\mu_n,\tilde\nu_m) - \mathrm{UOT}(\mu,\nu)
\Bigr].
\end{aligned}
\end{equation}
\end{proposition}

\begin{proof}[Proof of Proposition \ref{prop:uot_exact_empirical_excess}]
Since $\tilde\gamma$ is optimal for $\mathrm{UOT}(\tilde\mu_n,\tilde\nu_m)$, we have
\begin{align}
\label{eq:uot_empirical_objective_decomp_start}
\mathrm{UOT}(\tilde\mu_n,\tilde\nu_m)
=
\sum_{i=1}^n\sum_{j=1}^m c(X_i,Y_j)\tilde\gamma_{ij}
+ {D_\KL}(\tilde r \mid \tilde\mu)
+ {D_\KL}(\tilde s \mid \tilde\nu).
\end{align}
By the definition of $r_i^\star = e^{-\varphi_0(X_i)}\tilde\mu_i$, we obtain
\begin{align}
\label{eq:source_kl_rewrite}
\begin{aligned}
{D_\KL}(\tilde r \mid \tilde\mu)
&=
\sum_{i=1}^n
\left[
\tilde r_i \log \left(\frac{\tilde r_i}{r_i^\star}\right)
- \tilde r_i + r_i^\star
\right]
- \sum_{i=1}^n \tilde r_i \log \left(\frac{r_i^\star}{\tilde\mu_i}\right)
+ \sum_{i=1}^n (\tilde\mu_i-r_i^\star) \\
&=
{D_\KL}(\tilde r \mid r^\star)
- \sum_{i=1}^n \tilde r_i   \varphi_0(X_i)
+ \sum_{i=1}^n \zeta_0(X_i)\tilde\mu_i,
\end{aligned}
\end{align}
where the second equality follows
\begin{align}
    \log \left(\frac{r_i^\star}{\tilde\mu_i}\right)
    = -\varphi_0(X_i),
    \qquad
    \left(1-e^{-\varphi_0(X_i)}\right)
    = \zeta_0(X_i).
\end{align}
Similarly, we obtain
\begin{equation}
\label{eq:target_kl_rewrite}
{D_\KL}(\tilde s \mid \tilde\nu)
=
{D_\KL}(\tilde s \mid s^\star)
- \sum_{j=1}^m \tilde s_j   \psi_0(Y_j)
+ \sum_{j=1}^m \xi_0(Y_j)\tilde\nu_j .
\end{equation}
Substituting \eqref{eq:source_kl_rewrite} and \eqref{eq:target_kl_rewrite} into \eqref{eq:uot_empirical_objective_decomp_start}, and using the relations
\begin{align}
    \sum_{i=1}^n \tilde r_i \varphi_0(X_i)
    =
    \sum_{i=1}^n\sum_{j=1}^m \tilde\gamma_{ij}\varphi_0(X_i),
    \qquad
    \sum_{j=1}^m \tilde s_j \psi_0(Y_j)
    =
    \sum_{i=1}^n\sum_{j=1}^m \tilde\gamma_{ij}\psi_0(Y_j),
\end{align}
we obtain the exact identity \eqref{eq:uot_exact_empirical_excess}. Nonnegativity follows because the dual constraint gives
\begin{align}
    c(X_i,Y_j)-\varphi_0(X_i)-\psi_0(Y_j)\ge 0
\end{align}
for every $i,j$, and {both discrete $D_\KL$ terms are nonnegative.}

Finally, taking expectations in \eqref{eq:uot_exact_empirical_excess} and using the i.i.d.\ sampling model together with the oracle masses, we have
\begin{align}
    \Ep\left[\int \zeta_0   d\tilde\mu_n\right]
    = \int \zeta_0   d\mu,
    \qquad
    \Ep\left[\int \xi_0   d\tilde\nu_m\right]
    = \int \xi_0   d\nu.
\end{align}
Since we have the equality $\mathrm{UOT}(\mu,\nu)=\int \zeta_0   d\mu + \int \xi_0   d\nu$ by the duality, we obtain  \eqref{eq:uot_exact_empirical_excess_expectation}.
\end{proof}


\begin{corollary}[{Bound on $\Ep[\Delta_{nm}^{\mathrm{tr}}]$}]
\label{cor:uot_prop13_one_sided}
Assume that Assumption~\ref{assm:curvature} holds.
Let $\tilde\gamma=(\tilde\gamma_{ij})$ be any optimizer of $\mathrm{UOT}(\tilde\mu_n,\tilde\nu_m)$, and define
\begin{align}
    \Delta_{nm}^{\mathrm{tr}}
    \coloneqq
    \sum_{i=1}^n\sum_{j=1}^m
    \tilde\gamma_{ij}
    \|T_0(X_i)-Y_j\|^2.
\end{align}
Then, almost surely,
\begin{align}
\label{eq:uot_prop13_one_sided_as}
    \frac{\kappa}{2}\Delta_{nm}^{\mathrm{tr}}
    \le
    \mathrm{UOT}(\tilde\mu_n,\tilde\nu_m)
    - \int \zeta_0   d\tilde\mu_n
    - \int \xi_0   d\tilde\nu_m.
\end{align}
Consequently,
\begin{align}
\label{eq:uot_prop13_one_sided_expectation}
    \Ep[\Delta_{nm}^{\mathrm{tr}}]
    \le
    \frac{2}{\kappa}
    \Ep\Bigl[
        \mathrm{UOT}(\tilde\mu_n,\tilde\nu_m) - \mathrm{UOT}(\mu,\nu)
    \Bigr].
\end{align}
If, in addition, Assumptions~\ref{assm:curvature}, \ref{assm:supp_global}, and \ref{assm:positivity} hold, then with $C_\Lambda$ from Proposition~\ref{prop:two_sample_stability},
\begin{align}
\label{eq:uot_prop13_one_sided_rate}
    \Ep[\Delta_{nm}^{\mathrm{tr}}]
    \le
    \frac{2C_\Lambda}{\kappa}
    \left(
        \Ep[M_\mu W_2^2(\bar\mu_n,\bar\mu)] + \Ep[M_\nu W_2^2(\bar\nu_m,\bar\nu)]
    \right).
\end{align}
\end{corollary}


\begin{proof}[Proof of Corollary \ref{cor:uot_prop13_one_sided}]
By Proposition~\ref{prop:uot_exact_empirical_excess},
\begin{align}
\begin{aligned}
&\mathrm{UOT}(\tilde\mu_n,\tilde\nu_m)
- \int \zeta_0  d\tilde\mu_n
- \int \xi_0  d\tilde\nu_m \\
&\qquad=
\sum_{i=1}^n\sum_{j=1}^m
\tilde\gamma_{ij}
\Bigl(
c(X_i,Y_j)-\varphi_0(X_i)-\psi_0(Y_j)
\Bigr)
+ {D_\KL}(\tilde r \mid r^\star)
+ {D_\KL}(\tilde s \mid s^\star).
\end{aligned}
\end{align}
Lemma~\ref{lem:uot_gap_sufficient} and the nonnegativity of the {discrete $D_\KL$ terms} imply
\begin{align}
    \mathrm{UOT}(\tilde\mu_n,\tilde\nu_m)
    - \int \zeta_0  d\tilde\mu_n
    - \int \xi_0  d\tilde\nu_m
    &\ge
    \sum_{i=1}^n\sum_{j=1}^m
    \tilde\gamma_{ij}\frac{\kappa}{2}\|Y_j-T_0(X_i)\|^2 \\
    &= \frac{\kappa}{2}\Delta_{nm}^{\mathrm{tr}},
\end{align}
which proves \eqref{eq:uot_prop13_one_sided_as}. Taking expectations and using \eqref{eq:uot_exact_empirical_excess_expectation} yields \eqref{eq:uot_prop13_one_sided_expectation}.

Finally, Proposition~\ref{prop:two_sample_stability} applied to $(\tilde\mu_n,\tilde\nu_m)$ gives
\begin{align}
    &\mathrm{UOT}(\tilde\mu_n,\tilde\nu_m)-\mathrm{UOT}(\mu,\nu)\\
    &\le
    \int \zeta_0   d(\tilde\mu_n-\mu)
    + \int \xi_0   d(\tilde\nu_m-\nu)
    + C_\Lambda\Bigl(
        M_\mu W_2^2(\bar\mu_n,\bar\mu)+M_\nu W_2^2(\bar\nu_m,\bar\nu)
    \Bigr).
\end{align}
Taking expectations, the linear terms vanish by unbiasedness of the oracle empirical measures, and \eqref{eq:uot_prop13_one_sided_rate} follows.
\end{proof}

\begin{remark}
Corollary~\ref{cor:uot_prop13_one_sided} is the closest UOT analogue of the balanced case in \cite[Proposition~13]{manole2024plugin}.
The important difference from the balanced case is that the empirical excess objective contains, in addition to the transport mismatch $\Delta_{nm}^{\mathrm{tr}}$, {two positive $D_\KL$ terms encoding estimation error of the active source and target masses.}
Because of these extra terms, the argument above yields a clean \emph{one-sided} control of $\Ep[\Delta_{nm}^{\mathrm{tr}}]$ by the excess UOT objective, but not a reverse inequality without further information on the fitted empirical marginals.
\end{remark}

\subsection{From $\Delta_{nm}^{\mathrm{tr}}$ to barycentric and 1NN map errors}

We now show how Corollary~\ref{cor:uot_prop13_one_sided} feeds directly into concrete estimators of the transport map.
The first step is to pass from the row sums of $\tilde\gamma$ to the barycentric projection $\tilde T_i$.
The second step is to observe that the one-nearest-neighbor extension inherits the same in-sample error on the empirical active source measure.

\begin{theorem}[Oracle barycentric and 1NN bounds]
\label{thm:uot_barycentric_and_1nn_in_sample}
Assume that Assumption~\ref{assm:curvature} holds.
For each $i$ with $\tilde r_i>0$, define the barycentric projection of the $i$th row of $\tilde\gamma$ by
\begin{align}
    \tilde T_i \coloneqq \frac{1}{\tilde r_i}\sum_{j=1}^m \tilde\gamma_{ij}Y_j,
\end{align}
and set $\tilde T_i=X_i$ when $\tilde r_i=0$.
Define
\begin{align}
    \Delta_{nm}^{\mathrm{bar}}
    \coloneqq
    \sum_{i=1}^n \tilde r_i  \|\tilde T_i-T_0(X_i)\|^2.
\end{align}
Then, almost surely,
\begin{align}
\label{eq:uot_barycentric_delta_leq_transport_delta}
    \Delta_{nm}^{\mathrm{bar}}
    \le
    \Delta_{nm}^{\mathrm{tr}}.
\end{align}
Consequently, it holds that
\begin{align}
\label{eq:uot_barycentric_expectation_bound}
    \Ep\bigl[\Delta_{nm}^{\mathrm{bar}}\bigr]
    &\le
    \frac{2}{\kappa}
    \Ep\Bigl[
        \mathrm{UOT}(\tilde\mu_n,\tilde\nu_m)-\mathrm{UOT}(\mu,\nu)
    \Bigr],
\end{align}
and, if Assumptions~\ref{assm:curvature}-\ref{assm:positivity} hold, then
\begin{align}
\label{eq:uot_barycentric_rate_bound}
    \Ep\bigl[\Delta_{nm}^{\mathrm{bar}}\bigr]
    &\le
    \frac{2C_\Lambda}{\kappa}
    \left(
        \Ep[M_\mu W_2^2(\bar\mu_n,\bar\mu)] + \Ep[M_\nu W_2^2(\bar\nu_m,\bar\nu)]
    \right).
\end{align}

Now let $\tilde T^{\mathrm{1NN}}$ be the Voronoi extension defined above, and let
\begin{align}
    \tilde\gamma_0 \coloneqq \sum_{i=1}^n \tilde r_i   \delta_{X_i}
\end{align}
be the fitted empirical active source measure.
Then
 \begin{align}
 \label{eq:uot_1nn_identity_active_empirical}
    \int_\Omega \|\tilde T^{\mathrm{1NN}}(x)-T_{0,n}^{\mathrm{1NN}}(x)\|^2  d\tilde\gamma_0(x)
     =
     \Delta_{nm}^{\mathrm{bar}}.
 \end{align}
Hence the same expectation bounds \eqref{eq:uot_barycentric_expectation_bound}-\eqref{eq:uot_barycentric_rate_bound} hold for the in-sample error of the 1NN extension measured against $T_{0,n}^{\mathrm{1NN}}$ on $\tilde\gamma_0$.
 \end{theorem}

\begin{proof}[Proof of Theorem \ref{thm:uot_barycentric_and_1nn_in_sample}]
Fix $i\in\{1,\dots,n\}$.
If $\tilde r_i=0$, then the $i$th contribution to $\Delta_{nm}^{\mathrm{bar}}$ is zero and there is nothing to prove.
Assume therefore that $\tilde r_i>0$.
By the barycentric formula for the $i$th row of $\tilde\gamma$,
\begin{align}
    \tilde T_i-T_0(X_i)
    =
    \frac{1}{\tilde r_i}
    \sum_{j=1}^m \tilde\gamma_{ij}\bigl(Y_j-T_0(X_i)\bigr).
\end{align}
Since $z\mapsto \|z\|^2$ is convex, Jensen's inequality gives
\begin{align}
    \|\tilde T_i-T_0(X_i)\|^2
    &\le
    \frac{1}{\tilde r_i}
    \sum_{j=1}^m \tilde\gamma_{ij}  \|Y_j-T_0(X_i)\|^2.
\end{align}
Multiplying by $\tilde r_i$ and summing over $i$ yields
\begin{align}
    \Delta_{nm}^{\mathrm{bar}}
    = \sum_{i=1}^n \tilde r_i \|\tilde T_i-T_0(X_i)\|^2 \le \sum_{i=1}^n \sum_{j=1}^m \tilde\gamma_{ij}  \|Y_j-T_0(X_i)\|^2
    = \Delta_{nm}^{\mathrm{tr}},
\end{align}
which proves \eqref{eq:uot_barycentric_delta_leq_transport_delta}.
The expectation bounds \eqref{eq:uot_barycentric_expectation_bound} and \eqref{eq:uot_barycentric_rate_bound} now follow immediately from Corollary~\ref{cor:uot_prop13_one_sided}.

For the 1NN claim, note that the support of $\tilde\gamma_0$ is contained in $\{X_1,\dots,X_n\}$ and, for each $i$, one has
\begin{align}
    \tilde T^{\mathrm{1NN}}(X_i)=\tilde T_i,
    \qquad
    T_{0,n}^{\mathrm{1NN}}(X_i)=T_0(X_i).
\end{align}
Therefore,
\begin{align}
    \int_\Omega \|\tilde T^{\mathrm{1NN}}(x)-T_{0,n}^{\mathrm{1NN}}(x)\|^2  d\tilde\gamma_0(x)
    &= \sum_{i=1}^n \tilde r_i  \|\tilde T^{\mathrm{1NN}}(X_i)-T_{0,n}^{\mathrm{1NN}}(X_i)\|^2 \\
    &= \sum_{i=1}^n \tilde r_i  \|\tilde T_i-T_0(X_i)\|^2
    = \Delta_{nm}^{\mathrm{bar}},
\end{align}
which is \eqref{eq:uot_1nn_identity_active_empirical}.
The last sentence follows by substituting this identity into \eqref{eq:uot_barycentric_expectation_bound}-\eqref{eq:uot_barycentric_rate_bound}.
\end{proof}

Theorem~\ref{thm:uot_barycentric_and_1nn_in_sample} is deliberately stated for the oracle empirical plan $\tilde\gamma$ so that it matches Corollary~\ref{cor:uot_prop13_one_sided} without any additional bookkeeping.
For the actual fitted empirical plan based on $(\hat\mu_n,\hat\nu_m)$, one obtains the following estimated-mass variant.

\begin{proposition}[Bounds on $\hat\Delta_{nm}^{\mathrm{tr}}$ and $\hat\Delta_{nm}^{\mathrm{bar}}$]
\label{prop:uot_barycentric_and_1nn_estimated_mass}
Assume that Assumptions~\ref{assm:curvature}, \ref{assm:supp_global}, \ref{assm:positivity}, and \ref{assm:masses} hold.
Let $C_\Lambda$ be the stability constant from Proposition~\ref{prop:two_sample_stability}.
Define the empirical active measures
\begin{align}
    \hat r_i^\star \coloneqq e^{-\varphi_0(X_i)}\hat\mu_i,
    \qquad
    \hat s_j^\star \coloneqq e^{-\psi_0(Y_j)}\hat\nu_j,
\end{align}
and the fitted empirical transport and barycentric errors
\begin{align}
    \hat\Delta_{nm}^{\mathrm{tr}}
    &\coloneqq
    \sum_{i=1}^n\sum_{j=1}^m
    \hat\gamma_{ij}\|T_0(X_i)-Y_j\|^2, \\
    \hat\Delta_{nm}^{\mathrm{bar}}
    &\coloneqq
    \sum_{i=1}^n \hat r_i \|\hat T_i-T_0(X_i)\|^2.
\end{align}
Then, almost surely,
\begin{equation}
\label{eq:uot_hat_exact_empirical_excess}
\begin{aligned}
&\mathrm{UOT}(\hat\mu_n,\hat\nu_m)
- \int \zeta_0   d\hat\mu_n
- \int \xi_0   d\hat\nu_m \\
&\qquad=
\sum_{i=1}^n\sum_{j=1}^m
\hat\gamma_{ij}
\Bigl(
c(X_i,Y_j)-\varphi_0(X_i)-\psi_0(Y_j)
\Bigr)
+ {D_\KL}(\hat r \mid \hat r^\star)
+ {D_\KL}(\hat s \mid \hat s^\star).
\end{aligned}
\end{equation}
In particular,
\begin{align}
\label{eq:uot_hat_barycentric_delta_bound}
    \hat\Delta_{nm}^{\mathrm{bar}}
    \le
    \hat\Delta_{nm}^{\mathrm{tr}}
    \le
    \frac{2}{\kappa}
    \left[
        \mathrm{UOT}(\hat\mu_n,\hat\nu_m)
        - \int \zeta_0   d\hat\mu_n
        - \int \xi_0   d\hat\nu_m
    \right].
\end{align}
Now let
\begin{align}
    \hat T^{\mathrm{1NN}}(x)
    \coloneqq
    \sum_{i=1}^n \mathbf{1}\{x\in V_i\}\hat T_i,
    \qquad
    \hat\gamma_0 \coloneqq \sum_{i=1}^n \hat r_i \delta_{X_i}.
\end{align}
Then
\begin{align}
\label{eq:uot_hat_1nn_identity_active_empirical}
    \int_\Omega \|\hat T^{\mathrm{1NN}}(x)-T_{0,n}^{\mathrm{1NN}}(x)\|^2   d\hat\gamma_0(x)
    =
    \hat\Delta_{nm}^{\mathrm{bar}}.
\end{align}
Moreover, there exist constants $C_{\mathrm{mass}}>0$ and
$C_{\mathrm{log}}\coloneqq e\sqrt{M_\mu M_\nu}/\kappa$, depending only on
$M_\mu,M_\nu,\|\zeta_0\|_\infty,\|\xi_0\|_\infty$, such that the following bound holds almost surely:
\begin{align}
\label{eq:uot_hat_barycentric_rate_bound}
\begin{aligned}
    \hat\Delta_{nm}^{\mathrm{bar}}
    &\le
    \frac{2C_\Lambda}{\kappa}
    \left(
        M_\mu W_2^2(\bar\mu_n,\bar\mu)+M_\nu W_2^2(\bar\nu_m,\bar\nu)
    \right)
    + \frac{2C_{\mathrm{mass}}}{\kappa}
    \left(
        |\hat M_\mu-M_\mu|+|\hat M_\nu-M_\nu|
    \right) \\
    &\quad+ 2C_{\mathrm{log}}\bigl(|\log\alpha_n|+|\log\beta_m|\bigr) \\
    &\quad - \frac{2}{\kappa} \int\zeta_0 d(\hat\mu_n-\tilde\mu_n)
            - \frac{2}{\kappa} \int\xi_0 d(\hat\nu_m-\tilde\nu_m) \\
    &\quad - \frac{2}{\kappa} \int\zeta_0 d(\tilde\mu_n-\mu)
            - \frac{2}{\kappa} \int\xi_0 d(\tilde\nu_m-\nu)
\end{aligned}
\end{align}
Consequently, taking expectations and using
Assumption~\ref{assm:masses} (in particular the a.s.\ lower bound $\hat M_\mu\ge cM_\mu$), there exists a
constant $\widetilde C_{\mathrm{mass}}>0$ such that, for all sufficiently large $n,m$,
\begin{align}
\label{eq:uot_hat_barycentric_expectation_bound}
    \Ep[\hat\Delta_{nm}^{\mathrm{bar}}]
    &\le
    \frac{2C_\Lambda}{\kappa}
    \left(
        \Ep[M_\mu W_2^2(\bar\mu_n,\bar\mu)] + \Ep[M_\nu W_2^2(\bar\nu_m,\bar\nu)]
    \right)
    + \frac{2\widetilde C_{\mathrm{mass}}}{\kappa}\bigl(a_n+b_m\bigr).
\end{align}
The same bound therefore holds for the in-sample $1$NN error
\(
\int_\Omega \|\hat T^{\mathrm{1NN}}-T_{0,n}^{\mathrm{1NN}}\|^2 d\hat\gamma_0
\)
by \eqref{eq:uot_hat_1nn_identity_active_empirical}.
\end{proposition}

\begin{proof}[Proof of Proposition \ref{prop:uot_barycentric_and_1nn_estimated_mass}]
The algebra in Proposition~\ref{prop:uot_exact_empirical_excess} does not use that the reference masses are the true masses.
Therefore, replacing
\(
(\tilde\mu_n,\tilde\nu_m,\tilde\gamma,\tilde r,\tilde s,r^\star,s^\star)
\)
by
\(
(\hat\mu_n,\hat\nu_m,\hat\gamma,\hat r,\hat s,\hat r^\star,\hat s^\star)
\)
in that proof yields \eqref{eq:uot_hat_exact_empirical_excess}.

Lemma~\ref{lem:uot_gap_sufficient} and the nonnegativity of the {discrete $D_\KL$ terms} in \eqref{eq:uot_hat_exact_empirical_excess} imply
\begin{align}
    \frac{\kappa}{2}\hat\Delta_{nm}^{\mathrm{tr}}
    &\le
    \mathrm{UOT}(\hat\mu_n,\hat\nu_m)
    - \int \zeta_0   d\hat\mu_n
    - \int \xi_0   d\hat\nu_m.
\end{align}
Exactly the same Jensen argument as in Theorem~\ref{thm:uot_barycentric_and_1nn_in_sample}, now applied to the rows of $\hat\gamma$, gives
\begin{align}
    \hat\Delta_{nm}^{\mathrm{bar}} \le \hat\Delta_{nm}^{\mathrm{tr}},
\end{align}
which proves \eqref{eq:uot_hat_barycentric_delta_bound}.
The identity \eqref{eq:uot_hat_1nn_identity_active_empirical} is also the same computation as in Theorem~\ref{thm:uot_barycentric_and_1nn_in_sample}, with $\tilde\gamma_0$ replaced by $\hat\gamma_0$.

Next, by duality and \eqref{eq:estimated_mass_stability},
\begin{align}
\label{eq:uot_hat_excess_vs_oracle_bound}
\begin{aligned}
&\mathrm{UOT}(\hat\mu_n,\hat\nu_m)
- \int \zeta_0   d\hat\mu_n
- \int \xi_0   d\hat\nu_m \\
&\qquad=
\mathrm{UOT}(\hat\mu_n,\hat\nu_m)-\mathrm{UOT}(\mu,\nu)
- \int \zeta_0   d(\hat\mu_n-\mu)
- \int \xi_0   d(\hat\nu_m-\nu) \\
&\qquad\le
C_\Lambda\Bigl(
    M_\mu W_2^2(\bar\mu_n,\bar\mu)+M_\nu W_2^2(\bar\nu_m,\bar\nu)
\Bigr)
+ \bigl( |\hat M_\mu-M_\mu|+|\hat M_\nu-M_\nu| \bigr) \\
&\qquad\quad
+ e\sqrt{M_\mu M_\nu}\bigl(|\log\alpha_n|+|\log\beta_m|\bigr) \\
&\qquad\quad
- \int \zeta_0   d(\hat\mu_n-\tilde\mu_n)
- \int \xi_0   d(\hat\nu_m-\tilde\nu_m) \\
&\qquad\quad
- \int \zeta_0   d(\tilde\mu_n-\mu)
- \int \xi_0   d(\tilde\nu_m-\nu),
\end{aligned}
\end{align}
where we used $\mathrm{UOT}(\mu,\nu)=\int\zeta_0 d\mu+\int\xi_0 d\nu$ and rearranged the linear-in-$(\tilde\mu_n,\tilde\nu_m)$ terms produced by \eqref{eq:estimated_mass_stability}.
Since $\hat\mu_n-\tilde\mu_n$ and $\hat\nu_m-\tilde\nu_m$ are signed atomic measures carried by the observed supports,
\begin{align}
    \left|
        \int \zeta_0   d(\hat\mu_n-\tilde\mu_n)
    \right|
    \le
    \|\zeta_0\|_\infty |\hat M_\mu-M_\mu|,
    \qquad
    \left|
        \int \xi_0   d(\hat\nu_m-\tilde\nu_m)
    \right|
    \le
    \|\xi_0\|_\infty |\hat M_\nu-M_\nu|.
\end{align}
Substituting these bounds into \eqref{eq:uot_hat_excess_vs_oracle_bound}, multiplying by $2/\kappa$ via \eqref{eq:uot_hat_barycentric_delta_bound}, and absorbing constants proves the bound \eqref{eq:uot_hat_barycentric_rate_bound}.

Taking expectations of \eqref{eq:uot_hat_barycentric_rate_bound}, the terms $\Ep[\int\zeta_0 d(\tilde\mu_n-\mu)]$ and $\Ep[\int\xi_0 d(\tilde\nu_m-\nu)]$ vanish. By Assumption~\ref{assm:masses}, $\hat M_\mu\ge cM_\mu$ a.s.\ gives $|\log\alpha_n|\le|\alpha_n-1|/c$ a.s., hence $M_\mu\Ep[|\log\alpha_n|]\le \Ep[|\hat M_\mu-M_\mu|]/c\le a_n/c$, and analogously for the $\nu$-side, so
\(
e\sqrt{M_\mu M_\nu} \Ep[|\log\alpha_n|]\le c^{-1}e\sqrt{M_\nu/M_\mu} a_n
\)
and similarly $e\sqrt{M_\mu M_\nu}\Ep[|\log\beta_m|]\le c^{-1}e\sqrt{M_\mu/M_\nu} b_m$.
Setting $\widetilde C_{\mathrm{mass}}\coloneqq C_{\mathrm{mass}}+c^{-1}e\sqrt{M_\mu/M_\nu}+c^{-1}e\sqrt{M_\nu/M_\mu}$ proves \eqref{eq:uot_hat_barycentric_expectation_bound}.
The last sentence is immediate from \eqref{eq:uot_hat_1nn_identity_active_empirical}.
\end{proof}




\subsection{From the in-sample error to the population error}

To pass from the in-sample barycentric error to the population error of the 1NN extension under $\gamma_0$, we compare the random Voronoi cell masses with the fitted row masses, keeping the row-marginal $D_\KL$ error ${D_\KL}(\tilde r\mid r^\star)$ explicitly in the argument.



\begin{lemma}[Bounds on $R_n$ and $M_n$]
\label{lem:uot_voronoi_mass_radius}
Let $V_1,\dots,V_n$ be the Voronoi cells induced by the source sample $X_1,\dots,X_n$, and define
\begin{align*}
    M_n \coloneqq \max_{1\le i\le n} \bar\mu(V_i),
    \qquad
    R_n \coloneqq \max_{1\le i\le n}\sup_{x\in V_i}\|x-X_i\|.
\end{align*}
Assume Assumptions~\ref{assm:supp_global} and \ref{assm:positivity}. Then there exist constants $C_1,C_2>0$, depending only on $d$, $\beta_{\min}$, $\beta_{\max}$, $\epsilon_0$, and $\delta_0$, such that:
\begin{enumerate}[label=(\roman*)]
    \item for every $\delta\in(0,1)$,
    \begin{align}
\label{eq:uot_voronoi_maxmass_hp}
        \mathbb{P} \left(
            M_n
            \ge
            \frac{C_1}{n}\bigl[d\log n+\log(1/\delta)\bigr]
        \right)
        \le \delta;
    \end{align}
    \item
    \begin{align}
\label{eq:uot_voronoi_radius_moment}
        \Ep[R_n^2]
        \le
        C_2\left({(\log n)/n}\right)^{2/d}.
    \end{align}
\end{enumerate}
\end{lemma}


\begin{proof}[Proof of Lemma \ref{lem:uot_voronoi_mass_radius}]
Let $\bar p$ denote the density of $\bar\mu$.
Then $\bar p=p/M_\mu$ is bounded above and below by positive constants on $\Omega$.
Moreover, Assumption~\ref{assm:supp_global} states exactly that $\Omega$ is compact and satisfies the interior cone condition.
Hence the normalized sample $X_1,\dots,X_n\stackrel{\mathrm{i.i.d.}}{\sim}\bar\mu$ satisfies the hypotheses of \cite[Lemma~40]{manole2024plugin}.
Applying that lemma gives constants $C_1,C_2>0$ such that, for every $\delta\in(0,1)$,
\begin{align*}
    \mathbb{P}\left(
        M_n
        \ge
        \frac{C_1}{n}\bigl[d\log n+\log(1/\delta)\bigr]
    \right)
    \le
    \delta,
\end{align*}
and
\begin{align*}
    \Ep[R_n^2]
    \le
    C_2\left({(\log n)/n}\right)^{2/d}.
\end{align*}
These are precisely \eqref{eq:uot_voronoi_maxmass_hp} and \eqref{eq:uot_voronoi_radius_moment}.
\end{proof}



\begin{lemma}[Bound on $\sum_i\hat r_i^\star\hat b_i$]
\label{lem:uot_weighted_pinsker_rows_hat}
Define
\begin{align*}
    \hat\Delta_{nm}^{\mathrm{bar}}
    \coloneqq
    \sum_{i=1}^n \hat r_i \|\hat T_i-T_0(X_i)\|^2,
    \qquad
    \hat b_i \coloneqq \|\hat T_i-T_0(X_i)\|^2,
    \qquad
    C_\Omega \coloneqq \mathrm{diam}(\Omega)^2.
\end{align*}
Then, almost surely,
\begin{align}
    \sum_{i=1}^n \hat r_i^\star \hat b_i
    \le
    3\hat\Delta_{nm}^{\mathrm{bar}}
    + 2C_\Omega {D_\KL}(\hat r\mid \hat r^\star).
\end{align}
\end{lemma}


\begin{proof}[Proof of Lemma \ref{lem:uot_weighted_pinsker_rows_hat}]


Set
\begin{align}
    X_n^\star \coloneqq \sum_{i=1}^n \hat r_i^\star \hat b_i,
    \qquad
    Y_n \coloneqq \sum_{i=1}^n \hat r_i \hat b_i
    = \hat\Delta_{nm}^{\mathrm{bar}}.
\end{align}
We use the elementary scalar inequality
\begin{align}
\label{eq:scalar_kl_quadratic}
    u\log \left(\frac{u}{v}\right)-u+v
    \ge
    \frac{(u-v)^2}{2(u+v)},
    \qquad u,v\ge 0,
\end{align}
with the usual convention $0\log 0=0$.
Summing \eqref{eq:scalar_kl_quadratic} over $i$ gives
\begin{align}
\label{eq:rows_weighted_pinsker_basic}
    \sum_{i=1}^n \frac{(\hat r_i-\hat r_i^\star)^2}{\hat r_i+\hat r_i^\star}
    \le
    2 {D_\KL}(\hat r\mid \hat r^\star).
\end{align}
Since $0\le \hat b_i\le C_\Omega$, the Cauchy-Schwarz inequality and \eqref{eq:rows_weighted_pinsker_basic} imply
\begin{align}
\begin{aligned}
    X_n^\star-Y_n
    &= \sum_{i=1}^n (\hat r_i^\star-\hat r_i)\hat b_i \\
    &\le
    \left(
        \sum_{i=1}^n \frac{(\hat r_i-\hat r_i^\star)^2}{\hat r_i+\hat r_i^\star}
    \right)^{1/2}
    \left(
        \sum_{i=1}^n (\hat r_i+\hat r_i^\star)\hat b_i^2
    \right)^{1/2} \\
    &\le
    \sqrt{2{D_\KL}(\hat r\mid \hat r^\star)}
    \left(
        C_\Omega \sum_{i=1}^n (\hat r_i+\hat r_i^\star)\hat b_i
    \right)^{1/2} \\
    &=
    \sqrt{2C_\Omega {D_\KL}(\hat r\mid \hat r^\star)}
      \sqrt{X_n^\star+Y_n}.
\end{aligned}
\end{align}
Using $ab\le {1/2} a^2+{1/2} b^2$ with $a=\sqrt{2C_\Omega {D_\KL}(\hat r\mid \hat r^\star)}$ and $b=\sqrt{X_n^\star+Y_n}$, we obtain
\begin{align}
    X_n^\star-Y_n
    \le
    C_\Omega {D_\KL}(\hat r\mid \hat r^\star)
    + \tfrac{1}{2}(X_n^\star+Y_n).
\end{align}
Rearranging yields $X_n^\star \le 3Y_n + 2C_\Omega {D_\KL}(\hat r\mid \hat r^\star)$, which is the claim.

\end{proof}


\begin{theorem}[Map risk of $\hat T^{1\mathrm{NN}}$]
\label{thm:uot_complete_1nn_estimated_mass}
Assume that Assumptions~\ref{assm:curvature}, \ref{assm:supp_global}, \ref{assm:positivity}, and \ref{assm:masses} hold.
By Assumption~\ref{assm:curvature}, the map $T_0(x) = x - \nabla \varphi_0(x)$ is Lipschitz with constant $L_T \coloneqq \kappa^{-1}$.
Let
\begin{align}
    w(x)\coloneqq e^{-\varphi_0(x)},
    \qquad
    w_- \coloneqq \inf_{x\in\Omega} w(x),
    \qquad
    w_+ \coloneqq \sup_{x\in\Omega} w(x),
\end{align}
and define the mass-accuracy event
\begin{align}
    \mathcal{A}_{n,m}
    \coloneqq
    \left\{
        |\hat M_\mu-M_\mu|\le \frac{M_\mu}{2},
        \ |\hat M_\nu-M_\nu|\le \frac{M_\nu}{2}
    \right\}.
\end{align}
Then, on $\mathcal{A}_{n,m}$, almost surely,
\begin{align}
\label{eq:uot_complete_1nn_estimated_mass_as}
\begin{aligned}
    \int_\Omega \|\hat T^{\mathrm{1NN}}(x)-T_0(x)\|^2  d\gamma_0(x)
    &\le
    4\frac{w_+}{w_-}  n M_n
    \left(
        3\hat\Delta_{nm}^{\mathrm{bar}}
        + 2C_\Omega {D_\KL}(\hat r\mid \hat r^\star)
    \right) \\
    &\quad + 2L_T^2\gamma_0(\Omega)R_n^2.
\end{aligned}
\end{align}
Consequently, if
\(
\Ep[|\hat M_\mu-M_\mu|+|\hat M_\nu-M_\nu|]<\infty
\),
then there exists a constant $C>0$, depending only on
\(
d,\beta_{\min},\beta_{\max},\epsilon_0,\delta_0,M_\mu,M_\nu,\|\varphi_0\|_\infty,\|\zeta_0\|_\infty,\|\xi_0\|_\infty,L_T,\Omega
\),
such that for all sufficiently large $n,m$,
\begin{align}
\label{eq:uot_complete_1nn_estimated_mass_expectation}
\begin{aligned}
    \Ep\Bigl[
        \int_\Omega \|\hat T^{\mathrm{1NN}}(x)-T_0(x)\|^2  d\gamma_0(x)
    \Bigr]
    &\le
    C \log n
    \Bigl(
        \Ep[M_\mu W_2^2(\bar\mu_n,\bar\mu)]
        +
        \Ep[M_\nu W_2^2(\bar\nu_m,\bar\nu)]
        + a_n + b_m
    \Bigr) \\
    &\quad
    + C\left({(\log n)/n}\right)^{2/d}
    + C{(\log n)/n},
\end{aligned}
\end{align}
where $a_n,b_m$ are the rates from Assumption~\ref{assm:masses}.
\end{theorem}




\begin{proof}[Proof of Theorem \ref{thm:uot_complete_1nn_estimated_mass}]
Let
\begin{align}
    \hat b_i \coloneqq \|\hat T_i-T_0(X_i)\|^2.
\end{align}
For $x\in V_i$, one has $\hat T^{\mathrm{1NN}}(x)=\hat T_i$, hence
\begin{align}
    \|\hat T^{\mathrm{1NN}}(x)-T_0(x)\|^2
    \le
    2\hat b_i + 2\|T_0(X_i)-T_0(x)\|^2
    \le
    2\hat b_i + 2L_T^2\|X_i-x\|^2.
\end{align}
Integrating over $V_i$ with respect to $\gamma_0$ and summing over $i$ yields
\begin{align}
\label{eq:uot_complete_1nn_estimated_mass_start}
    \int_\Omega \|\hat T^{\mathrm{1NN}}(x)-T_0(x)\|^2  d\gamma_0(x)
    \le
    2\sum_{i=1}^n \gamma_0(V_i)\hat b_i
    + 2L_T^2\gamma_0(\Omega)R_n^2.
\end{align}
On $\mathcal{A}_{n,m}$, we have $\hat M_\mu\ge M_\mu/2$.
Therefore,
\begin{align}
    \gamma_0(V_i)
    = \int_{V_i} w  d\mu
    \le
    w_+ \mu(V_i)
    = M_\mu w_+\bar\mu(V_i)
    \le
    M_\mu w_+ M_n,
\end{align}
whereas
\begin{align}
    \hat r_i^\star = w(X_i){\hat M_\mu/n}
    \ge
    w_-\frac{M_\mu}{2n}.
\end{align}
Hence, on $\mathcal{A}_{n,m}$,
\begin{align}
    \gamma_0(V_i)
    \le
    2\frac{w_+}{w_-}  nM_n \hat r_i^\star.
\end{align}
Substituting this into \eqref{eq:uot_complete_1nn_estimated_mass_start} gives
\begin{align}
    \int_\Omega \|\hat T^{\mathrm{1NN}}(x)-T_0(x)\|^2  d\gamma_0(x)
    \le
    4\frac{w_+}{w_-}  nM_n \sum_{i=1}^n \hat r_i^\star\hat b_i
    + 2L_T^2\gamma_0(\Omega)R_n^2.
\end{align}
Applying Lemma~\ref{lem:uot_weighted_pinsker_rows_hat} proves
\eqref{eq:uot_complete_1nn_estimated_mass_as}.

For the expectation bound, note first that $\hat T_i\in\Omega$ for every $i$:
if $\hat r_i=0$, then $\hat T_i=X_i\in\Omega$, while if $\hat r_i>0$, then $\hat T_i$ is a convex combination of $Y_1,\dots,Y_m\in\Omega$ and $\Omega$ is convex.
Since also $T_0(x)\in\Omega$ for every $x\in\Omega$,
\begin{align}
    \int_\Omega \|\hat T^{\mathrm{1NN}}(x)-T_0(x)\|^2  d\gamma_0(x)
    \le
    C_\Omega \gamma_0(\Omega)
\end{align}
almost surely.
Therefore,
\begin{align}
\label{eq:uot_complete_1nn_estimated_mass_expect_split}
    \Ep\Bigl[
        \int_\Omega \|\hat T^{\mathrm{1NN}}-T_0\|^2  d\gamma_0
    \Bigr]
    \le
    \Ep\Bigl[
        \mathbf{1}_{\mathcal{A}_{n,m}}
        \int_\Omega \|\hat T^{\mathrm{1NN}}-T_0\|^2  d\gamma_0
    \Bigr]
    + C_\Omega \gamma_0(\Omega)\mathbb{P}(\mathcal{A}_{n,m}^c).
\end{align}
By Markov's inequality,
\begin{align}
\label{eq:uot_complete_1nn_estimated_mass_badmass_prob}
    \mathbb{P}(\mathcal{A}_{n,m}^c)
    \le
    \frac{2}{M_\mu}\Ep[|\hat M_\mu-M_\mu|]
    +
    \frac{2}{M_\nu}\Ep[|\hat M_\nu-M_\nu|].
\end{align}

Set
\begin{align}
    \hat Z_n
    \coloneqq
    3\hat\Delta_{nm}^{\mathrm{bar}}
    +2C_\Omega {D_\KL}(\hat r\mid \hat r^\star).
\end{align}
Then \eqref{eq:uot_complete_1nn_estimated_mass_as} yields
\begin{align}
\label{eq:uot_complete_1nn_estimated_mass_Z}
    \mathbf{1}_{\mathcal{A}_{n,m}}
    \int_\Omega \|\hat T^{\mathrm{1NN}}-T_0\|^2  d\gamma_0
    \le
    4\frac{w_+}{w_-}  nM_n \hat Z_n \mathbf{1}_{\mathcal{A}_{n,m}}
    + 2L_T^2\gamma_0(\Omega)R_n^2.
\end{align}
By Lemma~\ref{lem:uot_voronoi_mass_radius}, there exists $c>0$ such that
\begin{align}
    m_n \coloneqq c{(\log n)/n}
\end{align}
satisfies $\mathbb{P}(M_n\ge m_n)\le n^{-2}$ for all large $n$.
Hence
\begin{align}
\label{eq:uot_complete_1nn_estimated_mass_Mn_split}
\begin{aligned}
    \Ep[nM_n \hat Z_n \mathbf{1}_{\mathcal{A}_{n,m}}]
    &=
    \Ep[nM_n \hat Z_n \mathbf{1}_{\mathcal{A}_{n,m}}\mathbf{1}\{M_n<m_n\}] \\
    &\quad
    +
    \Ep[nM_n \hat Z_n \mathbf{1}_{\mathcal{A}_{n,m}}\mathbf{1}\{M_n\ge m_n\}].
\end{aligned}
\end{align}
The first term is bounded by
\begin{align}
\label{eq:uot_complete_1nn_estimated_mass_smallcells}
    \Ep[nM_n \hat Z_n \mathbf{1}_{\mathcal{A}_{n,m}}\mathbf{1}\{M_n<m_n\}]
    \le
    c\log n \Ep[\hat Z_n].
\end{align}

We next show that $\hat Z_n\mathbf{1}_{\mathcal{A}_{n,m}}\le C_0\log n$ for a deterministic constant $C_0$.
Let
\begin{align}
    \hat M_n^{row}\coloneqq \sum_{i=1}^n \hat r_i.
\end{align}
On $\mathcal{A}_{n,m}$, the zero plan is feasible for $\mathrm{UOT}(\hat\mu_n,\hat\nu_m)$ and therefore
\begin{align}
    \mathrm{UOT}(\hat\mu_n,\hat\nu_m)
    \le
    (\hat M_\mu+\hat M_\nu)
    \le
    \frac{3}{2}(M_\mu+M_\nu).
\end{align}
Since the transport term and the {column $D_\KL$ term} are nonnegative,
\begin{align}
    {D_\KL}(\hat r\mid \hat\mu)
    \le
    \frac{3}{2}(M_\mu+M_\nu).
\end{align}
By Jensen's inequality,
\begin{align}
    {D_\KL}(\hat r\mid \hat\mu)
    \ge
    \hat M_n^{row}\log\left(\frac{\hat M_n^{row}}{\hat M_\mu}\right)
    - \hat M_n^{row}+\hat M_\mu.
\end{align}
Because $\hat M_\mu\in[M_\mu/2,3M_\mu/2]$ on $\mathcal{A}_{n,m}$, there exists a deterministic constant $M_\star>0$, depending only on $M_\mu$ and $M_\nu$, such that
\begin{align}
    \hat M_n^{row}\mathbf{1}_{\mathcal{A}_{n,m}}
    \le
    M_\star.
\end{align}
Consequently,
\begin{align}
\label{eq:uot_complete_1nn_estimated_mass_bar_uniform}
    \hat\Delta_{nm}^{\mathrm{bar}}\mathbf{1}_{\mathcal{A}_{n,m}}
    \le
    C_\Omega M_\star.
\end{align}
Also, on $\mathcal{A}_{n,m}$,
\begin{align}
    \hat r_i^\star
    =
    w(X_i){\hat M_\mu/n}
    \ge
    \frac{w_-M_\mu}{2n},
    \qquad
    \sum_{i=1}^n \hat r_i^\star
    \le
    \frac{3}{2}M_\mu w_+.
\end{align}
Therefore,
\begin{align}
\label{eq:uot_complete_1nn_estimated_mass_kl_uniform}
\begin{aligned}
    {D_\KL}(\hat r\mid \hat r^\star)
    &=
    \sum_{i=1}^n
    \left[
        \hat r_i\log\left(\frac{\hat r_i}{\hat r_i^\star}\right)-\hat r_i+\hat r_i^\star
    \right] \\
    &\le
    \sum_{i=1}^n \hat r_i \log\left(\frac{\hat M_n^{row}}{w_-M_\mu/(2n)}\right)
    - \sum_{i=1}^n \hat r_i + \sum_{i=1}^n \hat r_i^\star \\
    &\le
    M_\star \log\left(\frac{2M_\star n}{w_-M_\mu}\right)
    + \frac{3}{2}M_\mu w_+.
\end{aligned}
\end{align}
Combining \eqref{eq:uot_complete_1nn_estimated_mass_bar_uniform} and
\eqref{eq:uot_complete_1nn_estimated_mass_kl_uniform}, we obtain
\begin{align}
    \hat Z_n\mathbf{1}_{\mathcal{A}_{n,m}}
    \le
    C_0\log n
\end{align}
for all large $n$.
Therefore,
\begin{align}
\label{eq:uot_complete_1nn_estimated_mass_largecells}
    \Ep[nM_n \hat Z_n \mathbf{1}_{\mathcal{A}_{n,m}}\mathbf{1}\{M_n\ge m_n\}]
    \le
    n \mathbb{P}(M_n\ge m_n) C_0\log n
    \le
    C{(\log n)/n}.
\end{align}
Combining \eqref{eq:uot_complete_1nn_estimated_mass_Mn_split},
\eqref{eq:uot_complete_1nn_estimated_mass_smallcells}, and
\eqref{eq:uot_complete_1nn_estimated_mass_largecells} yields
\begin{align}
\label{eq:uot_complete_1nn_estimated_mass_nMnZn}
    \Ep[nM_n \hat Z_n \mathbf{1}_{\mathcal{A}_{n,m}}]
    \le
    C\log n \Ep[\hat Z_n] + C{(\log n)/n}.
\end{align}

Next, Proposition~\ref{prop:uot_barycentric_and_1nn_estimated_mass} gives
\begin{align}
\label{eq:uot_complete_1nn_estimated_mass_bar_expect}
    \Ep[\hat\Delta_{nm}^{\mathrm{bar}}]
    \le
    C
    \Bigl(
        \Ep[M_\mu W_2^2(\bar\mu_n,\bar\mu)]
        +
        \Ep[M_\nu W_2^2(\bar\nu_m,\bar\nu)]
        + a_n + b_m
    \Bigr)
\end{align}
for a constant $C>0$ depending only on the model parameters, where $a_n,b_m$ are the rates from Assumption~\ref{assm:masses}.


Lemma~\ref{lem:uot_fitted_row_kl_expect} below provides the matching bound
\begin{align}
\label{eq:uot_complete_1nn_estimated_mass_kl_expect}
    \Ep[{D_\KL}(\hat r\mid \hat r^\star)]
    \le
    C
    \Bigl(
        \Ep[M_\mu W_2^2(\bar\mu_n,\bar\mu)]
        +
        \Ep[M_\nu W_2^2(\bar\nu_m,\bar\nu)]
        + a_n + b_m
    \Bigr)
\end{align}
for another constant $C>0$.
Substituting \eqref{eq:uot_complete_1nn_estimated_mass_bar_expect} and
\eqref{eq:uot_complete_1nn_estimated_mass_kl_expect} into
\eqref{eq:uot_complete_1nn_estimated_mass_nMnZn}, then combining
\eqref{eq:uot_complete_1nn_estimated_mass_expect_split},
\eqref{eq:uot_complete_1nn_estimated_mass_badmass_prob},
\eqref{eq:uot_complete_1nn_estimated_mass_Z},
\eqref{eq:uot_complete_1nn_estimated_mass_nMnZn}, and
Lemma~\ref{lem:uot_voronoi_mass_radius}(ii), proves
\eqref{eq:uot_complete_1nn_estimated_mass_expectation}.
\end{proof}



\subsection{Active-source and growth estimation for the 1NN extension}

The row marginals of a fitted UOT plan directly encode the active-source ratio
\(w_0=a_0^2=e^{-\varphi_0}\). They do not directly encode the Gaussian--Hellinger growth factor
\(
\lambda_0=w_0\exp(\|x-T_0(x)\|^2/4)
\).
Hence the proof separates two steps: first we estimate the active-source factor \(a_0\) from the row marginals, and then we transfer the active-factor and map errors to the corrected \(\lambda_0\)-risk.

\begin{lemma}[Square-root lower bound on $D_\KL$]
\label{lem:uot_kl_sqrt_lower_bound}
For every $a,b\ge 0$, with the conventions $0\log 0 = 0$ and $a\log(a/0)=+\infty$ for $a>0$, we have
\begin{align}
\label{eq:uot_kl_sqrt_lower_bound}
    a\log\left(\frac{a}{b}\right)-a+b
    \ge
    \left(\sqrt{a}-\sqrt{b}\right)^2.
\end{align}
\end{lemma}

\begin{proof}[Proof of Lemma \ref{lem:uot_kl_sqrt_lower_bound}]
If $b=0$, then \eqref{eq:uot_kl_sqrt_lower_bound} is immediate from the stated conventions, so we may assume that $b>0$.
Writing $a=tb$, it is enough to prove that
\begin{align*}
    t\log t - t + 1 \ge (\sqrt{t}-1)^2,
    \qquad t\ge 0.
\end{align*}
Set $s=\sqrt{t}$. Then
\begin{align*}
    t\log t-t+1-(\sqrt t-1)^2
    =2s^2\log s-2s^2+2s
    =2s(s\log s-s+1)\ge0,
\end{align*}
because $u\mapsto u\log u-u+1$ is nonnegative on $[0,\infty)$.
\end{proof}

\begin{lemma}[Pointwise bound on $|\hat\lambda_i-\lambda_0(X_i)|^2$]
\label{lem:discrete_gh_growth_transfer}
Let $x\in\Omega$, let $T,T^0\in\Omega$, and let $w^0\in[w_-,w_+]$. Define
\begin{align*}
    a^0=\sqrt{w^0},
    \qquad
    \lambda^0=w^0\exp\left(\frac14\|x-T^0\|^2\right).
\end{align*}
For any $u\ge0$, set
\begin{align*}
    \bar w=\clip_{[w_-,w_+]}(u),
    \qquad
    \bar\lambda=\bar w\exp\left(\frac14\|x-T\|^2\right).
\end{align*}
Then there exists a constant $C>0$, depending only on $w_-,w_+$ and $\operatorname{diam}(\Omega)$, such that
\begin{align}
\label{eq:discrete_gh_growth_transfer}
    |\bar\lambda-\lambda^0|^2
    \le
    C\left(|\sqrt{u}-a^0|^2+\|T-T^0\|^2\right).
\end{align}
\end{lemma}

\begin{proof}[Proof of Lemma \ref{lem:discrete_gh_growth_transfer}]
Since $w^0\in[w_-,w_+]$, clipping can only move $u$ closer to $w^0$; hence
\begin{align*}
    |\sqrt{\bar w}-\sqrt{w^0}|
    \le
    |\sqrt{u}-\sqrt{w^0}|.
\end{align*}
Moreover, $|\bar w-w^0|\le 2\sqrt{w_+}|\sqrt{\bar w}-\sqrt{w^0}|$. On the bounded domain, the map $T\mapsto\exp(\|x-T\|^2/4)$ is uniformly bounded and uniformly Lipschitz for $x,T\in\Omega$. Therefore
\begin{align*}
    |\bar\lambda-\lambda^0|
    &\le
    C|\bar w-w^0|+C\|T-T^0\| \\
    &\le
    C\left(|\sqrt{u}-a^0|+\|T-T^0\|\right),
\end{align*}
and squaring proves the claim.
\end{proof}

\begin{lemma}[Bound on $\int|\hat\lambda-\lambda_0|^2 d\hat\eta$]
\label{lem:continuous_gh_growth_transfer}
Let $\hat\eta$ be a finite positive measure on a bounded set $\Omega$, let $\hat a,a_0:\Omega\to[0,\infty)$, and let $\hat T,T_0:\Omega\to\Omega$ be measurable. Suppose $a_0^2=w_0\in[w_-,w_+]$ and set $\hat\gamma_0=\hat a^2\hat\eta$. Define
\begin{align*}
    \hat\lambda(x)
    \coloneqq
    \clip_{[w_-,w_+]}\bigl(\hat a(x)^2\bigr)
    \exp\left(\frac14\|x-\hat T(x)\|^2\right),
    \qquad
    \lambda_0(x)
    \coloneqq
    w_0(x)\exp\left(\frac14\|x-T_0(x)\|^2\right).
\end{align*}
Then there exists a constant $C>0$, depending only on $w_-,w_+$ and $\operatorname{diam}(\Omega)$, such that
\begin{align}
\label{eq:continuous_gh_growth_transfer}
    \int_\Omega |\hat\lambda(x)-\lambda_0(x)|^2 d\hat\eta(x)
    \le
    C\int_\Omega |\hat a(x)-a_0(x)|^2 d\hat\eta(x)
    + C\int_\Omega \|\hat T(x)-T_0(x)\|^2 d\hat\gamma_0(x).
\end{align}
\end{lemma}

\begin{proof}[Proof of Lemma \ref{lem:continuous_gh_growth_transfer}]
The pointwise argument in Lemma~\ref{lem:discrete_gh_growth_transfer} gives
\begin{align*}
    |\hat\lambda(x)-\lambda_0(x)|^2
    \le
    C|\hat a(x)-a_0(x)|^2+C\|\hat T(x)-T_0(x)\|^2.
\end{align*}
To integrate the last term with respect to $\hat\gamma_0$ instead of $\hat\eta$, we
prove the pointwise inequality
\begin{align}
\label{eq:gh_transfer_pointwise_helper}
    \|\hat T(x)-T_0(x)\|^2
    \le
    C'\hat a(x)^2\|\hat T(x)-T_0(x)\|^2+C'|\hat a(x)-a_0(x)|^2,
    \qquad x\in\Omega,
\end{align}
for $C'\coloneqq \max \left(\tfrac{4}{w_-}, \tfrac{4 \mathrm{diam}(\Omega)^2}{w_-}\right)$.
Fix $x\in\Omega$ and split into two cases.
\smallskip\\
\emph{Case 1: $\hat a(x)^2\ge w_-/4$.} Then
$\hat a(x)^2\|\hat T-T_0\|^2\ge (w_-/4)\|\hat T-T_0\|^2$, hence
$\|\hat T-T_0\|^2\le (4/w_-)\hat a(x)^2\|\hat T-T_0\|^2$.
\smallskip\\
\emph{Case 2: $\hat a(x)^2<w_-/4$.} Since $a_0^2(x)=w_0(x)\ge w_-$,
$a_0(x)\ge \sqrt{w_-}$, while $\hat a(x)<\sqrt{w_-}/2$, so
$|\hat a(x)-a_0(x)|\ge a_0(x)-\hat a(x)\ge \sqrt{w_-}/2$, i.e.\
$|\hat a(x)-a_0(x)|^2\ge w_-/4$. Combined with the trivial bound
$\|\hat T(x)-T_0(x)\|^2\le \mathrm{diam}(\Omega)^2$, we obtain
$\|\hat T-T_0\|^2\le (4 \mathrm{diam}(\Omega)^2/w_-)|\hat a-a_0|^2$.
\smallskip\\
In either case \eqref{eq:gh_transfer_pointwise_helper} holds.
Multiplying by $\hat\eta$ and integrating, and noting that
$\hat a(x)^2 d\hat\eta(x)=d\hat\gamma_0(x)$ by definition of $\hat\gamma_0$,
yields \eqref{eq:continuous_gh_growth_transfer}.
\end{proof}


\begin{lemma}[{Bound on $\Ep[D_\KL(\hat r\mid\hat r^\star)]$}]
\label{lem:uot_fitted_row_kl_expect}
Assume that Assumptions~\ref{assm:curvature}, \ref{assm:supp_global}, \ref{assm:positivity}, and \ref{assm:masses} hold.
Let $\hat\gamma$ be any optimizer of $\mathrm{UOT}(\hat\mu_n,\hat\nu_m)$, define
\begin{align*}
    \hat r_i \coloneqq \sum_{j=1}^m \hat\gamma_{ij},
    \qquad
    \hat\mu_i \coloneqq {\hat M_\mu/n},
    \qquad
    \hat r_i^\star \coloneqq e^{-\varphi_0(X_i)}\hat\mu_i,
\end{align*}
and write $\hat r=(\hat r_i)_{i=1}^n$ and $\hat r^\star=(\hat r_i^\star)_{i=1}^n$.
Then there exists a constant $C>0$, depending only on
\(
M_\mu,M_\nu,\|\zeta_0\|_\infty,\|\xi_0\|_\infty
\)
and $C_\Lambda$ from Proposition~\ref{prop:two_sample_stability}, such that for all sufficiently large $n,m$,
\begin{align}
\label{eq:uot_fitted_row_kl_expect}
    \Ep[{D_\KL}(\hat r\mid \hat r^\star)]
    \le
    C
    \Bigl(
        \Ep[M_\mu W_2^2(\bar\mu_n,\bar\mu)]
        + \Ep[M_\nu W_2^2(\bar\nu_m,\bar\nu)]
        + a_n + b_m
    \Bigr),
\end{align}
where $a_n,b_m$ are the rates from Assumption~\ref{assm:masses}.
\end{lemma}

\begin{proof}[Proof of Lemma \ref{lem:uot_fitted_row_kl_expect}]
By \eqref{eq:uot_hat_exact_empirical_excess} and the nonnegativity of the transport-slack term and of ${D_\KL}(\hat s\mid \hat s^\star)$, the inequality
\begin{align*}
    {D_\KL}(\hat r\mid \hat r^\star)
    \le
    \mathrm{UOT}(\hat\mu_n,\hat\nu_m)
    - \int \zeta_0 d\hat\mu_n
    - \int \xi_0 d\hat\nu_m
\end{align*}
holds.
By the stability bound \eqref{eq:estimated_mass_stability},
\begin{align*}
&\mathrm{UOT}(\hat\mu_n,\hat\nu_m)-\mathrm{UOT}(\mu,\nu) \\
&\qquad\le
\int \zeta_0 d(\tilde\mu_n-\mu)
+ \int \xi_0 d(\tilde\nu_m-\nu)
+ C_\Lambda\Bigl(
    M_\mu W_2^2(\bar\mu_n,\bar\mu)+M_\nu W_2^2(\bar\nu_m,\bar\nu)
\Bigr) \\
&\qquad\quad
+ \bigl( |\hat M_\mu-M_\mu|+|\hat M_\nu-M_\nu| \bigr)
+ e\sqrt{M_\mu M_\nu}\bigl(|\log\alpha_n|+|\log\beta_m|\bigr),
\end{align*}
using $M_\mu|\alpha_n-1|=|\hat M_\mu-M_\mu|$ and $M_\nu|\beta_m-1|=|\hat M_\nu-M_\nu|$.
Subtracting
\(
\int \zeta_0 d(\hat\mu_n-\mu)+\int \xi_0 d(\hat\nu_m-\nu)
\)
from both sides and using
\(
\mathrm{UOT}(\mu,\nu)=\int \zeta_0 d\mu+\int \xi_0 d\nu
\)
gives
\begin{align*}
    \mathrm{UOT}(\hat\mu_n,\hat\nu_m)
    - \int \zeta_0 d\hat\mu_n
    - \int \xi_0 d\hat\nu_m
    &\le
    C_\Lambda\Bigl(
        M_\mu W_2^2(\bar\mu_n,\bar\mu)+M_\nu W_2^2(\bar\nu_m,\bar\nu)
    \Bigr) \\
    &\quad +
    \bigl( |\hat M_\mu-M_\mu|+|\hat M_\nu-M_\nu| \bigr) \\
    &\quad +
    e\sqrt{M_\mu M_\nu}\bigl(|\log\alpha_n|+|\log\beta_m|\bigr) \\
    &\quad
    - \int \zeta_0 d(\hat\mu_n-\tilde\mu_n)
    - \int \xi_0 d(\hat\nu_m-\tilde\nu_m).
\end{align*}
The last two integrals are bounded in absolute value by
\begin{align*}
    \|\zeta_0\|_\infty |\hat M_\mu-M_\mu|
    +
    \|\xi_0\|_\infty |\hat M_\nu-M_\nu|.
\end{align*}
Taking expectations of the resulting inequality, the unbiased linear terms vanish:
$\Ep[\int\zeta_0 d(\tilde\mu_n-\mu)]=\Ep[\int\xi_0 d(\tilde\nu_m-\nu)]=0$.
By Assumption~\ref{assm:masses}, $\Ep[|\hat M_\mu-M_\mu|]\le a_n$ and the a.s.\ lower bound $\hat M_\mu\ge cM_\mu$ gives $|\log\alpha_n|\le|\alpha_n-1|/c$ a.s., so $M_\mu\Ep[|\log\alpha_n|]\le a_n/c$. Hence
\(
e\sqrt{M_\mu M_\nu} \Ep[|\log\alpha_n|]\le c^{-1}e\sqrt{M_\nu/M_\mu} a_n.
\)
The analogous bounds hold on the $\nu$-side with $b_m$.
Setting
\(
C\coloneqq C_\Lambda
+ \|\zeta_0\|_\infty+\|\xi_0\|_\infty+1
+ c^{-1}e\sqrt{M_\nu/M_\mu}+c^{-1}e\sqrt{M_\mu/M_\nu}
\)
and combining the displays proves \eqref{eq:uot_fitted_row_kl_expect}.
\end{proof}

\begin{lemma}[Bound on $\sum_i\hat\mu_i|\hat a_i-a_0(X_i)|^2$]
\label{lem:uot_complete_1nn_active_fitted}
We have
\begin{align}
\label{eq:uot_complete_1nn_active_fitted_empirical}
    \sum_{i=1}^n \hat\mu_i
    |\hat a_i-a_0(X_i)|^2
    \le
    {D_\KL}(\hat r\mid \hat r^\star).
\end{align}
Consequently, with $a_n,b_m$ as in Assumption~\ref{assm:masses},
\begin{align}
\label{eq:uot_complete_1nn_active_fitted_expectation}
    \Ep\left[
    \sum_{i=1}^n \hat\mu_i
    |\hat a_i-a_0(X_i)|^2
    \right]
    \le
    C
    \Bigl(
        \Ep[M_\mu W_2^2(\bar\mu_n,\bar\mu)]
        + \Ep[M_\nu W_2^2(\bar\nu_m,\bar\nu)]
        + a_n + b_m
    \Bigr).
\end{align}
\end{lemma}

\begin{proof}[Proof of Lemma \ref{lem:uot_complete_1nn_active_fitted}]
Since $\hat r_i=\hat\mu_i\hat a_i^2$ and $\hat r_i^\star=\hat\mu_i a_0(X_i)^2$, we have
\begin{align*}
    {D_\KL}(\hat r\mid \hat r^\star)
    =
    \sum_{i=1}^n\hat\mu_i
    \left[
        \hat a_i^2
        \log\left(\frac{\hat a_i^2}{a_0(X_i)^2}\right)
        -\hat a_i^2
        +a_0(X_i)^2
    \right].
\end{align*}
Lemma~\ref{lem:uot_kl_sqrt_lower_bound} gives \eqref{eq:uot_complete_1nn_active_fitted_empirical} termwise, and the expectation bound follows from Lemma~\ref{lem:uot_fitted_row_kl_expect}.
\end{proof}

\begin{lemma}[Bound on $\int|\hat\lambda^{1\mathrm{NN}}-\lambda_0|^2 d\mu$]
\label{lem:lambda_1nn_population_extension}
Assume Assumptions~\ref{assm:curvature}, \ref{assm:supp_global}, and \ref{assm:positivity}. Let $V_i,M_n,R_n$ be as in Lemma~\ref{lem:uot_voronoi_mass_radius}. On the mass-accuracy event $\mathcal A_{n,m}$ from \eqref{eq:uot_complete_1nn_estimated_mass_as}, there is a constant $C>0$ such that
\begin{align}
\label{eq:lambda_1nn_population_extension}
    \int_\Omega |\hat\lambda^{\mathrm{1NN}}(x)-\lambda_0(x)|^2d\mu(x)
    \le
    C nM_n\sum_{i=1}^n\hat\mu_i|\hat\lambda_i-\lambda_0(X_i)|^2
    + C R_n^2 .
\end{align}
\end{lemma}

\begin{proof}[Proof of Lemma~\ref{lem:lambda_1nn_population_extension}]
Assumption~\ref{assm:curvature}, compactness of $\Omega$, and the compactness of $\Omega$ imply that $\lambda_0=w_0\exp(\|x-T_0(x)\|^2/4)$ is Lipschitz; write $L_\lambda=\mathrm{Lip}(\lambda_0)$. For $x\in V_i$,
\[
|\hat\lambda_i-\lambda_0(x)|^2
\le
2|\hat\lambda_i-\lambda_0(X_i)|^2
+
2L_\lambda^2\|x-X_i\|^2 .
\]
Integrating over $V_i$ and using $\mu(V_i)=M_\mu\bar\mu(V_i)\le M_\mu M_n$ and, on $\mathcal A_{n,m}$, $\hat\mu_i=\hat M_\mu/n\ge M_\mu/(2n)$, gives
\[
\mu(V_i)\le 2nM_n\hat\mu_i .
\]
Summing over $i$ gives \eqref{eq:lambda_1nn_population_extension}.
\end{proof}

\begin{theorem}[Growth risk of $\hat\lambda^{1\mathrm{NN}}$]
\label{thm:uot_complete_1nn_lambda_fitted}
Assume that Assumptions~\ref{assm:curvature}, \ref{assm:supp_global}, \ref{assm:positivity}, and \ref{assm:masses} hold.
Let $\hat\lambda_i$ and $\hat\lambda^{\mathrm{1NN}}$ be defined by \eqref{eq:disc-gh-growth} and \eqref{eq:disc-1nn-map}. Then there exists a constant $C>0$ such that for all sufficiently large $n,m$,
\begin{align}
\label{eq:uot_complete_1nn_lambda_fitted_expectation}
\begin{aligned}
    \Ep\Bigl[
        \int_\Omega |\hat\lambda^{\mathrm{1NN}}(x)-\lambda_0(x)|^2  d\hat\mu_n(x)
    \Bigr]
    &\le
    C
    \Bigl(
        \Ep[M_\mu W_2^2(\bar\mu_n,\bar\mu)]
        + \Ep[M_\nu W_2^2(\bar\nu_m,\bar\nu)]
        + a_n + b_m
    \Bigr),
\end{aligned}
\end{align}
where $a_n,b_m$ are the rates from Assumption~\ref{assm:masses}.
\end{theorem}

\begin{proof}[Proof of Theorem \ref{thm:uot_complete_1nn_lambda_fitted}]
We first prove an in-sample bound and then pass to population loss through Lemma~\ref{lem:lambda_1nn_population_extension}.
Apply Lemma~\ref{lem:discrete_gh_growth_transfer} with
\begin{align*}
    u_i=\hat r_i/\hat\mu_i,
    \qquad
    w_i^0=w_0(X_i)=e^{-\varphi_0(X_i)},
    \qquad
    T_i=\hat T_i,
    \qquad
    T_i^0=T_0(X_i).
\end{align*}
This yields
\begin{align}
\label{eq:lambda_transfer_sum_plan}
    \sum_{i=1}^n\hat\mu_i |\hat\lambda_i-\lambda_0(X_i)|^2
    \le
    C\sum_{i=1}^n\hat\mu_i |\hat a_i-a_0(X_i)|^2
    + C\sum_{i=1}^n\hat\mu_i\|\hat T_i-T_0(X_i)\|^2.
\end{align}
The first term is controlled by Lemma~\ref{lem:uot_complete_1nn_active_fitted}. For the second term, since $w_0\ge w_-$,
\begin{align*}
    \sum_{i=1}^n\hat\mu_i\|\hat T_i-T_0(X_i)\|^2
    \le
    w_-^{-1}
    \sum_{i=1}^n\hat r_i^\star\|\hat T_i-T_0(X_i)\|^2.
\end{align*}
Lemma~\ref{lem:uot_weighted_pinsker_rows_hat} gives
\begin{align*}
    \sum_{i=1}^n\hat r_i^\star\|\hat T_i-T_0(X_i)\|^2
    \le
    3\hat\Delta_{nm}^{\mathrm{bar}}
    +2C_\Omega {D_\KL}(\hat r\mid \hat r^\star).
\end{align*}
Taking expectations and using Proposition~\ref{prop:uot_barycentric_and_1nn_estimated_mass} together with Lemma~\ref{lem:uot_fitted_row_kl_expect} proves \eqref{eq:uot_complete_1nn_lambda_fitted_expectation}.
\end{proof}

\begin{proof}[Proof of Theorem \ref{thm:main_plan_based_rates}]
The map bound in \eqref{eq:section4_plan_map_rate_abstract} follows from Theorem~\ref{thm:uot_complete_1nn_estimated_mass}. The growth bound in the same display is the bound for the correct Gaussian--Hellinger factor $\lambda_0$ and follows from Theorem~\ref{thm:uot_complete_1nn_lambda_fitted}. It remains only to substitute explicit rates for the weighted empirical measures.

By Assumptions~\ref{assm:supp_global} and \ref{assm:positivity}, the normalized measures $\bar\mu$ and $\bar\nu$ are supported on a bounded subset of $\mathbb{R}^d$ and have finite moments of every order. Choose any $q>4$. Theorem~1 of \cite{fournier2015rate}, applied with $p=2$, yields
\begin{align*}
    \Ep[W_2^2(\bar\mu_n,\bar\mu)]
    &\le C_\mu \mathfrak R_n^{\mathrm{emp}}(d),
    \qquad
    \Ep[W_2^2(\bar\nu_m,\bar\nu)]
    \le C_\nu \mathfrak R_m^{\mathrm{emp}}(d).
\end{align*}
Substituting these bounds together with the mass-estimation bounds from Assumption~\ref{assm:masses} into Theorems~\ref{thm:uot_complete_1nn_estimated_mass} and \ref{thm:uot_complete_1nn_lambda_fitted} yields \eqref{eq:section4_plan_map_rate_abstract}. The additional logarithm in the theorem statement covers the nearest-neighbor population extension term in the map bound.

Additionally, since \(d\gamma_0=w_0d\mu\) with
\(w_0=e^{-\varphi_0}\) and, by Assumption~\ref{assm:curvature},
\(w_0\ge w_->0\) on \(\Omega\), we have
\[
    \int_\Omega \|\hat T^{\mathrm{1NN}}-T_0\|^2 d\mu
    =
    \int_\Omega \|\hat T^{\mathrm{1NN}}-T_0\|^2 w_0^{-1} d\gamma_0
    \le
    w_-^{-1}
    \int_\Omega \|\hat T^{\mathrm{1NN}}-T_0\|^2 d\gamma_0 .
\]
Taking expectations and absorbing \(w_-^{-1}\) into the constant gives the desired
\(d\mu\)-risk bound.
\end{proof}


\section{Proof of Theorem \ref{thm:main_cube_kernel_rates}}\label{sec:proof_cube_kernel_rates}

Here, we additionally define
\begin{align}
    \widehat{\mathrm{UOT}}_{nm}^{\mathrm{ker}}
    &\coloneqq
    \mathrm{UOT}(\hat\mu_n^{\mathrm{ker}},\hat\nu_m^{\mathrm{ker}}).
\end{align}
After positive-part renormalization, let
\begin{align}
    \hat p_n^{\mathrm{ker}}(x)
    &\coloneqq
    \frac{(\tilde p_n^{\mathrm{ker}}(x))_+}
    {\int_{[0,1]^d}(\tilde p_n^{\mathrm{ker}}(u))_+ du},
    \qquad
    \hat q_m^{\mathrm{ker}}(y)
    \coloneqq
    \frac{(\tilde q_m^{\mathrm{ker}}(y))_+}
    {\int_{[0,1]^d}(\tilde q_m^{\mathrm{ker}}(v))_+ dv}.
\end{align}
We then define the oracle equal-mass measures
\begin{align}
    \tilde\mu_n^{\mathrm{ker}}(\cdot)
    &\coloneqq
    M_\mu \hat p_n^{\mathrm{ker}}(x) dx,
    \qquad
    \tilde\nu_m^{\mathrm{ker}}(\cdot)
    \coloneqq
    M_\nu \hat q_m^{\mathrm{ker}}(y) dy,
\end{align}
We also define a useful rate
\begin{align}
    \mathfrak L_n(\alpha)
    \coloneqq
    n^{{-(\alpha-1)/(2(\alpha-1)+d)}}.
\end{align}



\subsection{Convergence rates of boundary-adapted kernel density estimator}

For a zero-mean $r\in L^2([0,1]^d)$, we write
\begin{align}
\label{eq:neumann_hminus1_norm}
    \|r\|_{H^{-1}_N([0,1]^d)}^2
    \coloneqq
    \sum_{k\in\mathbb{N}_0^d\setminus\{0\}}
    \frac{|\langle r,\eta_k\rangle_{L^2([0,1]^d)}|^2}{\lambda_k}.
\end{align}
This is the natural negative-order Sobolev norm associated with the Neumann Laplacian eigenbasis.

\begin{lemma}[Neumann $H^{-1}$ controls $W_2$ on the cube]
\label{lem:cube_hminus1_to_w2}
Let $f$ and $g$ be probability densities on $[0,1]^d$ satisfying
\begin{align}
    0<\underline\beta \le f(x),g(x) \le \overline\beta < \infty,
    \qquad x\in[0,1]^d.
\end{align}
Then
\begin{align}
\label{eq:cube_hminus1_to_w2}
    W_2^2(f dx,g dx)
    \le
    \underline\beta^{-1}
    \|f-g\|_{H^{-1}_N([0,1]^d)}^2.
\end{align}
\end{lemma}

\begin{proof}[Proof of Lemma \ref{lem:cube_hminus1_to_w2}]
Set $r\coloneqq g-f$.
Because $\int_{[0,1]^d}r=0$, there exists a unique weak solution $u$ of the Neumann problem
\begin{align*}
    -\Delta u = r,
    \qquad
    \partial_\nu u = 0 \text{ on } \partial[0,1]^d,
    \qquad
    \int_{[0,1]^d}u = 0.
\end{align*}
Writing $u=\sum_{k\neq 0}u_k\eta_k$ gives $u_k = \langle r,\eta_k\rangle/\lambda_k$, and therefore
\begin{align*}
    \int_{[0,1]^d}\|\nabla u(x)\|^2 dx
    =
    \sum_{k\neq 0}\lambda_k |u_k|^2
    =
    \sum_{k\neq 0}\frac{|\langle r,\eta_k\rangle|^2}{\lambda_k}
    =
    \|r\|_{H^{-1}_N([0,1]^d)}^2.
\end{align*}
Now set $f_t=(1-t)f+t g$ for $t\in[0,1]$ and define
\begin{align*}
    v_t(x) \coloneqq \frac{\nabla u(x)}{f_t(x)}.
\end{align*}
Since $f_t\ge \underline\beta$, we have
\begin{align*}
    \partial_t f_t + \nabla \cdot(f_t v_t)
    =
    (g-f)+\nabla \cdot(\nabla u)
    =0,
\end{align*}
so $(f_t,v_t)$ is an admissible Benamou-Brenier path from $f dx$ to $g dx$.
Hence
\begin{align*}
    W_2^2(f dx,g dx)
    &\le
    \int_0^1\int_{[0,1]^d} f_t(x)\|v_t(x)\|^2 dx dt \\
    &=
    \int_0^1\int_{[0,1]^d}\frac{\|\nabla u(x)\|^2}{f_t(x)} dx dt \\
    &\le
    \underline\beta^{-1}\int_{[0,1]^d}\|\nabla u(x)\|^2 dx
    =
    \underline\beta^{-1}\|f-g\|_{H^{-1}_N([0,1]^d)}^2,
\end{align*}
which proves \eqref{eq:cube_hminus1_to_w2}.
\end{proof}


\begin{proposition}[Cube kernel density rates]
\label{prop:cube_kernel_density_rates}
Set $s\coloneqq \alpha-1>0$.
Assume that $\Omega=[0,1]^d$, that  Assumptions \ref{assm:cube_kernel_smoothness} hold, and that
\begin{align*}
    \beta_{\min} \le p(x),q(x) \le \beta_{\max}
    \quad \text{for all } x\in[0,1]^d.
\end{align*}
Let
\begin{align}
    L_n \asymp n^{1/(d+2s)},
    \qquad
    L_m \asymp m^{1/(d+2s)}.
\end{align}
Then there exists a constant $C>0$, depending only on $d,\alpha,M,\beta_{\min},\beta_{\max}$ and the cutoff $\tau$, such that
\begin{align}
\label{eq:cube_kernel_density_rate_source_w2}
    \Ep\bigl[W_2^2(\tilde\mu_n^{\mathrm{ker}},\mu)\bigr]
    &\le
    C M_\mu\mathfrak R_n^{\mathrm{ker}}(\alpha),
    \\
\label{eq:cube_kernel_density_rate_target_w2}
    \Ep\bigl[W_2^2(\tilde\nu_m^{\mathrm{ker}},\nu)\bigr]
    &\le
    C M_\nu\mathfrak R_m^{\mathrm{ker}}(\alpha),
\end{align}
and
\begin{align}
\label{eq:cube_kernel_density_rate_source_l1}
    \Ep\bigl[\|\hat p_n^{\mathrm{ker}}-p\|_{L^1([0,1]^d)}\bigr]
    &\le
    C\mathfrak L_n(\alpha),
    \\
\label{eq:cube_kernel_density_rate_target_l1}
    \Ep\bigl[\|\hat q_m^{\mathrm{ker}}-q\|_{L^1([0,1]^d)}\bigr]
    &\le
    C\mathfrak L_m(\alpha).
\end{align}
\end{proposition}




\begin{proof}[Proof of Proposition \ref{prop:cube_kernel_density_rates}]
We prove the source bounds; the target bounds are identical. For brevity write
\begin{align*}
    m_n(k)
    \coloneqq
    m_{L_n}^{\otimes}(k)
    =
    \prod_{r=1}^d
    \tau\left(\frac{\pi^2 k_r^2}{L_n^2}\right),
    \qquad k\in\mathbb N_0^d.
\end{align*}
Let
\begin{align*}
    \theta_k \coloneqq \int_{[0,1]^d} p(x)\eta_k(x) dx,
    \qquad
    \hat\theta_k \coloneqq {1/n}\sum_{i=1}^n \eta_k(X_i),
\end{align*}
so that $\theta_0=1$ and
\begin{equation}\label{eq:tildeker}
    \tilde p_n^{\mathrm{ker}}(x)
    =
    \sum_{k\in\mathbb N_0^d}
    m_n(k)\hat\theta_k\eta_k(x).
\end{equation}
Write
\begin{align*}
    p_{L_n}(x)
    \coloneqq
    \Ep\bigl[\tilde p_n^{\mathrm{ker}}(x)\bigr]
    =
    \sum_{k\in\mathbb N_0^d}
    m_n(k)\theta_k\eta_k(x).
\end{align*}
Equivalently, $p_{L_n}$ is the restriction to $[0,1]^d$ of the tensor-product smooth cosine cutoff on the doubled torus with coordinatewise multiplier $m_n$.
By Definition~\ref{def:neumann_holder_class}, the even reflection $\mathcal Ep$ belongs to $C^s(\mathbb T_2^d)$ with $\|\mathcal Ep\|_{C^s(\mathbb T_2^d)}\le M$.
Let $P_N$ denote the sharp Fourier projection onto frequencies $\|k\|_\infty\le N$ on $\mathbb T_2^d$.
Since $\tau\in C_c^\infty([0,\infty))$ with $\tau\equiv 1$ on a neighborhood of $0$, there exists $c>0$ such that $m_n(k)=1$ whenever $\|k\|_\infty\le cL_n$, and $|1-m_n(k)|\le \mathbf 1\{\|k\|_\infty>cL_n\}$ for all $k$.
Consequently, by Parseval's identity on $\mathbb T_2^d$,
\begin{align*}
    \|p_{L_n}-p\|_{L^2([0,1]^d)}
    \le
    \|\mathcal Ep-P_{cL_n}\mathcal Ep\|_{L^2(\mathbb T_2^d)}.
\end{align*}
Applying the multivariate Fourier truncation estimate \cite[Eq.~(5.8.4)]{canuto2006spectral} to $\mathcal Ep$ on $\mathbb T_2^d$ gives
\begin{align}
\label{eq:cube_kernel_bias_l2}
    \|p_{L_n}-p\|_{L^2([0,1]^d)}
    \le
    C L_n^{-s}.
\end{align}
Moreover, if the $k$th coefficient of $p_{L_n}-p$ is nonzero, then $m_n(k)\neq 1$, so for some coordinate $r$ one has $\pi^2k_r^2/L_n^2\ge 1$ and therefore $\lambda_k\ge \pi^2k_r^2\ge L_n^2$.
Hence, by \eqref{eq:neumann_hminus1_norm},
\begin{align}
\label{eq:cube_kernel_bias_hminus1}
    \|p_{L_n}-p\|_{H^{-1}_N([0,1]^d)}^2
    &=
    \sum_{k\neq 0}
    \frac{|1-m_n(k)|^2|\theta_k|^2}{\lambda_k}
    \\
    &\le
    C L_n^{-2}
    \sum_{k\neq 0}|1-m_n(k)|^2|\theta_k|^2
    \\
    &=
    C L_n^{-2}\|p_{L_n}-p\|_{L^2([0,1]^d)}^2
    \le
    C L_n^{-2(s+1)}
    =
    C L_n^{-2\alpha},
\end{align}
where the last inequality uses \eqref{eq:cube_kernel_bias_l2}.

For the stochastic term, orthonormality of $(\eta_k)$ and the upper bound on $p$ imply
\begin{align}
\label{eq:cube_kernel_variance_l2}
\begin{aligned}
    \Ep\bigl[\|\tilde p_n^{\mathrm{ker}}-p_{L_n}\|_{L^2([0,1]^d)}^2\bigr]
    &=
    \sum_{k\neq 0}
    m_n(k)^2
    \mathrm{Var}(\hat\theta_k)
    \\
    &\le
    \frac{C}{n}
    \#\Bigl\{k\in\mathbb N_0^d:m_n(k)\neq 0\Bigr\}
    \\
    &\le
    C\frac{L_n^d}{n}.
\end{aligned}
\end{align}
Likewise,
\begin{align}
\label{eq:cube_kernel_variance_hminus1}
\begin{aligned}
    \Ep\bigl[\|\tilde p_n^{\mathrm{ker}}-p_{L_n}\|_{H^{-1}_N([0,1]^d)}^2\bigr]
    &=
    \sum_{k\neq 0}
    \frac{m_n(k)^2}{\lambda_k}
    \mathrm{Var}(\hat\theta_k)
    \\
    &\le
    \frac{C}{n}
    \sum_{\substack{k\in\mathbb N_0^d\setminus\{0\}\\0\le k_r\lesssim L_n}}
    |k|^{-2}
    \\
    &\le
    C
    \begin{cases}
        n^{-1}, & d=1,
        \\
        (\log L_n)n^{-1}, & d=2,
        \\
        L_n^{d-2}n^{-1}, & d\ge 3.
    \end{cases}
\end{aligned}
\end{align}
Combining \eqref{eq:cube_kernel_bias_l2}, \eqref{eq:cube_kernel_bias_hminus1}, \eqref{eq:cube_kernel_variance_l2}, and \eqref{eq:cube_kernel_variance_hminus1}, and using Cauchy-Schwarz together with Jensen's inequality, we obtain
\begin{align}
    \Ep\bigl[\|\tilde p_n^{\mathrm{ker}}-p\|_{L^1([0,1]^d)}\bigr]
    &\le
    C\left(
        L_n^{-s}
        +
        \sqrt{\frac{L_n^d}{n}}
    \right),
    \\
\label{eq:cube_kernel_preliminary_hminus1}
    \Ep\bigl[\|\tilde p_n^{\mathrm{ker}}-p\|_{H^{-1}_N([0,1]^d)}^2\bigr]
    &\le
    C\left(
        L_n^{-2\alpha}
        +
        \begin{cases}
            n^{-1}, & d=1,
            \\
            (\log L_n)n^{-1}, & d=2,
            \\
            L_n^{d-2}n^{-1}, & d\ge 3
        \end{cases}
    \right).
\end{align}
We now derive the corresponding $L^1$ bound for the renormalized estimator $\hat p_n^{\mathrm{ker}}$.
Since $\tau(0)=1$ and $\eta_0\equiv 1$ on $[0,1]^d$, the spectral form \eqref{eq:tildeker} together with $\hat\theta_0={1/n}\sum_{i=1}^n\eta_0(X_i)=1$ gives
\begin{align*}
    \int_{[0,1]^d}\tilde p_n^{\mathrm{ker}}(x) dx
    = m_n(0)\hat\theta_0
    = 1.
\end{align*}
Set $u\coloneqq(\tilde p_n^{\mathrm{ker}})_+$ and $Z\coloneqq\int_{[0,1]^d}u$. Since $\int_{[0,1]^d}\tilde p_n^{\mathrm{ker}}=1$, we have
\begin{align*}
    Z-1 = \int_{[0,1]^d}(\tilde p_n^{\mathrm{ker}})_-.
\end{align*}
Moreover, on the set $\{\tilde p_n^{\mathrm{ker}}<0\}$ we have $p-\tilde p_n^{\mathrm{ker}}\ge -\tilde p_n^{\mathrm{ker}}=(\tilde p_n^{\mathrm{ker}})_-$, hence
\begin{align*}
    Z-1 = \int_{[0,1]^d}(\tilde p_n^{\mathrm{ker}})_- \le \|\tilde p_n^{\mathrm{ker}}-p\|_{L^1([0,1]^d)}.
\end{align*}
Also,
\begin{align*}
    \|u-\tilde p_n^{\mathrm{ker}}\|_{L^1([0,1]^d)}
    = \int_{[0,1]^d}(\tilde p_n^{\mathrm{ker}})_-
    = Z-1,
\end{align*}
so by the triangle inequality,
\begin{align*}
    \|u-p\|_{L^1([0,1]^d)}
    \le \|u-\tilde p_n^{\mathrm{ker}}\|_{L^1([0,1]^d)} + \|\tilde p_n^{\mathrm{ker}}-p\|_{L^1([0,1]^d)}
    \le 2\|\tilde p_n^{\mathrm{ker}}-p\|_{L^1([0,1]^d)}.
\end{align*}
Finally,
\begin{align*}
    \left\|\frac{u}{Z}-u\right\|_{L^1([0,1]^d)}
    = |1-Z| = Z-1
    \le \|\tilde p_n^{\mathrm{ker}}-p\|_{L^1([0,1]^d)},
\end{align*}
which together with the previous display implies
\begin{align}
\label{eq:cube_kernel_positive_part_l1}
    \|\hat p_n^{\mathrm{ker}}-p\|_{L^1([0,1]^d)}
    \le 3\|\tilde p_n^{\mathrm{ker}}-p\|_{L^1([0,1]^d)}.
\end{align}
Taking expectations and combining with the bound on $\Ep[\|\tilde p_n^{\mathrm{ker}}-p\|_{L^1([0,1]^d)}]$ above yields
\begin{align}
\label{eq:cube_kernel_l1_after_renormalization}
    \Ep\bigl[\|\hat p_n^{\mathrm{ker}}-p\|_{L^1([0,1]^d)}\bigr]
    \le
    C\left(
        L_n^{-s}
        +
        \sqrt{\frac{L_n^d}{n}}
    \right).
\end{align}

It remains to prove the Wasserstein bound.
Let $\underline p \coloneqq \inf_{x\in[0,1]^d} p(x)>0$.
Since $\tau\in C_c^\infty([0,\infty))$ with $\tau(0)=1$, the convolution kernel of the multiplier operator $T_{m_n}\coloneqq f\mapsto \sum_k m_n(k)\hat f(k)e^{i\pi k\cdot( \cdot )}$ on $\mathbb T_2^d$ is the periodization of $L_n^d \psi^{\otimes d}(L_n \cdot )$ for a fixed Schwartz function $\psi$ independent of $n$, and therefore has uniformly bounded $L^1(\mathbb T_2^d)$ norm; by Young's convolution inequality, $T_{m_n}$ is uniformly bounded on $L^\infty(\mathbb T_2^d)$. Moreover, $T_{m_n}$ acts as the identity on cosine polynomials of degree at most $cL_n$. Since $\mathcal Ep$ is continuous on $\mathbb T_2^d$, the Weierstrass approximation theorem produces a cosine polynomial $\phi_n$ of degree $\le cL_n$ with $\|\mathcal Ep-\phi_n\|_{L^\infty(\mathbb T_2^d)}\to 0$, and writing $\mathcal Ep-T_{m_n}\mathcal Ep=(\mathcal Ep-\phi_n)-T_{m_n}(\mathcal Ep-\phi_n)$ gives
\begin{align*}
    \|p_{L_n}-p\|_{L^\infty([0,1]^d)}
    =
    \|\mathcal Ep-T_{m_n}\mathcal Ep\|_{L^\infty(\mathbb T_2^d)}
    \to 0
    \qquad \text{as }n\to\infty.
\end{align*}
Hence $\inf p_{L_n}\ge \underline p/2$ for all sufficiently large $n$.
Moreover, the product form \eqref{eq:uot_cube_kernel} implies
\begin{align*}
    \sup_{x\in[0,1]^d} K_{L_n}(x,x)
    =
    \sup_{x\in[0,1]^d}
    \prod_{r=1}^d \kappa_{L_n}(x_r,x_r)
    \le
    C L_n^d.
\end{align*}
Define
\begin{align*}
    P f \coloneqq \int f(z)p(z)dz,
    \qquad
    P_n f \coloneqq {1/n}\sum_{i=1}^n f(X_i),
    \qquad
    f_x(z) \coloneqq K_{L_n}(x,z).
\end{align*}
Then
\begin{align*}
    \|\tilde p_n^{\mathrm{ker}}-p_{L_n}\|_{L^\infty([0,1]^d)}
    =
    \sup_{x\in[0,1]^d}|(P_n-P)f_x|.
\end{align*}
For every $x,z\in[0,1]^d$, Cauchy-Schwarz in the cosine expansion gives
\begin{align*}
    |f_x(z)|
    =
    |K_{L_n}(x,z)|
    \le
    K_{L_n}(x,x)^{1/2}K_{L_n}(z,z)^{1/2}
    \le
    C L_n^d,
\end{align*}
so the class $\mathcal F_n\coloneqq\{f_x:x\in[0,1]^d\}$ has envelope $U_n\le C L_n^d$.
Likewise, using $p\le \beta_{\max}$, orthonormality of $(\eta_k)$, and $0\le m_n(k)\le 1$,
\begin{align*}
    \mathrm{Var}(f_x(X_1))
    &\le
    \Ep[f_x(X_1)^2]
    \\
    &=
    \int K_{L_n}(x,z)^2 p(z)dz
    \\
    &\le
    \beta_{\max}\int K_{L_n}(x,z)^2 dz
    \\
    &=
    \beta_{\max}\sum_{k\in\mathbb N_0^d} m_n(k)^2\eta_k(x)^2
    \le
    \beta_{\max}K_{L_n}(x,x)
    \le
    C L_n^d.
\end{align*}
Hence $\sigma_n^2\coloneqq \sup_{x}\mathrm{Var}(f_x(X_1))\le C L_n^d$.
For the mean supremum, observe that $\mathcal F_n$ is contained in the finite-dimensional subspace $V_{L_n}\coloneqq\mathrm{span}\{\eta_k:m_n(k)\ne 0\}$ of $L^2([0,1]^d)$, of dimension at most $CL_n^d$. By a standard chaining bound for empirical processes indexed by a uniformly bounded subset of a finite-dimensional class with envelope $U_n$ and variance $\sigma_n^2$ \cite[Corollary~3.5.8]{gine2016mathematical},
\begin{align*}
    \Ep\Big[\sup_x |(P_n-P)f_x|\Big]
    \le
    C\sqrt{\frac{L_n^d\log n}{n}}.
\end{align*}
Applying Bernstein's inequality to $\mathcal F_n$ therefore yields, for every $t\ge 1$,
\begin{align}
    \Pr\left(
        \|\tilde p_n^{\mathrm{ker}}-p_{L_n}\|_{L^\infty([0,1]^d)}
        >
        C\left(
            \sqrt{\frac{L_n^d t}{n}}
            +
            \frac{L_n^d t}{n}
        \right)
    \right)
    \le
    2e^{-t}.
\end{align}

Taking $t=A\log n$ and using $L_n^d\log n/n\to 0$, we obtain
\begin{equation}\label{eq:bern}
    \Pr\left(
        \|\tilde p_n^{\mathrm{ker}}-p_{L_n}\|_{L^\infty([0,1]^d)}
        >
        \frac{\underline p}{4}
    \right)
    \le
    C_A n^{-A}
\end{equation}
for every fixed $A>0$ and all sufficiently large $n$.
Therefore the event
\begin{align*}
    \mathcal G_n
    \coloneqq
    \left\{
        \|\tilde p_n^{\mathrm{ker}}-p\|_{L^\infty([0,1]^d)}
        \le
        \frac{\underline p}{2}
    \right\}
\end{align*}
satisfies $\Pr(\mathcal G_n^c)\le C_A n^{-A}$.
On $\mathcal G_n$ we have $\tilde p_n^{\mathrm{ker}}\ge \underline p/2>0$, so the positive-part renormalization is inactive and
\begin{align*}
    \hat p_n^{\mathrm{ker}}
    =
    \tilde p_n^{\mathrm{ker}}.
\end{align*}
Since $p\ge \underline p$ and $\tilde p_n^{\mathrm{ker}}\ge \underline p/2$ on $\mathcal G_n$, Lemma~\ref{lem:cube_hminus1_to_w2} gives
\begin{align*}
    W_2^2\bigl(\hat p_n^{\mathrm{ker}}dx,p dx\bigr)
    =
    W_2^2\bigl(\tilde p_n^{\mathrm{ker}}dx,p dx\bigr)
    \le
    \frac{2}{\underline p}
    \|\tilde p_n^{\mathrm{ker}}-p\|_{H^{-1}_N([0,1]^d)}^2
    \qquad \text{on }\mathcal G_n.
\end{align*}
On the complement, both measures are supported in $[0,1]^d$, so
\begin{align*}
    W_2^2\bigl(\hat p_n^{\mathrm{ker}}dx,p dx\bigr)
    \le d.
\end{align*}
Taking expectations and using \eqref{eq:cube_kernel_preliminary_hminus1} therefore yields
\begin{align}
\label{eq:cube_kernel_w2_probability}
    \Ep\bigl[W_2^2(\hat p_n^{\mathrm{ker}}dx,p dx)\bigr]
    \le
    C\left(
        L_n^{-2\alpha}
        +
        \begin{cases}
            n^{-1}, & d=1,
            \\
            (\log L_n)n^{-1}, & d=2,
            \\
            L_n^{d-2}n^{-1}, & d\ge 3
        \end{cases}
    \right).
\end{align}
Now choose $L_n\asymp n^{1/(d+2s)}$.
Then \eqref{eq:cube_kernel_l1_after_renormalization} gives
\begin{align*}
    \Ep\bigl[\|\hat p_n^{\mathrm{ker}}-p\|_{L^1([0,1]^d)}\bigr]
    \le
    C n^{{-s/(2s+d)}}
    =
    C\mathfrak L_n(\alpha),
\end{align*}
while \eqref{eq:cube_kernel_w2_probability} becomes
\begin{align*}
    \Ep\bigl[W_2^2(\hat p_n^{\mathrm{ker}}dx,p dx)\bigr]
    \le
    C\mathfrak R_n^{\mathrm{ker}}(\alpha).
\end{align*}
Finally,
\begin{align*}
    W_2^2(\tilde\mu_n^{\mathrm{ker}},\mu)
    =
    M_\mu W_2^2(\hat p_n^{\mathrm{ker}}dx,p dx),
\end{align*}
which proves \eqref{eq:cube_kernel_density_rate_source_w2} and \eqref{eq:cube_kernel_density_rate_source_l1}.
The target bounds \eqref{eq:cube_kernel_density_rate_target_w2}-\eqref{eq:cube_kernel_density_rate_target_l1} are identical.
\end{proof}


\subsection{Proof of the upper bound}



One key ingredient to the proof is the following result from \cite{Cordero1999} that guarantees the Lipschitzness of the optimal transport map on the torus, given that the densities are bounded above and away from zero:

\begin{proposition}[{Regularity of OT on the doubled torus}]
\label{prop:torus_ot_doubled}
Let $P, Q \in \mathcal{P}_{\mathrm{ac}}([0,1]^d)$ admit densities $p, q$ satisfying $0 < \beta_{\min} \le p(x), q(x) \le \beta_{\max} < \infty$. Let $\mathcal{E}P, \mathcal{E}Q \in \mathcal{P}_{\mathrm{ac}}(\mathbb{T}_2^d)$ denote their respective even $2\mathbb{Z}^d$-periodic reflections on the doubled torus $\mathbb{T}_2^d = (\mathbb{R}/2\mathbb{Z})^d$, with densities $\mathcal{E}p, \mathcal{E}q$.

Then, there exists an optimal transport map $T_{\mathcal{E}} = \nabla \Psi$ from $\mathcal{E}P$ to $\mathcal{E}Q$, where the Brenier potential $\Psi: \mathbb{R}^d \to \mathbb{R}$ is a convex function satisfying the following properties:
\begin{enumerate}
    \item \emph{Periodicity}: The map $x \mapsto \frac{1}{2}\|x\|^2 - \Psi(x)$ is $2\mathbb{Z}^d$-periodic, and $T_{\mathcal{E}}(x+2k) = T_{\mathcal{E}}(x) + 2k$ for almost every $x \in \mathbb{R}^d$ and $k \in \mathbb{Z}^d$.
    \item \emph{Symmetry and Restriction}: Because the densities $\mathcal{E}p$ and $\mathcal{E}q$ are coordinatewise even, the map $T_{\mathcal{E}}$ leaves the hypercube $[0,1]^d$ invariant. The restriction of $T_{\mathcal{E}}$ to $[0,1]^d$ uniquely determines the optimal transport map $T_0$ from $P$ to $Q$.
    \item \emph{Lipschitz Continuity}: There exists a constant $\lambda \in (0,1)$ such that $\Psi$ is strongly convex and has a uniformly bounded Hessian:
    \begin{align}
        \lambda I_d \preceq \nabla^2 \Psi(x) \preceq \lambda^{-1} I_d, \quad \text{for all } x \in \mathbb{R}^d.
    \end{align}
    Consequently, the restricted transport map $T_0 = \nabla \Psi|_{[0,1]^d}$ and its optimal inverse $S_0 = \nabla \Psi^*|_{[0,1]^d}$ are Lipschitz on the hypercube $[0,1]^d$.
\end{enumerate}
\end{proposition}

We are now in a position to prove the main bounds of the kernel-based estimator.

\begin{corollary}[Kernel-based estimator's plug-in rate]
\label{cor:uot_cube_kernel_rate_transfer}
Assume that $\Omega=[0,1]^d$, $c(x,y)=\tfrac{1}{2}\|x-y\|^2$, and that Assumptions~\ref{assm:curvature}, \ref{assm:positivity}, and \ref{assm:masses} hold.
Let
\[
    \hat\mu_n^{\mathrm{ker}},\ \hat\nu_m^{\mathrm{ker}},\ \hat T_{nm}^{\mathrm{ker}},\ \hat\lambda_{nm}^{\mathrm{ker}},\ \widehat{\mathrm{UOT}}_{nm}^{\mathrm{ker}}.
\]
Let $a_n,b_m$ be the sequences from Assumption~\ref{assm:masses}.
Suppose that, for some deterministic sequences $r_n,r_m,\ell_n,\ell_m\ge 0$, where $r_n \ge 1/n$ and $r_m > 1/m$,
\begin{align}
    \Ep\bigl[W_2^2(\tilde\mu_n^{\mathrm{ker}},\mu)\bigr] \le r_n,
    \qquad
    \Ep\bigl[W_2^2(\tilde\nu_m^{\mathrm{ker}},\nu)\bigr] \le r_m,
\end{align}
\begin{align}
    \Ep\bigl[\|\hat p_n^{\mathrm{ker}}-p\|_{L^1([0,1]^d)}\bigr] \le \ell_n,
    \qquad
    \Ep\bigl[\|\hat q_m^{\mathrm{ker}}-q\|_{L^1([0,1]^d)}\bigr] \le \ell_m.
\end{align}
Then there exists a constant $C>0$, depending only on the constants in Assumption~\ref{assm:curvature}, such that
\begin{align}
\label{eq:uot_cube_kernel_rate_transfer_map}
    \Ep\left[
        \int_{[0,1]^d}
        \|\hat T_{nm}^{\mathrm{ker}}(x)-T_0(x)\|^2 d\mu(x)
    \right]
    &\le C(r_n+r_m+a_n+b_m),
\\
\label{eq:uot_cube_kernel_rate_transfer_lambda}
    \Ep\left[
        \int_{[0,1]^d}
        |\hat\lambda_{nm}^{\mathrm{ker}}(x)-\lambda_0(x)|^2 d\mu(x)
    \right]
    &\le C(r_n+r_m+a_n+b_m),
\end{align}
and
\begin{align}
\label{eq:uot_cube_kernel_rate_transfer_abs_value}
\begin{aligned}
    \Ep\left[
        \left|
            \widehat{\mathrm{UOT}}_{nm}^{\mathrm{ker}}
            -
            \mathrm{UOT}(\mu,\nu)
        \right|
    \right]
    &\le M_\mu\|\zeta_0\|_{L^\infty([0,1]^d)}\ell_n
    + M_\nu\|\xi_0\|_{L^\infty([0,1]^d)}\ell_m \\
    &\qquad + C(r_n+r_m+a_n+b_m).
\end{aligned}
\end{align}
\end{corollary}




\begin{proof}[Proof of Corollary \ref{cor:uot_cube_kernel_rate_transfer}]
Let $\hat\gamma_{nm}^{\mathrm{ker}}$ be an optimal plan for
$\widehat{\mathrm{UOT}}_{nm}^{\mathrm{ker}}$, with marginals
$\hat\gamma_{0,nm}^{\mathrm{ker}}$ and $\hat\gamma_{1,nm}^{\mathrm{ker}}$.

We first derive an algebraic identity that decomposes the empirical UOT cost into a transport-cost term and two $D_\KL$ penalties against oracle reference measures. By optimality of $\hat\gamma_{nm}^{\mathrm{ker}}$ for $\mathrm{UOT}(\hat\mu_n^{\mathrm{ker}},\hat\nu_m^{\mathrm{ker}})$,
\begin{align}
\label{eq:uot_cube_kernel_exact_excess_start}
    \widehat{\mathrm{UOT}}_{nm}^{\mathrm{ker}}
    =
    \int \tfrac{1}{2}\|x-y\|^2 d\hat\gamma_{nm}^{\mathrm{ker}}(x,y)
    + {D_\KL}\bigl(\hat\gamma_{0,nm}^{\mathrm{ker}}\mid \hat\mu_n^{\mathrm{ker}}\bigr)
    + {D_\KL}\bigl(\hat\gamma_{1,nm}^{\mathrm{ker}}\mid \hat\nu_m^{\mathrm{ker}}\bigr).
\end{align}
We rewrite the source $D_\KL$ term with the new reference measure $\gamma_0^{\mathrm{or}}\coloneqq e^{-\varphi_0}\hat\mu_n^{\mathrm{ker}}$. If $f=d\hat\gamma_{0,nm}^{\mathrm{ker}}/d\hat\mu_n^{\mathrm{ker}}$, then $d\hat\gamma_{0,nm}^{\mathrm{ker}}/d\gamma_0^{\mathrm{or}} = fe^{\varphi_0}$, and therefore
\begin{align*}
    {D_\KL}\bigl(\hat\gamma_{0,nm}^{\mathrm{ker}}\mid \gamma_0^{\mathrm{or}}\bigr)
    &= \int\Bigl[f\log\bigl(fe^{\varphi_0}\bigr)-f+e^{-\varphi_0}\Bigr]d\hat\mu_n^{\mathrm{ker}} \\
    &= {D_\KL}\bigl(\hat\gamma_{0,nm}^{\mathrm{ker}}\mid \hat\mu_n^{\mathrm{ker}}\bigr)
    + \int \varphi_0 d\hat\gamma_{0,nm}^{\mathrm{ker}}
    - \int \zeta_0 d\hat\mu_n^{\mathrm{ker}},
\end{align*}
where we used $\zeta_0=-(e^{-\varphi_0}-1)$. Hence
\begin{align}
\label{eq:uot_cube_kernel_exact_excess_source}
    {D_\KL}\bigl(\hat\gamma_{0,nm}^{\mathrm{ker}}\mid \hat\mu_n^{\mathrm{ker}}\bigr)
    =
    {D_\KL}\bigl(\hat\gamma_{0,nm}^{\mathrm{ker}}\mid \gamma_0^{\mathrm{or}}\bigr)
    - \int \varphi_0 d\hat\gamma_{0,nm}^{\mathrm{ker}}
    + \int \zeta_0 d\hat\mu_n^{\mathrm{ker}}.
\end{align}
Exactly the same computation with $\gamma_1^{\mathrm{or}}\coloneqq e^{-\psi_0}\hat\nu_m^{\mathrm{ker}}$ yields
\begin{align}
\label{eq:uot_cube_kernel_exact_excess_target}
    {D_\KL}\bigl(\hat\gamma_{1,nm}^{\mathrm{ker}}\mid \hat\nu_m^{\mathrm{ker}}\bigr)
    =
    {D_\KL}\bigl(\hat\gamma_{1,nm}^{\mathrm{ker}}\mid \gamma_1^{\mathrm{or}}\bigr)
    - \int \psi_0 d\hat\gamma_{1,nm}^{\mathrm{ker}}
    + \int \xi_0 d\hat\nu_m^{\mathrm{ker}}.
\end{align}
Substituting \eqref{eq:uot_cube_kernel_exact_excess_source} and \eqref{eq:uot_cube_kernel_exact_excess_target} into \eqref{eq:uot_cube_kernel_exact_excess_start} and using
\begin{align*}
    \int \varphi_0 d\hat\gamma_{0,nm}^{\mathrm{ker}} + \int \psi_0 d\hat\gamma_{1,nm}^{\mathrm{ker}}
    =
    \int \bigl(\varphi_0(x)+\psi_0(y)\bigr) d\hat\gamma_{nm}^{\mathrm{ker}}(x,y),
\end{align*}
we obtain
\begin{align}
\label{eq:uot_cube_kernel_exact_excess}
\begin{aligned}
&\widehat{\mathrm{UOT}}_{nm}^{\mathrm{ker}}
- \int \zeta_0 d\hat\mu_n^{\mathrm{ker}}
- \int \xi_0 d\hat\nu_m^{\mathrm{ker}} \\
&\qquad=
\int_{[0,1]^d\times[0,1]^d}
\Bigl(
\tfrac{1}{2}\|x-y\|^2-\varphi_0(x)-\psi_0(y)
\Bigr)
 d\hat\gamma_{nm}^{\mathrm{ker}}(x,y) \\
&\qquad\quad
+ {D_\KL}\bigl(\hat\gamma_{0,nm}^{\mathrm{ker}}\mid e^{-\varphi_0}\hat\mu_n^{\mathrm{ker}}\bigr)
+ {D_\KL}\bigl(\hat\gamma_{1,nm}^{\mathrm{ker}}\mid e^{-\psi_0}\hat\nu_m^{\mathrm{ker}}\bigr).
\end{aligned}
\end{align}
By the stability bound \eqref{eq:estimated_mass_stability_linearized}, defining the value-functional excess
\begin{align}
\label{eq:uot_cube_kernel_value_gap}
    V_{nm}^{\mathrm{ker}}
    \coloneqq
    \widehat{\mathrm{UOT}}_{nm}^{\mathrm{ker}}
    -\mathrm{UOT}(\mu,\nu)
    -\int\zeta_0 d(\hat\mu_n^{\mathrm{ker}}-\mu)
    -\int\xi_0 d(\hat\nu_m^{\mathrm{ker}}-\nu),
\end{align}
its expectation satisfies
\begin{equation}
\label{eq:uot_cube_kernel_value_expansion_expect}
\Ep\bigl[V_{nm}^{\mathrm{ker}}\bigr]
\le
C(r_n+r_m+a_n+b_m).
\end{equation}

Applying Lemma~\ref{lem:uot_gap_sufficient} to the integrand in \eqref{eq:uot_cube_kernel_exact_excess}, we obtain the empirical active-marginal bound:
\begin{align}
\label{eq:uot_cube_kernel_active_marginal_bound}
    \frac{\kappa}{2}
    \int_{[0,1]^d}
    \|\hat T_{nm}^{\mathrm{ker}}(x)-T_0(x)\|^2 d\hat\gamma_{0,nm}^{\mathrm{ker}}(x)
    \le V_{nm}^{\mathrm{ker}}.
\end{align}
Directly calculating ${D_\KL}\bigl(\hat\gamma_{0,nm}^{\mathrm{ker}}\mid e^{-\varphi_0}\hat\mu_n^{\mathrm{ker}}\bigr)$ and using \eqref{eq:uot_cube_kernel_exact_excess} yield the following bound for the active-source factors:
\begin{align}
\label{eq:uot_cube_kernel_active_factor_bound}
    \int_{[0,1]^d}
    |\hat a_{nm}^{\mathrm{ker}}(x)-a_0(x)|^2 d\hat\mu_n^{\mathrm{ker}}(x)
    \le
    {D_\KL}\bigl(\hat\gamma_{0,nm}^{\mathrm{ker}}\mid e^{-\varphi_0}\hat\mu_n^{\mathrm{ker}}\bigr)
    \le V_{nm}^{\mathrm{ker}}.
\end{align}

Since $\varphi_0$ is continuous, it is uniformly bounded on $[0,1]^d$, which guarantees the existence of a constant $w_- > 0$ such that $
    a_0(x)^2 = e^{-\varphi_0(x)} \ge w_-$, for all $x \in [0,1]^d$. Also, since the images of both $\hat T_{nm}^{\mathrm{ker}}$ and $T_0$ are contained in $[0,1]^d$, we have
$
    \|\hat T_{nm}^{\mathrm{ker}}(x) - T_0(x)\|^2 \le \mathrm{diam}([0,1]^d)^2 = d.
$. Consequently,
\begin{align}
\label{eq:helper_derivation_step1}
    \|\hat T_{nm}^{\mathrm{ker}}(x) - T_0(x)\|^2 
    &\le \frac{1}{w_-} a_0(x)^2 \|\hat T_{nm}^{\mathrm{ker}}(x) - T_0(x)\|^2. \\
    &\le \frac{2}{w_-} \hat a_{nm}^{\mathrm{ker}}(x)^2 \|\hat T_{nm}^{\mathrm{ker}}(x) - T_0(x)\|^2 \\
    &\quad + \frac{2}{w_-} | \hat a_{nm}^{\mathrm{ker}}(x) - a_0(x) |^2 \|\hat T_{nm}^{\mathrm{ker}}(x) - T_0(x)\|^2 \\
    &\le
    C_{\mathrm{tr}} \hat a_{nm}^{\mathrm{ker}}(x)^2 \|\hat T_{nm}^{\mathrm{ker}}(x)-T_0(x)\|^2
    +
    C_{\mathrm{tr}} |\hat a_{nm}^{\mathrm{ker}}(x)-a_0(x)|^2. 
\end{align}
For some constant $C_{\mathrm{tr}}>0$. Integrating this bound with respect to $d\hat\mu_n^{\mathrm{ker}}(x)$, and noting that $d\hat\gamma_{0,nm}^{\mathrm{ker}}(x) = \hat a_{nm}^{\mathrm{ker}}(x)^2 d\hat\mu_n^{\mathrm{ker}}(x)$ by definition, we transfer the map error from the active marginal to the empirical measure:
\begin{align}
\label{eq:uot_cube_kernel_T_to_mu_n}
\begin{aligned}
    \int_{[0,1]^d}
    \|\hat T_{nm}^{\mathrm{ker}}(x)-T_0(x)\|^2 d\hat\mu_n^{\mathrm{ker}}(x)
    &\le
    C_{\mathrm{tr}} \int_{[0,1]^d} \|\hat T_{nm}^{\mathrm{ker}}(x)-T_0(x)\|^2 d\hat\gamma_{0,nm}^{\mathrm{ker}}(x) \\
    &\quad + C_{\mathrm{tr}} \int_{[0,1]^d} |\hat a_{nm}^{\mathrm{ker}}(x)-a_0(x)|^2 d\hat\mu_n^{\mathrm{ker}}(x) \\
    &\le C_{\mathrm{tr}} \bigl( 2/\kappa + 1 \bigr) V_{nm}^{\mathrm{ker}}.
\end{aligned}
\end{align}

Using the notations from Proposition~\ref{prop:cube_kernel_density_rates}, we define the high-probability event $\mathcal{G}_{nm}$ by
\begin{align}
    \mathcal{G}_{nm}
    \coloneqq
    \left\{
        \inf_{x \in [0,1]^d} \tilde p_n^{\mathrm{ker}}(x) \ge \frac{\underline p}{2}
        \quad \text{and} \quad
        \inf_{y \in [0,1]^d} \tilde q_m^{\mathrm{ker}}(y) \ge \frac{\underline q}{2}
    \right\},
\end{align}
where $\tilde p_n^{\mathrm{ker}}$ is defined in \eqref{eq:tildeker}, $\underline p = \inf_{x\in[0,1]^d} p(x)$, and $\tilde q_n^{\mathrm{ker}}$ and $\underline q$ are defined analogously. Conditional on $\mathcal{G}_{nm}$, we have $\hat p_n^{\mathrm{ker}} = \tilde p_n^{\mathrm{ker}}$ and $\hat q_m^{\mathrm{ker}} = \tilde q_m^{\mathrm{ker}}$, and the densities are uniformly bounded away from zero. Because the estimators are constructed using the Neumann-compatible cosine kernel \eqref{eq:uot_cube_kernel}, $\hat p_n^{\mathrm{ker}}$ and $\hat q_m^{\mathrm{ker}}$ coincide exactly with the restrictions of their even-reflected smooth extensions $\mathcal{E}\hat p_n^{\mathrm{ker}}$ and $\mathcal{E}\hat q_m^{\mathrm{ker}}$ on the doubled torus $\mathbb{T}_2^d$.

Consequently, the empirical active marginals $\hat\gamma_{0,nm}^{\mathrm{ker}} = e^{-\hat\varphi_{nm}^{\mathrm{ker}}} \hat\mu_n^{\mathrm{ker}}$ and $ \hat\gamma_{1,nm}^{\mathrm{ker}} = e^{-\hat\psi_{nm}^{\mathrm{ker}}} \hat\nu_m^{\mathrm{ker}}$ are strictly bounded away from zero and infinity on $\mathcal{G}_{nm}$. We may therefore apply the torus regularity result of Proposition~\ref{prop:torus_ot_doubled}. Let $\hat\Psi_{nm}^{\mathrm{ker}}$ be the optimal Brenier potential for the extended balanced transport problem on $\mathbb{T}_2^d$. By Proposition~\ref{prop:torus_ot_doubled}, there exists a uniform constant $\hat\lambda \in (0,1)$ such that the Hessian is uniformly bounded:
\begin{align}
    \hat\lambda I_d \preceq \nabla^2 \hat\Psi_{nm}^{\mathrm{ker}}(x) \preceq \hat\lambda^{-1} I_d, \quad \text{for all } x \in \mathbb{T}_2^d.
\end{align}
Because $\mathcal{E}\hat p_n^{\mathrm{ker}}$ and $\mathcal{E}\hat q_m^{\mathrm{ker}}$ are coordinatewise even, the symmetry guarantees that the transport map $\nabla\hat\Psi_{nm}^{\mathrm{ker}}$ leaves the sub-domain $[0,1]^d$ invariant. By uniqueness, its restriction to the hypercube is precisely our empirical active transport map:
\begin{align}
    \hat T_{nm}^{\mathrm{ker}}(x) = \nabla\hat\Psi_{nm}^{\mathrm{ker}}(x), \quad \text{for } x \in [0,1]^d.
\end{align}
The uniform Hessian bound directly implies that $\hat T_{nm}^{\mathrm{ker}}$ is Lipschitz on $[0,1]^d$ with constant $L_T \coloneqq \hat\lambda^{-1}$. By an identical application of Proposition~\ref{prop:torus_ot_doubled} to the population densities $p, q \in \mathcal{C}_N^{\alpha-1}([0,1]^d; M)$, the population map $T_0$ is also Lipschitz on $[0,1]^d$ with constant $L_0 \coloneqq \lambda^{-1}$.

Let $\hat\gamma_n^{\mathrm{ker}}$ be the optimal transport plan between $\mu$ and $\hat\mu_n^{\mathrm{ker}}$ for the squared Euclidean cost. Squaring the triangle inequality and integrating over $\hat\gamma_n^{\mathrm{ker}}(x,x')$, we obtain on $\mathcal{G}_{nm}$:
\begin{align}
\label{eq:uot_cube_w2_coupling}
\begin{aligned}
    \int_{[0,1]^d} \|\hat T_{nm}^{\mathrm{ker}}(x)-T_0(x)\|^2 d\mu(x) 
    &= \int_{[0,1]^d \times [0,1]^d} \|\hat T_{nm}^{\mathrm{ker}}(x)-T_0(x)\|^2 d\hat\gamma_n^{\mathrm{ker}}(x,x') \\
    &\le 3 \int \|\hat T_{nm}^{\mathrm{ker}}(x) - \hat T_{nm}^{\mathrm{ker}}(x')\|^2 d\hat\gamma_n^{\mathrm{ker}}(x,x') \\
    &\quad + 3 \int \|\hat T_{nm}^{\mathrm{ker}}(x') - T_0(x')\|^2 d\hat\mu_n^{\mathrm{ker}}(x') \\
    &\quad + 3 \int \|T_0(x') - T_0(x)\|^2 d\hat\gamma_n^{\mathrm{ker}}(x,x') \\
    &\le 3 L_T^2 W_2^2(\mu, \hat\mu_n^{\mathrm{ker}}) 
    + 3 C_{\mathrm{tr}} \bigl( 2/\kappa + 1 \bigr) V_{nm}^{\mathrm{ker}} 
    + 3 L_0^2 W_2^2(\mu, \hat\mu_n^{\mathrm{ker}}).
\end{aligned}
\end{align}
Now we focus on $\mathcal{G}_{nm}$. Since both the empirical map $\hat T_{nm}^{\mathrm{ker}}$ and the population map $T_0$ take values in the unit hypercube $[0,1]^d$, the squared map error is trivially bounded by the squared Euclidean diameter of the domain, $d$. Hence, it follows from the tail bound \eqref{eq:bern} and the union bound that:
\begin{align}
    \Ep\left[ \int_{[0,1]^d} \|\hat T_{nm}^{\mathrm{ker}}(x)-T_0(x)\|^2 d\mu(x) \mathbf{1}_{\mathcal{G}_{nm}^c} \right] &\le d\Pr(\mathcal{G}^c_{nm}) \\ 
    &\le d \cdot C_A(n^{-A} + m^{-A}).
\end{align}
By choosing $A \ge 2$, the expectation decays as $O(n^{-2} + m^{-2})$, which is absorbed by the density estimation rates in $W^2_2$, namely $O(r_n + r_m)$. 

Taking expectations of \eqref{eq:uot_cube_w2_coupling} and substituting \eqref{eq:uot_cube_kernel_value_expansion_expect} yields~\eqref{eq:uot_cube_kernel_rate_transfer_map}.

For the growth factor, Lemma~\ref{lem:continuous_gh_growth_transfer} provides the bound under the empirical measure:
\begin{align}
\label{eq:uot_cube_kernel_growth_empirical}
    \int_{[0,1]^d}
    |\hat\lambda_{nm}^{\mathrm{ker}}(x')-\lambda_0(x')|^2 d\hat\mu_n^{\mathrm{ker}}(x')
    &\le
    C {D_\KL}\bigl(\hat\gamma_{0,nm}^{\mathrm{ker}}\mid e^{-\varphi_0}\hat\mu_n^{\mathrm{ker}}\bigr) \\
    &\quad + C \int_{[0,1]^d}
    \|\hat T_{nm}^{\mathrm{ker}}(x')-T_0(x')\|^2 d\hat\gamma_{0,nm}^{\mathrm{ker}}(x') \nonumber \\
    &\le C V_{nm}^{\mathrm{ker}}.
\end{align}
Since $\hat T_{nm}^{\mathrm{ker}}$ and $\hat\varphi_{nm}^{\mathrm{ker}}$ have bounded derivatives on $\mathcal{G}_{nm}$, the clipped growth factor $\hat\lambda_{nm}^{\mathrm{ker}}$ is also globally Lipschitz with some constant $L_\lambda$. Reapplying the same Wasserstein coupling $\hat\gamma_n^{\mathrm{ker}}(x,x')$:
\begin{align}
\begin{aligned}
    \int_{[0,1]^d} |\hat\lambda_{nm}^{\mathrm{ker}}(x)-\lambda_0(x)|^2 d\mu(x) 
    &\le 3 L_\lambda^2 W_2^2(\mu, \hat\mu_n^{\mathrm{ker}}) 
    + 3 \int_{[0,1]^d} |\hat\lambda_{nm}^{\mathrm{ker}}(x')-\lambda_0(x')|^2 d\hat\mu_n^{\mathrm{ker}}(x') \\
    &\quad + 3 L_{\lambda_0}^2 W_2^2(\mu, \hat\mu_n^{\mathrm{ker}}).
\end{aligned}
\end{align}
Taking expectations and applying \eqref{eq:uot_cube_kernel_value_expansion_expect} and \eqref{eq:uot_cube_kernel_growth_empirical} proves~\eqref{eq:uot_cube_kernel_rate_transfer_lambda}.

Finally, taking absolute values in the linearized estimated-mass bound yields
\begin{align*}
    \left|
        \widehat{\mathrm{UOT}}_{nm}^{\mathrm{ker}}
        -
        \mathrm{UOT}(\mu,\nu)
    \right|
    &\le
    \left|
        \int \zeta_0 d(\tilde\mu_n^{\mathrm{ker}}-\mu)
    \right|
    +
    \left|
        \int \xi_0 d(\tilde\nu_m^{\mathrm{ker}}-\nu)
    \right| \\
    &\qquad +
    C_\Lambda\Bigl(
    W_2^2(\tilde\mu_n^{\mathrm{ker}},\mu)
    +
    W_2^2(\tilde\nu_m^{\mathrm{ker}},\nu)
    \Bigr) \\
    &\qquad +
    C\Bigl(
    |\hat M_\mu-M_\mu|+|\hat M_\nu-M_\nu|
    \Bigr).
\end{align*}
Because
\begin{align*}
    \left|\int \zeta_0 d(\tilde\mu_n^{\mathrm{ker}}-\mu)\right|
    &\le
    M_\mu\|\zeta_0\|_{L^\infty([0,1]^d)}
    \|\hat p_n^{\mathrm{ker}}-p\|_{L^1([0,1]^d)},
    \\
    \left|\int \xi_0 d(\tilde\nu_m^{\mathrm{ker}}-\nu)\right|
    &\le
    M_\nu\|\xi_0\|_{L^\infty([0,1]^d)}
    \|\hat q_m^{\mathrm{ker}}-q\|_{L^1([0,1]^d)},
\end{align*}
taking expectations proves \eqref{eq:uot_cube_kernel_rate_transfer_abs_value}.
\end{proof}



\begin{proof}[Proof of Theorem \ref{thm:main_cube_kernel_rates}]
Apply Corollary~\ref{cor:uot_cube_kernel_rate_transfer} with
\begin{align*}
    r_n = C M_\mu\mathfrak R_n^{\mathrm{ker}}(\alpha),
    \qquad
    r_m = C M_\nu\mathfrak R_m^{\mathrm{ker}}(\alpha),
\end{align*}
and
\begin{align*}
    \ell_n = C\mathfrak L_n(\alpha),
    \qquad
    \ell_m = C\mathfrak L_m(\alpha),
\end{align*}
which are available from Proposition~\ref{prop:cube_kernel_density_rates}. Assumption~\ref{assm:masses} supplies the bounds $\Ep[|\hat M_\mu-M_\mu|]\le a_n$ and $\Ep[|\hat M_\nu-M_\nu|]\le b_m$, so the corollary yields the two bounds collected in \eqref{eq:section4_cube_kernel_map_rate}. 
%The final sentence follows because, when $n\asymp m\asymp N$, the leading nonparametric term is $\mathfrak R_N^{\mathrm{ker}}(\alpha)$.
\end{proof}



\section{Boundary-corrected wavelet estimator on the hypercube}\label{sec:wavelet}

\subsection{Definition}

In this section, we replace the periodic wavelet construction by a boundary-corrected wavelet construction on the unit hypercube $\Omega=[0,1]^d$. This allows the smooth plugin estimator to operate directly on non-periodic data while keeping the same UOT reduction as before. 


\begin{definition}[Boundary-corrected wavelet basis]
\label{def:bc_wavelet_basis}
Fix an integer $N\ge 2$ and a base scale $j_0\ge \lceil \log_2 N\rceil$. Let $\zeta_0^{\mathrm{db}},\xi_0^{\mathrm{db}}$ denote the compactly supported $N$-regular Daubechies scaling and wavelet functions on $\mathbb R$. We follow the standard Cohen-Daubechies-Vial boundary correction.

For each level $j\ge j_0$, define the \emph{interior} one-dimensional translates
\begin{align*}
    \zeta_{j,k}^{\mathrm{int}}(t) \coloneqq 2^{j/2}\zeta_0^{\mathrm{db}}(2^jt-k),
    \qquad
    \xi_{j,k}^{\mathrm{int}}(t) \coloneqq 2^{j/2}\xi_0^{\mathrm{db}}(2^jt-k),
\end{align*}
for $N\le k\le 2^j-N-1$. The boundary-corrected construction replaces the first and last $N$ translates at each level by \emph{edge functions}
\begin{align*}
    \zeta_{j,k}^{L},\ \xi_{j,k}^{L},
    \qquad
    \zeta_{j,k}^{R},\ \xi_{j,k}^{R},
    \qquad k=0,\ldots,N-1,
\end{align*}
satisfying
\begin{align*}
    \operatorname{supp}(\zeta_{j,k}^{L}),\operatorname{supp}(\xi_{j,k}^{L})
    &\subset [0,(2N-1)2^{-j}],\\
    \operatorname{supp}(\zeta_{j,k}^{R}),\operatorname{supp}(\xi_{j,k}^{R})
    &\subset [1-(2N-1)2^{-j},1].
\end{align*}
We then define the one-dimensional boundary-corrected scaling and wavelet families by
\begin{align}
    \zeta_{j,k}^{\mathrm{bc},1}(t)
    &\coloneqq
    \begin{cases}
        \zeta_{j,k}^{L}(t), & 0\le k\le N-1,\\
        \zeta_{j,k}^{\mathrm{int}}(t), & N\le k\le 2^j-N-1,\\
        \zeta_{j,2^j-1-k}^{R}(t), & 2^j-N\le k\le 2^j-1,
    \end{cases}\\
    \xi_{j,k}^{\mathrm{bc},1}(t)
    &\coloneqq
    \begin{cases}
        \xi_{j,k}^{L}(t), & 0\le k\le N-1,\\
        \xi_{j,k}^{\mathrm{int}}(t), & N\le k\le 2^j-N-1,\\
        \xi_{j,2^j-1-k}^{R}(t), & 2^j-N\le k\le 2^j-1.
    \end{cases}
\end{align}
Thus the correction is local: only the $2N$ boundary indices are modified, while all interior indices coincide with the usual Daubechies translates.

Now write $K(j)\coloneqq \{0,\ldots,2^j-1\}^d$. For $k=(k_1,\ldots,k_d)\in K(j_0)$, $x=(x_1,\ldots,x_d)\in \Omega$, and $\ell\in\{0,1\}^d\setminus\{0\}$, define
\begin{align*}
    \zeta_{j_0,k}^{\mathrm{bc}}(x)
    &\coloneqq
    \prod_{r=1}^d \zeta_{j_0,k_r}^{\mathrm{bc},1}(x_r),\\
    \xi_{j,k,\ell}^{\mathrm{bc}}(x)
    &\coloneqq
    \prod_{r:\ell_r=0}\zeta_{j,k_r}^{\mathrm{bc},1}(x_r)
    \prod_{r:\ell_r=1}\xi_{j,k_r}^{\mathrm{bc},1}(x_r),
    \qquad j\ge j_0.
\end{align*}
Set
\begin{align*}
    \Phi^{\mathrm{bc}}
    &\coloneqq
    \left\{\zeta_{j_0,k}^{\mathrm{bc}}:k\in K(j_0)\right\},\\
    \Psi_j^{\mathrm{bc}}
    &\coloneqq
    \left\{\xi_{j,k,\ell}^{\mathrm{bc}}:k\in K(j),\ \ell\in\{0,1\}^d\setminus\{0\}\right\},
    \qquad j\ge j_0.
\end{align*}
The resulting family
\begin{align*}
    \Psi^{\mathrm{bc}}
    \coloneqq
    \Phi^{\mathrm{bc}}\cup\bigcup_{j\ge j_0}\Psi_j^{\mathrm{bc}}
\end{align*}
is the boundary-corrected tensor-product wavelet basis on $\Omega=[0,1]^d$. By the Cohen-Daubechies-Vial construction, it is an orthonormal basis of $L^2(\Omega)$, and the scaling space $\mathrm{span}(\Phi^{\mathrm{bc}})$ contains all tensor-product polynomials of coordinatewise degree at most $N-1$; in particular, it contains the constants.
\end{definition}

\begin{remark}
At every scale $j$, the basis agrees with the ordinary Daubechies system on the interior cells $k=N,\ldots,2^j-N-1$. The only modification is at the first and last $N$ cells, where the translates that would otherwise cross the boundary are replaced by edge functions supported inside $[0,(2N-1)2^{-j}]$ and $[1-(2N-1)2^{-j},1]$. In particular, the cube estimator does not periodize the data or wrap information across opposite faces of $[0,1]^d$; the non-periodic behavior is handled locally at the boundary.
\end{remark}


\begin{definition}[Boundary-corrected wavelet plugin estimator]
\label{def:uot_wavelet_plugin}
Let
    $\Psi^{\mathrm{bc}} = \Phi^{\mathrm{bc}}\cup \bigcup_{j\ge j_0}\Psi_j^{\mathrm{bc}}$
be the boundary-corrected tensor-product wavelet basis from Definition~\ref{def:bc_wavelet_basis}. For $\zeta\in\Phi^{\mathrm{bc}}$ and $\xi\in\Psi_j^{\mathrm{bc}}$, define the empirical coefficients
$
    \hat\alpha_\zeta \coloneqq \tfrac{1}{n}\sum_{i=1}^n \zeta(X_i),
    \hat\beta_\zeta \coloneqq \tfrac{1}{m}\sum_{j=1}^m \zeta(Y_j),
$
and
$
    \hat\alpha_\xi \coloneqq \tfrac{1}{n}\sum_{i=1}^n \xi(X_i),
    \hat\beta_\xi \coloneqq \tfrac{1}{m}\sum_{j=1}^m \xi(Y_j).
$
Given resolution levels $J_n,J_m\ge j_0$, the preliminary boundary-corrected wavelet estimators of the normalized densities $p$ and $q$ are
\begin{align}
    \tilde p_n^{\mathrm{wav}}(x)
    &\coloneqq
    \sum_{\zeta\in\Phi^{\mathrm{bc}}}\hat\alpha_\zeta\zeta(x)
    + \sum_{j=j_0}^{J_n}\sum_{\xi\in\Psi_j^{\mathrm{bc}}}\hat\alpha_\xi\xi(x), \quad 
    \tilde q_m^{\mathrm{wav}}(y)
    \coloneqq
    \sum_{\zeta\in\Phi^{\mathrm{bc}}}\hat\beta_\zeta\zeta(y)
    + \sum_{j=j_0}^{J_m}\sum_{\xi\in\Psi_j^{\mathrm{bc}}}\hat\beta_\xi\xi(y).
\end{align}
Since these preliminary estimators may take negative values, we define the positive-part renormalizations
\begin{align}
\label{eq:uot_wavelet_density_estimator}
    \hat p_n^{\mathrm{wav}}(x)
    &\coloneqq
    \frac{(\tilde p_n^{\mathrm{wav}}(x))_+}{\int_{\Omega}(\tilde p_n^{\mathrm{wav}}(u))_+ du},
    \qquad
    \hat q_m^{\mathrm{wav}}(y)
    \coloneqq
    \frac{(\tilde q_m^{\mathrm{wav}}(y))_+}{\int_{\Omega}(\tilde q_m^{\mathrm{wav}}(v))_+ dv}.
\end{align}
We then set
    $\hat\mu_n^{\mathrm{wav}} \coloneqq M_\mu\hat p_n^{\mathrm{wav}}dx$ and
    $\hat\nu_m^{\mathrm{wav}} \coloneqq M_\nu\hat q_m^{\mathrm{wav}}dx$.
The \emph{boundary-corrected wavelet plugin UOT estimator} is defined as the fitted transport-growth pair
\begin{align}
    \bigl(\hat T_{nm}^{\mathrm{wav}},\hat\lambda_{nm}^{\mathrm{wav}}\bigr)
    \coloneqq
    \bigl(T_{\hat\mu_n^{\mathrm{wav}},\hat\nu_m^{\mathrm{wav}}},\lambda_{\hat\mu_n^{\mathrm{wav}},\hat\nu_m^{\mathrm{wav}}}\bigr),
\end{align}
whenever the fitted problem between $\hat\mu_n^{\mathrm{wav}}$ and $\hat\nu_m^{\mathrm{wav}}$ is Monge-type in the above sense.
\end{definition}

\begin{remark}
The positive-part renormalization in \eqref{eq:uot_wavelet_density_estimator} is included to guarantee that the fitted object is a probability density on $\Omega$. Theorem~\ref{thm:uot_plugin_density_rates} below shows that this truncation does not alter the $L^1$ rate. If the masses $M_\mu$ and $M_\nu$ are unknown, they may again be replaced by external estimators $\hat M_\mu$ and $\hat M_\nu$ without changing the definition of the fitted pair.
\end{remark}



\subsection{Theoretical Analysis}



For the boundary-corrected wavelet construction on the cube, the rate statement becomes completely explicit. Write
\begin{align}
    \mathfrak R_n^{\mathrm{wav}}(\alpha)
    \coloneqq
    \begin{cases}
        n^{-1}, & d=1,\\
        (\log n)^2 n^{-1}, & d=2,\\
        n^{{-2\alpha/(2(\alpha-1)+d)}}, & d\ge 3.
    \end{cases}
\end{align}

\begin{theorem}[Wavelet plugin rate]
\label{thm:main_wavelet_rates}
Assume that $\Omega=[0,1]^d$, that $c(x,y)=\tfrac{1}{2}\|x-y\|^2$, and let
\begin{align*}
    \hat\mu_n^{\mathrm{wav}},\qquad \hat\nu_m^{\mathrm{wav}},\qquad \hat T_{nm}^{\mathrm{wav}},\qquad \hat\lambda_{nm}^{\mathrm{wav}},\qquad \widehat{\mathrm{UOT}}_{nm}^{\mathrm{wav}}
\end{align*}
be the boundary-corrected wavelet plugin objects from Definition~\ref{def:uot_wavelet_plugin}.
{Assume Assumptions~\ref{assm:curvature} and \ref{assm:positivity}} for the unit cube, Assumption~\ref{assm:supp_global} holds automatically.
Assume moreover that $\alpha>1$ and $\alpha\notin\mathbb N$, that the densities satisfy
\begin{align}
    \Gamma^{-1} \le p(x),q(x) \le \Gamma,
    \qquad
    \|p\|_{C^{\alpha-1}(\Omega)}+\|q\|_{C^{\alpha-1}(\Omega)} \le M,
\end{align}
that the boundary-corrected wavelet basis in Definition~\ref{def:bc_wavelet_basis} has regularity strictly greater than $\alpha-1$, that the fitted wavelet problems are Monge-type, and write
\begin{align*}
    \hat\gamma_{0,nm}^{\mathrm{wav}}
    \coloneqq
    (\hat a_{nm}^{\mathrm{wav}})^2\hat\mu_n^{\mathrm{wav}}.
\end{align*}
Assume also that
    $2^{J_n} \asymp n^{{1} / ({d+2(\alpha-1)})}$ and
    $2^{J_m} \asymp m^{{1} / ({d+2(\alpha-1)})}$.
Then there exists a constant $C>0$ such that
\begin{align}
    \Ep\left[
        \int_{\Omega}
        \|\hat T_{nm}^{\mathrm{wav}}(x)-T_0(x)\|^2   d\mu(x)
    \right]
    &\le
    C\Bigl(M_\mu \mathfrak R_n^{\mathrm{wav}}(\alpha) + M_\nu \mathfrak R_m^{\mathrm{wav}}(\alpha)\Bigr),
\end{align}
\begin{align}
    \Ep\left[
        \int_{\Omega}
        |\hat\lambda_{nm}^{\mathrm{wav}}(x)-\lambda_0(x)|^2   d\mu(x)
    \right]
    &\le
    C\Bigl(M_\mu \mathfrak R_n^{\mathrm{wav}}(\alpha) + M_\nu \mathfrak R_m^{\mathrm{wav}}(\alpha)\Bigr),
\end{align}
and
\begin{align}
\begin{aligned}
    \Ep\left[
        \left|
            \widehat{\mathrm{UOT}}_{nm}^{\mathrm{wav}} - \mathrm{UOT}(\mu,\nu)
        \right|
    \right]
    &\le
    C\Bigl(
        \mathfrak L_n(\alpha)+\mathfrak L_m(\alpha) + M_\mu \mathfrak R_n^{\mathrm{wav}}(\alpha)+M_\nu \mathfrak R_m^{\mathrm{wav}}(\alpha)
    \Bigr).
\end{aligned}
\end{align}
\end{theorem}



\section{Proof of Theorem \ref{thm:main_wavelet_rates}} \label{sec:proof_wavelet}

\subsection{Convergence rates of wavelet density estimator}

We derive an upper bound for the convergence rate of the wavelet density estimator. Throughout this subsection we write
\begin{align}
    s \coloneqq \alpha-1 > 0,
\end{align}
and we assume that the population densities satisfy
\begin{align}
    \Gamma^{-1} \le p(x), q(x) \le \Gamma,
    \qquad
    \|p\|_{C^{s}(\Omega)} + \|q\|_{C^{s}(\Omega)} \le M,
\end{align}
for some $M,\Gamma>0$. We also define the rate templates
\begin{align}
    \mathfrak R_n^{\mathrm{wav}}(\alpha)
    \coloneqq
    \begin{cases}
        n^{-1}, & d=1,\\
        (\log n)^2 n^{-1}, & d=2,\\
        n^{{-2\alpha/(2(\alpha-1)+d)}}, & d\ge 3,
    \end{cases}
\end{align}
and the common $L^1$ density-estimation rate
\begin{align}
    \mathfrak L_n(\alpha)
    \coloneqq
    n^{{-(\alpha-1)/(2(\alpha-1)+d)}}.
\end{align}

\begin{theorem}[Wavelet density rates]
\label{thm:uot_plugin_density_rates}
Assume $\Omega=[0,1]^d$, $\alpha>1$, and $\alpha\notin\mathbb N$.
Assume moreover that
\begin{align*}
    \Gamma^{-1} \le p(x),q(x) \le \Gamma,
    \qquad
    \|p\|_{C^{\alpha-1}(\Omega)}+\|q\|_{C^{\alpha-1}(\Omega)} \le M,
\end{align*}
that the boundary-corrected wavelet basis in Definition~\ref{def:bc_wavelet_basis} has regularity strictly greater than $\alpha-1$, and choose
\begin{align*}
    2^{J_n} \asymp n^{{1/(d+2(\alpha-1))}},
    \qquad
    2^{J_m} \asymp m^{{1/(d+2(\alpha-1))}}.
\end{align*}
Then there exists a constant $C>0$, depending only on $d,\alpha,M,\Gamma$ and the chosen basis, such that
\begin{align}
\label{eq:uot_plugin_wavelet_w2_source}
    \Ep\bigl[W_2^2(\hat\mu_n^{\mathrm{wav}},\mu)\bigr]
    &\le C M_\mu\mathfrak R_n^{\mathrm{wav}}(\alpha),
    \\
    \Ep\bigl[W_2^2(\hat\nu_m^{\mathrm{wav}},\nu)\bigr]
    &\le C M_\nu\mathfrak R_m^{\mathrm{wav}}(\alpha),
\end{align}
and
\begin{align}
    \Ep\bigl[\|\hat p_n^{\mathrm{wav}}-p\|_{L^1(\Omega)}\bigr]
    &\le C\mathfrak L_n(\alpha),
    \\
\label{eq:uot_plugin_wavelet_l1_target}
    \Ep\bigl[\|\hat q_m^{\mathrm{wav}}-q\|_{L^1(\Omega)}\bigr]
    &\le C\mathfrak L_m(\alpha).
\end{align}
\end{theorem}


\begin{proof}[Proof of Theorem \ref{thm:uot_plugin_density_rates}]
We first prove the $L^1$ bounds.

\smallskip
\noindent\emph{Wavelet $L^1$ error.}
Let
\begin{align*}
    p_{J_n}^{\mathrm{bc}}(x)
    \coloneqq
    \sum_{\zeta\in\Phi^{\mathrm{bc}}}\alpha_\zeta\zeta(x)
    + \sum_{j=j_0}^{J_n}\sum_{\xi\in\Psi_j^{\mathrm{bc}}}\alpha_\xi\xi(x),
\end{align*}
where $\alpha_\zeta=\int_{\Omega}\zeta(x)p(x) dx$ and $\alpha_\xi=\int_{\Omega}\xi(x)p(x) dx$.
Let
\begin{align*}
    V_{J_n}^{\mathrm{bc}}
    \coloneqq
    \mathrm{span}(\Phi^{\mathrm{bc}})
    \oplus
    \bigoplus_{j=j_0}^{J_n}\mathrm{span}(\Psi_j^{\mathrm{bc}}),
\end{align*}
so that $p_{J_n}^{\mathrm{bc}}$ is exactly the $L^2$-orthogonal projection of $p$ onto $V_{J_n}^{\mathrm{bc}}$.
Because the chosen basis is the Cohen-Daubechies-Vial boundary-corrected basis from Definition~\ref{def:bc_wavelet_basis}, has regularity strictly larger than $s=\alpha-1$, and $\alpha\notin\mathbb N$ implies $C^s(\Omega)=\mathcal B^s_{\infty,\infty}(\Omega)$ with equivalent norms, the bias bound established in the proof of \cite[Lemma 30]{manole2024plugin} applies to the boundary-corrected projection $p_{J_n}^{\mathrm{bc}}=P_{V_{J_n}^{\mathrm{bc}}}p$ and gives
\begin{align*}
    \|p-p_{J_n}^{\mathrm{bc}}\|_{L^\infty(\Omega)}
    =
    \|p-P_{V_{J_n}^{\mathrm{bc}}}p\|_{L^\infty(\Omega)}
    \le
    C2^{-J_ns},
\end{align*}
and therefore also
\begin{align*}
    \|p-p_{J_n}^{\mathrm{bc}}\|_{L^1(\Omega)}
    \le
    C2^{-J_ns}.
\end{align*}

For the stochastic part, orthonormality and Parseval imply
\begin{align*}
    \Ep\bigl[\|\tilde p_n^{\mathrm{wav}}-p_{J_n}^{\mathrm{bc}}\|_{L^2(\Omega)}^2\bigr]
    &= \sum_{\zeta\in\Phi^{\mathrm{bc}}}\mathrm{Var}(\hat\alpha_\zeta)
    + \sum_{j=j_0}^{J_n}\sum_{\xi\in\Psi_j^{\mathrm{bc}}}\mathrm{Var}(\hat\alpha_\xi) \\
    &\le \frac{C}{n}\left(\#\Phi^{\mathrm{bc}} + \sum_{j=j_0}^{J_n}\#\Psi_j^{\mathrm{bc}}\right) \\
    &= \frac{C}{n}\left(2^{dj_0} + (2^d-1)\sum_{j=j_0}^{J_n}2^{dj}\right)
    \le C\frac{2^{dJ_n}}{n}.
\end{align*}
Since $|\Omega|=1$, $\|f\|_{L^1(\Omega)}\le \|f\|_{L^2(\Omega)}$; hence by Jensen,
\begin{align*}
    \Ep\bigl[\|\tilde p_n^{\mathrm{wav}}-p_{J_n}^{\mathrm{bc}}\|_{L^1(\Omega)}\bigr]
    \le
    \left(\Ep\bigl[\|\tilde p_n^{\mathrm{wav}}-p_{J_n}^{\mathrm{bc}}\|_{L^2(\Omega)}^2\bigr]\right)^{1/2}
    \le C\sqrt{\frac{2^{dJ_n}}{n}}.
\end{align*}
Therefore
\begin{align*}
    \Ep\bigl[\|\tilde p_n^{\mathrm{wav}}-p\|_{L^1(\Omega)}\bigr]
    \le C\left(2^{-J_ns}+\sqrt{\frac{2^{dJ_n}}{n}}\right).
\end{align*}
Balancing the two terms with $2^{J_n}\asymp n^{1/(d+2s)}$ yields
\begin{align*}
    \Ep\bigl[\|\tilde p_n^{\mathrm{wav}}-p\|_{L^1(\Omega)}\bigr]
    \le Cn^{{-s/(2s+d)}} = C\mathfrak L_n(\alpha).
\end{align*}
By Definition~\ref{def:bc_wavelet_basis}, the scaling space $\mathrm{span}(\Phi^{\mathrm{bc}})$ contains the constants, so there exist coefficients $(c_\zeta)_{\zeta\in\Phi^{\mathrm{bc}}}$ such that $1=\sum_{\zeta\in\Phi^{\mathrm{bc}}}c_\zeta\zeta$. Hence
\begin{align*}
    \int_{\Omega}\tilde p_n^{\mathrm{wav}}(x) dx
    &= \sum_{\zeta\in\Phi^{\mathrm{bc}}} c_\zeta\hat\alpha_\zeta
    = {1/n}\sum_{i=1}^n\sum_{\zeta\in\Phi^{\mathrm{bc}}}c_\zeta\zeta(X_i)
    = {1/n}\sum_{i=1}^n 1 = 1.
\end{align*}
The same positive-part renormalization argument as in the proof of Proposition~\ref{prop:cube_kernel_density_rates} (yielding \eqref{eq:cube_kernel_positive_part_l1}) therefore gives
\begin{align*}
    \Ep\bigl[\|\hat p_n^{\mathrm{wav}}-p\|_{L^1(\Omega)}\bigr]
    \le 3\Ep\bigl[\|\tilde p_n^{\mathrm{wav}}-p\|_{L^1(\Omega)}\bigr]
    \le C\mathfrak L_n(\alpha).
\end{align*}
The same argument gives the target bound for $\hat q_m^{\mathrm{wav}}$.

\smallskip
\noindent\emph{Wasserstein error.}
The probability measures induced by $\hat p_n^{\mathrm{wav}}$ and $\hat q_m^{\mathrm{wav}}$ are exactly the boundary-corrected one-sample wavelet estimators studied in \cite[Lemma 30]{manole2024plugin}, whose Wasserstein bound yields
\begin{align*}
    \Ep\bigl[W_2^2(\hat p_n^{\mathrm{wav}}dx,p dx)\bigr]
    \le C\mathfrak R_n^{\mathrm{wav}}(\alpha),
    \qquad
    \Ep\bigl[W_2^2(\hat q_m^{\mathrm{wav}}dx,q dx)\bigr]
    \le C\mathfrak R_m^{\mathrm{wav}}(\alpha).
\end{align*}
Finally, since the source and target masses are deterministic,
\begin{align*}
    W_2^2(\hat\mu_n^{\mathrm{wav}},\mu) = M_\mu W_2^2(\hat p_n^{\mathrm{wav}}dx,p dx),
    \qquad
    W_2^2(\hat\nu_m^{\mathrm{wav}},\nu) = M_\nu W_2^2(\hat q_m^{\mathrm{wav}}dx,q dx),
\end{align*}
which proves \eqref{eq:uot_plugin_wavelet_w2_source}-\eqref{eq:uot_plugin_wavelet_l1_target}.
\end{proof}


\begin{lemma}[High-probability lower bound for the wavelet density]
\label{lem:uot_wavelet_lower_bound_event}
Under the assumptions of Theorem~\ref{thm:uot_plugin_density_rates}, write $\underline p\coloneqq \inf_{x\in\Omega}p(x)>0$ and $\underline q\coloneqq \inf_{y\in\Omega}q(y)>0$, and define the event
\begin{align*}
    \mathcal G_{nm}^{\mathrm{wav}}
    \coloneqq
    \left\{
        \inf_{x\in\Omega}\tilde p_n^{\mathrm{wav}}(x)\ge \frac{\underline p}{2}
        \quad \text{and}\quad
        \inf_{y\in\Omega}\tilde q_m^{\mathrm{wav}}(y)\ge \frac{\underline q}{2}
    \right\}.
\end{align*}
Then for every $A>0$ there exists $C_A>0$ such that, for all sufficiently large $n,m$,
\begin{align}
\label{eq:uot_wavelet_lower_bound_event}
    \Pr\bigl(\mathcal G_{nm}^{\mathrm{wav},c}\bigr) \le C_A\bigl(n^{-A}+m^{-A}\bigr).
\end{align}
On $\mathcal G_{nm}^{\mathrm{wav}}$, the positive-part renormalization is inactive, so $\hat p_n^{\mathrm{wav}}=\tilde p_n^{\mathrm{wav}}$ and $\hat q_m^{\mathrm{wav}}=\tilde q_m^{\mathrm{wav}}$, and the source and target densities are dominated by the wavelet fits:
\begin{align}
\label{eq:wavelet_density_domination}
    d\mu \le \frac{2\Gamma}{\underline p} d\hat\mu_n^{\mathrm{wav}},
    \qquad
    d\nu \le \frac{2\Gamma}{\underline q} d\hat\nu_m^{\mathrm{wav}}.
\end{align}
\end{lemma}

\begin{proof}[Proof of Lemma \ref{lem:uot_wavelet_lower_bound_event}]
Let
\begin{align*}
    K_{J_n}^{\mathrm{wav}}(x,z)
    \coloneqq
    \sum_{\zeta\in\Phi^{\mathrm{bc}}}\zeta(x)\zeta(z)
    + \sum_{j=j_0}^{J_n}\sum_{\xi\in\Psi_j^{\mathrm{bc}}}\xi(x)\xi(z),
\end{align*}
so that $\tilde p_n^{\mathrm{wav}}(x) = \tfrac{1}{n}\sum_{i=1}^n K_{J_n}^{\mathrm{wav}}(x,X_i)$. Orthonormality and the bound $\#\Phi^{\mathrm{bc}}+\sum_{j=j_0}^{J_n}\#\Psi_j^{\mathrm{bc}}\le C 2^{J_nd}$ give
\begin{align*}
    \sup_{x\in\Omega} K_{J_n}^{\mathrm{wav}}(x,x) \le C 2^{J_nd},
\end{align*}
and Cauchy-Schwarz in the wavelet expansion yields $|K_{J_n}^{\mathrm{wav}}(x,z)|\le K_{J_n}^{\mathrm{wav}}(x,x)^{1/2}K_{J_n}^{\mathrm{wav}}(z,z)^{1/2}\le C 2^{J_nd}$. Since $\{K_{J_n}^{\mathrm{wav}}(x,\cdot):x\in\Omega\}$ is contained in $V_{J_n}^{\mathrm{bc}}$, a subspace of $L^2(\Omega)$ of dimension at most $C2^{J_nd}$, the same chaining and Bernstein argument used to establish \eqref{eq:bern} in the proof of Proposition~\ref{prop:cube_kernel_density_rates} yields
\begin{align*}
    \Pr\left(
        \|\tilde p_n^{\mathrm{wav}}-p_{J_n}^{\mathrm{bc}}\|_{L^\infty(\Omega)}
        > \frac{\underline p}{4}
    \right)
    \le C_A n^{-A}
\end{align*}
for every $A>0$ and all sufficiently large $n$. Since $\|p_{J_n}^{\mathrm{bc}}-p\|_{L^\infty(\Omega)}\le C 2^{-J_n s}\to 0$, the same conclusion holds for $\tilde p_n^{\mathrm{wav}}-p$ in place of $\tilde p_n^{\mathrm{wav}}-p_{J_n}^{\mathrm{bc}}$. The target bound is identical, and the union bound proves \eqref{eq:uot_wavelet_lower_bound_event}.

On $\mathcal G_{nm}^{\mathrm{wav}}$, $\tilde p_n^{\mathrm{wav}}\ge \underline p/2>0$ pointwise, so the positive part is the function itself and $\int_{\Omega}(\tilde p_n^{\mathrm{wav}})_+=\int_{\Omega}\tilde p_n^{\mathrm{wav}}=1$, giving $\hat p_n^{\mathrm{wav}}=\tilde p_n^{\mathrm{wav}}$. Hence $\hat\mu_n^{\mathrm{wav}}=M_\mu\hat p_n^{\mathrm{wav}}dx\ge M_\mu(\underline p/2)dx$, while $\mu = M_\mu p dx\le M_\mu \Gamma dx$, which combine to give the source bound in \eqref{eq:wavelet_density_domination}; the target bound is identical.
\end{proof}


\subsection{Proof of the upper bound}

Throughout this section we work on the Euclidean domain $\Omega=[0,1]^d$ and write
\begin{align*}
    \hat\mu_n^{\mathrm{wav}},\qquad \hat\nu_m^{\mathrm{wav}},\qquad \hat T_{nm}^{\mathrm{wav}},\qquad \hat\lambda_{nm}^{\mathrm{wav}},\qquad \widehat{\mathrm{UOT}}_{nm}^{\mathrm{wav}}
\end{align*}
for the boundary-corrected wavelet plugin objects from Definition~\ref{def:uot_wavelet_plugin}.

\begin{lemma}[Plug-in UOT excess identity]
\label{lem:uot_plugin_exact_excess}
Work on $\Omega=[0,1]^d$, and let
\[
    \hat\mu_n^{\mathrm{wav}},\ \hat\nu_m^{\mathrm{wav}},\ \widehat{\mathrm{UOT}}_{nm}^{\mathrm{wav}}
\]
be the boundary-corrected wavelet plugin objects from Definition~\ref{def:uot_wavelet_plugin}.
Let $(\varphi_0,\psi_0)$ be an optimal dual pair for $\mathrm{UOT}(\mu,\nu)$, define
\begin{align*}
    \zeta_0(x) \coloneqq -\bigl(e^{-\varphi_0(x)}-1\bigr),
    \qquad
    \xi_0(y) \coloneqq -\bigl(e^{-\psi_0(y)}-1\bigr),
\end{align*}
and let $\hat\gamma_{nm}^{\mathrm{wav}}$ be any optimal plan for $\widehat{\mathrm{UOT}}_{nm}^{\mathrm{wav}}$ with marginals $(\hat\gamma_{0,nm}^{\mathrm{wav}},\hat\gamma_{1,nm}^{\mathrm{wav}})$.
Define the oracle active fitted marginals by
\begin{align*}
    \gamma_{0,nm}^{\mathrm{wav},\mathrm{or}}
    &\coloneqq e^{-\varphi_0}\hat\mu_n^{\mathrm{wav}},
    \\
    \gamma_{1,nm}^{\mathrm{wav},\mathrm{or}}
    &\coloneqq e^{-\psi_0}\hat\nu_m^{\mathrm{wav}}.
\end{align*}
Then
\begin{equation}
\label{eq:uot_plugin_exact_excess}
\begin{aligned}
&\widehat{\mathrm{UOT}}_{nm}^{\mathrm{wav}}
- \int \zeta_0 d\hat\mu_n^{\mathrm{wav}}
- \int \xi_0 d\hat\nu_m^{\mathrm{wav}} \\
&\qquad=
\int_{\Omega\times\Omega}
\Bigl(
\tfrac{1}{2}\|x-y\|^2-\varphi_0(x)-\psi_0(y)
\Bigr)
  d\hat\gamma_{nm}^{\mathrm{wav}}(x,y) \\
&\qquad\quad
+ {D_\KL}(\hat\gamma_{0,nm}^{\mathrm{wav}}\mid \gamma_{0,nm}^{\mathrm{wav},\mathrm{or}})
+ {D_\KL}(\hat\gamma_{1,nm}^{\mathrm{wav}}\mid \gamma_{1,nm}^{\mathrm{wav},\mathrm{or}}).
\end{aligned}
\end{equation}
In particular,
\begin{align}
\label{eq:uot_plugin_exact_excess_nonnegative}
    \widehat{\mathrm{UOT}}_{nm}^{\mathrm{wav}}
    - \mathrm{UOT}(\mu,\nu)
    - \int \zeta_0 d(\hat\mu_n^{\mathrm{wav}}-\mu)
    - \int \xi_0 d(\hat\nu_m^{\mathrm{wav}}-\nu)
    \ge 0.
\end{align}
\end{lemma}


\begin{proof}[Proof of Lemma \ref{lem:uot_plugin_exact_excess}]
The derivation of \eqref{eq:uot_plugin_exact_excess} is identical to that of \eqref{eq:uot_cube_kernel_exact_excess} in the proof of Corollary~\ref{cor:uot_cube_kernel_rate_transfer}, with $(\hat\mu_n^{\mathrm{ker}},\hat\nu_m^{\mathrm{ker}},\hat\gamma_{nm}^{\mathrm{ker}})$ replaced by $(\hat\mu_n^{\mathrm{wav}},\hat\nu_m^{\mathrm{wav}},\hat\gamma_{nm}^{\mathrm{wav}})$: it relies only on optimality of $\hat\gamma_{nm}^{\mathrm{wav}}$ for $\mathrm{UOT}(\hat\mu_n^{\mathrm{wav}},\hat\nu_m^{\mathrm{wav}})$ and the rewriting of the source and target $D_\KL$ terms with reference measures $e^{-\varphi_0}\hat\mu_n^{\mathrm{wav}}$ and $e^{-\psi_0}\hat\nu_m^{\mathrm{wav}}$, neither of which depends on the specific construction of the fitted measures.

Finally, duality gives
\begin{align*}
    \mathrm{UOT}(\mu,\nu)
    = \int \zeta_0 d\mu + \int \xi_0 d\nu,
\end{align*}
and \eqref{eq:uot_plugin_exact_excess_nonnegative} follows from the nonnegativity of the dual slack and of the {$D_\KL$ terms}.
\end{proof}

\begin{theorem}[Plug-in UOT stability bound]
\label{thm:uot_plugin_stability_value}
Assume Assumptions~\ref{assm:curvature} and \ref{assm:positivity}, and work on $\Omega=[0,1]^d$.
Then
\begin{equation}
\label{eq:uot_plugin_value_expansion}
\begin{aligned}
0
&\le
\widehat{\mathrm{UOT}}_{nm}^{\mathrm{wav}}
- \mathrm{UOT}(\mu,\nu)
- \int \zeta_0 d(\hat\mu_n^{\mathrm{wav}}-\mu)
- \int \xi_0 d(\hat\nu_m^{\mathrm{wav}}-\nu) \\
&\le
C_\Lambda\Bigl(
W_2^2(\hat\mu_n^{\mathrm{wav}},\mu)
+
W_2^2(\hat\nu_m^{\mathrm{wav}},\nu)
\Bigr).
\end{aligned}
\end{equation}
In particular,
\begin{equation}
\label{eq:uot_plugin_value_absolute}
\begin{aligned}
\left|
\widehat{\mathrm{UOT}}_{nm}^{\mathrm{wav}} - \mathrm{UOT}(\mu,\nu)
\right|
&\le
M_\mu\|\zeta_0\|_{L^\infty(\Omega)}
\|\hat p_n^{\mathrm{wav}}-p\|_{L^1(\Omega)} \\
&\quad+
M_\nu\|\xi_0\|_{L^\infty(\Omega)}
\|\hat q_m^{\mathrm{wav}}-q\|_{L^1(\Omega)} \\
&\quad+
C_\Lambda\Bigl(
W_2^2(\hat\mu_n^{\mathrm{wav}},\mu)
+
W_2^2(\hat\nu_m^{\mathrm{wav}},\nu)
\Bigr).
\end{aligned}
\end{equation}
\end{theorem}

\begin{proof}[Proof of Theorem \ref{thm:uot_plugin_stability_value}]
The lower bound in \eqref{eq:uot_plugin_value_expansion} is exactly \eqref{eq:uot_plugin_exact_excess_nonnegative} from Lemma~\ref{lem:uot_plugin_exact_excess}.
Since $\Omega=[0,1]^d$ is compact, convex, and satisfies the interior cone condition, Proposition~\ref{prop:two_sample_stability} applies to the pair $(\hat\mu_n^{\mathrm{wav}},\hat\nu_m^{\mathrm{wav}})$, which has the same masses $M_\mu$ and $M_\nu$ as $(\mu,\nu)$. This gives the upper bound in \eqref{eq:uot_plugin_value_expansion}.

For \eqref{eq:uot_plugin_value_absolute}, write
\begin{align*}
    A_{n,m}^{\mathrm{wav}}
    \coloneqq
    \int \zeta_0 d(\hat\mu_n^{\mathrm{wav}}-\mu)
    + \int \xi_0 d(\hat\nu_m^{\mathrm{wav}}-\nu).
\end{align*}
By \eqref{eq:uot_plugin_value_expansion},
\begin{align*}
    \widehat{\mathrm{UOT}}_{nm}^{\mathrm{wav}} - \mathrm{UOT}(\mu,\nu)
    = A_{n,m}^{\mathrm{wav}} + R_{n,m}^{\mathrm{wav}},
    \qquad
    0\le R_{n,m}^{\mathrm{wav}} \le C_\Lambda\Bigl(W_2^2(\hat\mu_n^{\mathrm{wav}},\mu)+W_2^2(\hat\nu_m^{\mathrm{wav}},\nu)\Bigr).
\end{align*}
Therefore
\begin{align*}
    \left|\widehat{\mathrm{UOT}}_{nm}^{\mathrm{wav}} - \mathrm{UOT}(\mu,\nu)\right|
    \le |A_{n,m}^{\mathrm{wav}}| + C_\Lambda\Bigl(W_2^2(\hat\mu_n^{\mathrm{wav}},\mu)+W_2^2(\hat\nu_m^{\mathrm{wav}},\nu)\Bigr).
\end{align*}
Since $\hat\mu_n^{\mathrm{wav}}=M_\mu\hat p_n^{\mathrm{wav}}dx$ and $\mu=M_\mu p dx$,
\begin{align*}
    \left|\int \zeta_0 d(\hat\mu_n^{\mathrm{wav}}-\mu)\right|
    &= M_\mu\left|\int_{\Omega} \zeta_0(x)(\hat p_n^{\mathrm{wav}}(x)-p(x)) dx\right| \\
    &\le M_\mu\|\zeta_0\|_{L^\infty(\Omega)}\|\hat p_n^{\mathrm{wav}}-p\|_{L^1(\Omega)}.
\end{align*}
The target term is treated in the same way, giving \eqref{eq:uot_plugin_value_absolute}.
\end{proof}

\begin{theorem}[Plug-in transport map risk]
\label{thm:uot_plugin_map_risk}
Assume Assumptions~\ref{assm:curvature} and \ref{assm:positivity}, and work on $\Omega=[0,1]^d$.
Suppose that the fitted problem between $(\hat\mu_n^{\mathrm{wav}},\hat\nu_m^{\mathrm{wav}})$ is Monge-type, with optimal plan
\begin{align*}
    \hat\gamma_{nm}^{\mathrm{wav}} = (\mathrm{id},\hat T_{nm}^{\mathrm{wav}})_\#\hat\gamma_{0,nm}^{\mathrm{wav}},
    \qquad
    \hat\gamma_{0,nm}^{\mathrm{wav}} = (\hat a_{nm}^{\mathrm{wav}})^2\hat\mu_n^{\mathrm{wav}}
    = e^{-\hat\varphi_{nm}^{\mathrm{wav}}}\hat\mu_n^{\mathrm{wav}}.
\end{align*}
Then
\begin{equation}
\label{eq:uot_plugin_map_risk_main}
\begin{aligned}
\frac{\kappa}{2}
\int_{\Omega}
\|\hat T_{nm}^{\mathrm{wav}}(x)-T_0(x)\|^2 d\hat\gamma_{0,nm}^{\mathrm{wav}}(x)
&\le
\widehat{\mathrm{UOT}}_{nm}^{\mathrm{wav}}
- \mathrm{UOT}(\mu,\nu) \\
&\quad - \int \zeta_0 d(\hat\mu_n^{\mathrm{wav}}-\mu)
- \int \xi_0 d(\hat\nu_m^{\mathrm{wav}}-\nu) \\
&\le
C_\Lambda\Bigl(
W_2^2(\hat\mu_n^{\mathrm{wav}},\mu)
+
W_2^2(\hat\nu_m^{\mathrm{wav}},\nu)
\Bigr).
\end{aligned}
\end{equation}
Consequently,
\begin{align}
\label{eq:uot_plugin_map_risk_rate}
    \int_{\Omega}
    \|\hat T_{nm}^{\mathrm{wav}}(x)-T_0(x)\|^2 d\hat\gamma_{0,nm}^{\mathrm{wav}}(x)
    \le
    \frac{2C_\Lambda}{\kappa}
    \Bigl(
        W_2^2(\hat\mu_n^{\mathrm{wav}},\mu)
        +
        W_2^2(\hat\nu_m^{\mathrm{wav}},\nu)
    \Bigr).
\end{align}
Moreover, the active-mass mismatch is controlled by the same remainder:
\begin{equation}
\label{eq:uot_plugin_active_mass_mismatch}
\begin{aligned}
&{D_\KL}(\hat\gamma_{0,nm}^{\mathrm{wav}}\mid e^{-\varphi_0}\hat\mu_n^{\mathrm{wav}})
+
{D_\KL}(\hat\gamma_{1,nm}^{\mathrm{wav}}\mid e^{-\psi_0}\hat\nu_m^{\mathrm{wav}}) \le
C_\Lambda\Bigl(
W_2^2(\hat\mu_n^{\mathrm{wav}},\mu)
+
W_2^2(\hat\nu_m^{\mathrm{wav}},\nu)
\Bigr),
\end{aligned}
\end{equation}
and the active-source factor satisfies
\begin{align}
\label{eq:uot_plugin_active_factor_bound}
    \int_{\Omega}
    |\hat a_{nm}^{\mathrm{wav}}(x)-a_0(x)|^2 d\hat\mu_n^{\mathrm{wav}}(x)
    \le
    {D_\KL}(\hat\gamma_{0,nm}^{\mathrm{wav}}\mid e^{-\varphi_0}\hat\mu_n^{\mathrm{wav}}).
\end{align}
Finally, there exists a constant $C_{\mathrm{tr}}>0$, depending only on $\|\varphi_0\|_{L^\infty(\Omega)}$ and $\mathrm{diam}(\Omega)$, such that
\begin{align}
\label{eq:uot_plugin_T_to_mu_n}
    \int_{\Omega}
    \|\hat T_{nm}^{\mathrm{wav}}(x)-T_0(x)\|^2 d\hat\mu_n^{\mathrm{wav}}(x)
    \le
    C_{\mathrm{tr}}(2/\kappa+1) C_\Lambda\Bigl(
        W_2^2(\hat\mu_n^{\mathrm{wav}},\mu)
        +
        W_2^2(\hat\nu_m^{\mathrm{wav}},\nu)
    \Bigr).
\end{align}
\end{theorem}

\begin{proof}[Proof of Theorem \ref{thm:uot_plugin_map_risk}]
By Lemma~\ref{lem:uot_plugin_exact_excess},
\begin{align}
\label{eq:uot_plugin_map_risk_proof_start}
\begin{aligned}
&\widehat{\mathrm{UOT}}_{nm}^{\mathrm{wav}}
- \mathrm{UOT}(\mu,\nu)
- \int \zeta_0 d(\hat\mu_n^{\mathrm{wav}}-\mu)
- \int \xi_0 d(\hat\nu_m^{\mathrm{wav}}-\nu) \\
&\qquad=
\int
\Bigl(
\tfrac{1}{2}\|x-\hat T_{nm}^{\mathrm{wav}}(x)\|^2
- \varphi_0(x)-\psi_0(\hat T_{nm}^{\mathrm{wav}}(x))
\Bigr)
  d\hat\gamma_{0,nm}^{\mathrm{wav}}(x) \\
&\qquad\quad + {D_\KL}(\hat\gamma_{0,nm}^{\mathrm{wav}}\mid e^{-\varphi_0}\hat\mu_n^{\mathrm{wav}})
+ {D_\KL}(\hat\gamma_{1,nm}^{\mathrm{wav}}\mid e^{-\psi_0}\hat\nu_m^{\mathrm{wav}}).
\end{aligned}
\end{align}
Lemma~\ref{lem:uot_gap_sufficient} gives
\begin{align*}
    \tfrac{1}{2}\|x-\hat T_{nm}^{\mathrm{wav}}(x)\|^2
    - \varphi_0(x)-\psi_0(\hat T_{nm}^{\mathrm{wav}}(x))
    \ge \frac{\kappa}{2}\|\hat T_{nm}^{\mathrm{wav}}(x)-T_0(x)\|^2
\end{align*}
for all $x\in\Omega$.
Substituting this lower bound into \eqref{eq:uot_plugin_map_risk_proof_start} and discarding the {nonnegative $D_\KL$ terms} yields the left inequality in \eqref{eq:uot_plugin_map_risk_main}. The right inequality in \eqref{eq:uot_plugin_map_risk_main} is exactly Theorem~\ref{thm:uot_plugin_stability_value}. Hence \eqref{eq:uot_plugin_map_risk_rate} follows immediately.
Keeping the {$D_\KL$ terms} in \eqref{eq:uot_plugin_map_risk_proof_start} and using once again the upper bound from Theorem~\ref{thm:uot_plugin_stability_value} proves \eqref{eq:uot_plugin_active_mass_mismatch}.

For \eqref{eq:uot_plugin_active_factor_bound}, the active-source identities $\hat\gamma_{0,nm}^{\mathrm{wav}}=(\hat a_{nm}^{\mathrm{wav}})^2\hat\mu_n^{\mathrm{wav}}$ and $e^{-\varphi_0}\hat\mu_n^{\mathrm{wav}}=a_0^2\hat\mu_n^{\mathrm{wav}}$ allow us to apply Lemma~\ref{lem:uot_kl_sqrt_lower_bound} pointwise with $a=(\hat a_{nm}^{\mathrm{wav}})^2$ and $b=a_0^2$; integrating against $\hat\mu_n^{\mathrm{wav}}$ yields
\begin{align*}
    \int_{\Omega}|\hat a_{nm}^{\mathrm{wav}}(x)-a_0(x)|^2 d\hat\mu_n^{\mathrm{wav}}(x)
    \le {D_\KL}(\hat\gamma_{0,nm}^{\mathrm{wav}}\mid e^{-\varphi_0}\hat\mu_n^{\mathrm{wav}}).
\end{align*}

To prove \eqref{eq:uot_plugin_T_to_mu_n}, we transfer the map error from the active marginal to $\hat\mu_n^{\mathrm{wav}}$. Since $\varphi_0$ is continuous and $\Omega$ is compact, there exists a constant $w_->0$ such that $a_0(x)^2=e^{-\varphi_0(x)}\ge w_-$ for all $x\in\Omega$. Because $\hat T_{nm}^{\mathrm{wav}}$ and $T_0$ both take values in $\Omega=[0,1]^d$, we also have $\|\hat T_{nm}^{\mathrm{wav}}(x)-T_0(x)\|^2\le \mathrm{diam}(\Omega)^2 = d$. Consequently,
\begin{align*}
    \|\hat T_{nm}^{\mathrm{wav}}(x)-T_0(x)\|^2
    &\le \frac{1}{w_-} a_0(x)^2 \|\hat T_{nm}^{\mathrm{wav}}(x)-T_0(x)\|^2 \\
    &\le \frac{2}{w_-}\hat a_{nm}^{\mathrm{wav}}(x)^2\|\hat T_{nm}^{\mathrm{wav}}(x)-T_0(x)\|^2 \\
    &\quad + \frac{2}{w_-}|\hat a_{nm}^{\mathrm{wav}}(x)-a_0(x)|^2\|\hat T_{nm}^{\mathrm{wav}}(x)-T_0(x)\|^2 \\
    &\le
    C_{\mathrm{tr}} \hat a_{nm}^{\mathrm{wav}}(x)^2 \|\hat T_{nm}^{\mathrm{wav}}(x)-T_0(x)\|^2
    + C_{\mathrm{tr}} |\hat a_{nm}^{\mathrm{wav}}(x)-a_0(x)|^2,
\end{align*}
for a constant $C_{\mathrm{tr}}>0$ depending only on $w_-$ and $\mathrm{diam}(\Omega)$. Integrating against $d\hat\mu_n^{\mathrm{wav}}$ and using $d\hat\gamma_{0,nm}^{\mathrm{wav}}=\hat a_{nm}^{\mathrm{wav}}(x)^2 d\hat\mu_n^{\mathrm{wav}}(x)$ together with \eqref{eq:uot_plugin_map_risk_rate}, \eqref{eq:uot_plugin_active_mass_mismatch}, and \eqref{eq:uot_plugin_active_factor_bound} yields \eqref{eq:uot_plugin_T_to_mu_n}.
\end{proof}

\begin{corollary}[Plug-in growth map risk]
\label{cor:uot_plugin_growth_risk}
Assume that $\Omega=[0,1]^d$, that Assumptions~\ref{assm:curvature} and \ref{assm:positivity} hold.
Let
\[
    \hat\mu_n^{\mathrm{wav}},\ \hat\nu_m^{\mathrm{wav}},\ \hat T_{nm}^{\mathrm{wav}},\ \hat a_{nm}^{\mathrm{wav}},\ \hat\lambda_{nm}^{\mathrm{wav}}
\]
be the boundary-corrected wavelet plugin objects from Definition~\ref{def:uot_wavelet_plugin}. Suppose that the fitted problem between $(\hat\mu_n^{\mathrm{wav}},\hat\nu_m^{\mathrm{wav}})$ is Monge-type, with optimal plan
\begin{align*}
    \hat\gamma_{nm}^{\mathrm{wav}} = (\mathrm{id},\hat T_{nm}^{\mathrm{wav}})_\#\hat\gamma_{0,nm}^{\mathrm{wav}},
    \qquad
    \hat\gamma_{0,nm}^{\mathrm{wav}} = (\hat a_{nm}^{\mathrm{wav}})^2\hat\mu_n^{\mathrm{wav}}.
\end{align*}
Then
\begin{align}
\label{eq:uot_plugin_growth_risk}
    \int_{\Omega}
    |\hat\lambda_{nm}^{\mathrm{wav}}(x)-\lambda_0(x)|^2 d\hat\mu_n^{\mathrm{wav}}(x)
    \le
    C_\Lambda\Bigl(
        W_2^2(\hat\mu_n^{\mathrm{wav}},\mu)
        +
        W_2^2(\hat\nu_m^{\mathrm{wav}},\nu)
    \Bigr).
\end{align}
\end{corollary}

\begin{proof}[Proof of Corollary \ref{cor:uot_plugin_growth_risk}]
By Lemma~\ref{lem:continuous_gh_growth_transfer} applied with $\hat\eta=\hat\mu_n^{\mathrm{wav}}$, $\hat a=\hat a_{nm}^{\mathrm{wav}}$, $\hat T=\hat T_{nm}^{\mathrm{wav}}$, and $\hat\gamma_0=\hat\gamma_{0,nm}^{\mathrm{wav}}$,
\begin{align*}
    \int_{\Omega}|\hat\lambda_{nm}^{\mathrm{wav}}-\lambda_0|^2 d\hat\mu_n^{\mathrm{wav}}
    \le
    C\int_{\Omega}|\hat a_{nm}^{\mathrm{wav}}-a_0|^2 d\hat\mu_n^{\mathrm{wav}}
    +C\int_{\Omega}\|\hat T_{nm}^{\mathrm{wav}}-T_0\|^2 d\hat\gamma_{0,nm}^{\mathrm{wav}}.
\end{align*}
The first term on the right is bounded by combining \eqref{eq:uot_plugin_active_factor_bound} with \eqref{eq:uot_plugin_active_mass_mismatch}, and the second is bounded by \eqref{eq:uot_plugin_map_risk_rate}.
\end{proof}

\begin{corollary}[Wavelet plugin rate transfer]
\label{cor:uot_plugin_rate_transfer}
Assume that $\Omega=[0,1]^d$, that $c(x,y)=\tfrac{1}{2}\|x-y\|^2$, and that Assumptions~\ref{assm:curvature} and \ref{assm:positivity} hold, together with the smoothness assumptions of Theorem~\ref{thm:uot_plugin_density_rates} so that the high-probability event $\mathcal G_{nm}^{\mathrm{wav}}$ of Lemma~\ref{lem:uot_wavelet_lower_bound_event} is available.
Let
\[
    \hat\mu_n^{\mathrm{wav}},\ \hat\nu_m^{\mathrm{wav}},\ \hat T_{nm}^{\mathrm{wav}},\ \hat\lambda_{nm}^{\mathrm{wav}},\ \widehat{\mathrm{UOT}}_{nm}^{\mathrm{wav}}
\]
be the boundary-corrected wavelet plugin objects from Definition~\ref{def:uot_wavelet_plugin}, and write
\begin{align*}
    \hat\gamma_{0,nm}^{\mathrm{wav}}
    \coloneqq
    (\hat a_{nm}^{\mathrm{wav}})^2\hat\mu_n^{\mathrm{wav}}.
\end{align*}
Suppose that the fitted problem between $(\hat\mu_n^{\mathrm{wav}},\hat\nu_m^{\mathrm{wav}})$ is Monge-type.
Assume that for some deterministic sequences $r_{n,\mathrm{wav}},r_{m,\mathrm{wav}}\ge 0$, where $r_{n,\mathrm{wav}}\ge 1/n$ and $r_{m,\mathrm{wav}}\ge 1/m$,
\begin{align*}
    \Ep\bigl[W_2^2(\hat\mu_n^{\mathrm{wav}},\mu)\bigr] \le r_{n,\mathrm{wav}},
    \qquad
    \Ep\bigl[W_2^2(\hat\nu_m^{\mathrm{wav}},\nu)\bigr] \le r_{m,\mathrm{wav}}.
\end{align*}
Then there exists a constant $C>0$, depending only on the constants in Assumptions~\ref{assm:curvature} and \ref{assm:positivity}, such that
\begin{align}
\label{eq:uot_plugin_rate_transfer_map}
    \Ep\left[
        \int_{\Omega}
        \|\hat T_{nm}^{\mathrm{wav}}(x)-T_0(x)\|^2 d\mu(x)
    \right]
    \le
    C(r_{n,\mathrm{wav}}+r_{m,\mathrm{wav}}),
\end{align}
\begin{align}
\label{eq:uot_plugin_rate_transfer_lambda}
    \Ep\left[
        \int_{\Omega}
        |\hat\lambda_{nm}^{\mathrm{wav}}(x)-\lambda_0(x)|^2 d\mu(x)
    \right]
    \le
    C(r_{n,\mathrm{wav}}+r_{m,\mathrm{wav}}),
\end{align}
and
\begin{align}
\label{eq:uot_plugin_rate_transfer_value}
    \Ep\left[
        \left|
        \widehat{\mathrm{UOT}}_{nm}^{\mathrm{wav}}
        - \mathrm{UOT}(\mu,\nu)
        - \int \zeta_0 d(\hat\mu_n^{\mathrm{wav}}-\mu)
        - \int \xi_0 d(\hat\nu_m^{\mathrm{wav}}-\nu)
        \right|
    \right]
    \le
    C_\Lambda(r_{n,\mathrm{wav}}+r_{m,\mathrm{wav}}).
\end{align}
If, in addition,
\begin{align*}
    \Ep\bigl[\|\hat p_n^{\mathrm{wav}}-p\|_{L^1(\Omega)}\bigr] \le \ell_{n,\mathrm{wav}},
    \qquad
    \Ep\bigl[\|\hat q_m^{\mathrm{wav}}-q\|_{L^1(\Omega)}\bigr] \le \ell_{m,\mathrm{wav}},
\end{align*}
then
\begin{align}
\label{eq:uot_plugin_rate_transfer_abs_value}
\begin{aligned}
    \Ep\left[
        \left|
        \widehat{\mathrm{UOT}}_{nm}^{\mathrm{wav}} - \mathrm{UOT}(\mu,\nu)
        \right|
    \right]
    &\le M_\mu\|\zeta_0\|_{L^\infty}\ell_{n,\mathrm{wav}}
    + M_\nu\|\xi_0\|_{L^\infty}\ell_{m,\mathrm{wav}} \\
    &\quad + C_\Lambda(r_{n,\mathrm{wav}}+r_{m,\mathrm{wav}}).
\end{aligned}
\end{align}
\end{corollary}


\begin{proof}[Proof of Corollary \ref{cor:uot_plugin_rate_transfer}]
We first prove \eqref{eq:uot_plugin_rate_transfer_map}. On the high-probability event $\mathcal G_{nm}^{\mathrm{wav}}$ of Lemma~\ref{lem:uot_wavelet_lower_bound_event}, the density-domination relation \eqref{eq:wavelet_density_domination} gives
\begin{align*}
    \int_{\Omega}\|\hat T_{nm}^{\mathrm{wav}}(x)-T_0(x)\|^2 d\mu(x)
    \le \frac{2\Gamma}{\underline p}
    \int_{\Omega}\|\hat T_{nm}^{\mathrm{wav}}(x)-T_0(x)\|^2 d\hat\mu_n^{\mathrm{wav}}(x).
\end{align*}
Taking expectations in this inequality on $\mathcal G_{nm}^{\mathrm{wav}}$ and applying \eqref{eq:uot_plugin_T_to_mu_n} together with the assumed Wasserstein bounds yields
\begin{align*}
    \Ep\left[
        \int_{\Omega}\|\hat T_{nm}^{\mathrm{wav}}(x)-T_0(x)\|^2 d\mu(x) 
        \mathbf 1_{\mathcal G_{nm}^{\mathrm{wav}}}
    \right]
    \le C(r_{n,\mathrm{wav}}+r_{m,\mathrm{wav}}).
\end{align*}
On the complement, both $\hat T_{nm}^{\mathrm{wav}}$ and $T_0$ take values in $\Omega=[0,1]^d$, so $\|\hat T_{nm}^{\mathrm{wav}}(x)-T_0(x)\|^2\le d$, and \eqref{eq:uot_wavelet_lower_bound_event} gives
\begin{align*}
    \Ep\left[
        \int_{\Omega}\|\hat T_{nm}^{\mathrm{wav}}(x)-T_0(x)\|^2 d\mu(x) 
        \mathbf 1_{\mathcal G_{nm}^{\mathrm{wav},c}}
    \right]
    \le d M_\mu\Pr(\mathcal G_{nm}^{\mathrm{wav},c})
    \le C_A(n^{-A}+m^{-A}).
\end{align*}
Choosing $A\ge 2$, this contribution is $O(n^{-2}+m^{-2})$ and is absorbed by $C(r_{n,\mathrm{wav}}+r_{m,\mathrm{wav}})$, proving \eqref{eq:uot_plugin_rate_transfer_map}.

For the growth factor, the same density-domination on $\mathcal G_{nm}^{\mathrm{wav}}$ together with \eqref{eq:uot_plugin_growth_risk} gives
\begin{align*}
    \int_{\Omega}|\hat\lambda_{nm}^{\mathrm{wav}}(x)-\lambda_0(x)|^2 d\mu(x)
    \le \frac{2\Gamma}{\underline p} 
    C_\Lambda\Bigl(W_2^2(\hat\mu_n^{\mathrm{wav}},\mu)+W_2^2(\hat\nu_m^{\mathrm{wav}},\nu)\Bigr)
    \quad \text{on }\mathcal G_{nm}^{\mathrm{wav}}.
\end{align*}
On the complement, $\hat\lambda_{nm}^{\mathrm{wav}}$ is clipped to $[w_-,w_+]$ and $\lambda_0$ is bounded above and below on $\Omega$ by Assumption~\ref{assm:curvature}, so $|\hat\lambda_{nm}^{\mathrm{wav}}-\lambda_0|^2$ is uniformly bounded; the same union-bound argument as above absorbs the bad-event contribution into $C(r_{n,\mathrm{wav}}+r_{m,\mathrm{wav}})$, proving \eqref{eq:uot_plugin_rate_transfer_lambda}.

For \eqref{eq:uot_plugin_rate_transfer_value}, \eqref{eq:uot_plugin_value_expansion} implies
\begin{align*}
    0
    &\le
    \widehat{\mathrm{UOT}}_{nm}^{\mathrm{wav}}
    - \mathrm{UOT}(\mu,\nu)
    - \int \zeta_0 d(\hat\mu_n^{\mathrm{wav}}-\mu)
    - \int \xi_0 d(\hat\nu_m^{\mathrm{wav}}-\nu)\\
    &\le
    C_\Lambda\Bigl(W_2^2(\hat\mu_n^{\mathrm{wav}},\mu)+W_2^2(\hat\nu_m^{\mathrm{wav}},\nu)\Bigr),
\end{align*}
so taking expectations gives the claim. Finally, taking expectations in \eqref{eq:uot_plugin_value_absolute} and using the assumed $L^1$ bounds yields \eqref{eq:uot_plugin_rate_transfer_abs_value}.
\end{proof}


\begin{proof}[Proof of Theorem \ref{thm:main_wavelet_rates}]
Apply Corollary~\ref{cor:uot_plugin_rate_transfer} with
\begin{align*}
    r_{n,\mathrm{wav}} = C M_\mu\mathfrak R_n^{\mathrm{wav}}(\alpha),
    \quad
    r_{m,\mathrm{wav}} = C M_\nu\mathfrak R_m^{\mathrm{wav}}(\alpha),
    \quad
    \ell_{n,\mathrm{wav}} = C\mathfrak L_n(\alpha),
    \quad
    \ell_{m,\mathrm{wav}} = C\mathfrak L_m(\alpha),
\end{align*}
as supplied by Theorem~\ref{thm:uot_plugin_density_rates}. The constants $M_\mu\|\zeta_0\|_{L^\infty}$ and $M_\nu\|\xi_0\|_{L^\infty}$ are absorbed into the generic constant $C$.
\end{proof}




\section{Proof of Theorem~\ref{thm:lb-ubot}}

\subsection{Statistical model and lower-bound class}

Let $\Omega=[0,1]^d$. We assume that $M_\mu,M_\nu$ are known and observe two independent samples $X_1,\dots,X_n\overset{\mathrm{i.i.d.}}{\sim}\bar\mu\coloneqq\mu/M_\mu$ and $Y_1,\dots,Y_n\overset{\mathrm{i.i.d.}}{\sim}\bar\nu\coloneqq\nu/M_\nu$. We write $\mathbb P_{\mu,\nu}^n$ and $\mathbb E_{\mu,\nu}^n$ for the joint law and expectation. The KL divergence between two pairs of source and target measures $(\mu,\nu)$ and $(\mu',\mu')$
\begin{align}
    \label{eq:iid_kl}
    D_\KL(\mathbb P_{\mu,\nu}^n\|\mathbb P_{\mu',\nu'}^n)
    =
    n\bigl[D_\KL(\bar\mu\|\bar\mu')+D_\KL(\bar\nu\|\bar\nu')\bigr].
\end{align}

For the appendix proof, write $\mathcal U_\alpha$ for $\mathcal U_\alpha(M,B,\Lambda)$ from Section~\ref{sec:minimax_lower_bound}. We introduce the shorthand minimax risks
\begin{align}
    \mathfrak M_n^T(\hat T)
    &\coloneqq
    \sup_{(\mu,\nu)\in\mathcal U_\alpha}
    \mathbb E_{\mu,\nu}^n\left[
        \int_{[0,1]^d}\|\hat T(x)-T_0(x)\|^2 d\mu(x)
    \right],\\
    \mathfrak M_n^\lambda(\hat\lambda)
    &\coloneqq
    \sup_{(\mu,\nu)\in\mathcal U_\alpha}
    \mathbb E_{\mu,\nu}^n\left[
        \int_{[0,1]^d}|\hat\lambda(x)-\lambda_0(x)|^2 d\mu(x)
    \right].
\end{align}

\subsection{Tools for minimax lower-bound}

Throughout, we use standard tools for minimax lower bounds. See, e.g., \cite[Theorem~2.2 and~2.5]{tsybakov2009nonparametric}, restated here for convenience.

\begin{external}[Lower bound from two hypotheses, \cite{tsybakov2009nonparametric}]\label{thm:tsy-twopoint}
    If probability measures $P_0,P_1$ satisfy $D_\KL(P_1\|P_0)\le \alpha<\infty$ and $d(\theta_0,\theta_1)\ge 2s$, then there exists a constant $c_\alpha>0$ such that $\inf_{\hat\theta}\max_{i=0,1}P_i(d(\hat\theta,\theta_i)\ge s)\ge c_\alpha$.
\end{external}
\begin{external}[Lower bound from multiple hypotheses, \cite{tsybakov2009nonparametric}]\label{thm:tsy-fano}
    Let $(\Theta,d)$ be a pseudometric space and $\{P_\theta\}_{\theta\in\Theta}$ a family of probability measures. Suppose there exist $\theta_0,\theta_1,\dots,\theta_K\in\Theta$ and $s>0$ such that $P_{\theta_k} \ll P_{\theta_0}$ for $k=1,\ldots,K$, $d(\theta_k,\theta_{k'})\ge 2s$ for $0\le k\ne k'\le K$, and $K^{-1}\sum_{k=1}^KD_\KL(P_{\theta_k}\|P_{\theta_0})\le (\log K)/9$, with $K\ge2$. Then there exists a constant $c>0$ such that $\inf_{\hat\theta}\sup_k P_{\theta_k}(d(\hat\theta,\theta_k)\ge s)\ge c$.
\end{external}
\begin{external}[Varshamov-Gilbert bound, {\cite[Lemma~2.9]{tsybakov2009nonparametric}}]\label{lem:vg}
    For every integer $N_J\ge 8$ there exist binary strings $\omega^{(0)},\omega^{(1)},\dots,\omega^{(K)}\in\{0,1\}^{N_J}$ with $\omega^{(0)}=\mathbf 0$, cardinality $K\ge 2^{N_J/8}$, and pairwise Hamming separation
    \begin{align}
        \|\omega^{(k)}-\omega^{(k')}\|_2^2 = \#\{i:\omega^{(k)}_i\ne\omega^{(k')}_i\} \ge N_J/8 \qquad \text{for all }0\le k\ne k'\le K.
    \end{align}
\end{external}

\subsection{Proof of Theorem~\ref{thm:lb-ubot}}
\begin{proof}
To establish the minimax lower bound, we will appeal to both External result~\ref{thm:tsy-twopoint} and \ref{thm:tsy-fano}; both involve constructing a collection of hypotheses $(\mu,\nu_k)\in\mathcal U_\alpha$ for which estimation is hard. Throughout, $O_\infty(r_m)$ denotes a function $f:\Omega\to\R$ with $\|f\|_{L^\infty(\Omega)}\le C r_m$, where $C$ is a constant independent of $k$ and $J$. For any matrix-valued function $A:[0,1]^d \to \R^{d\times d}$, we denote $\| A \|_{\mathrm{op},\infty} \coloneqq \sup_{x\in [0,1]^d} \| A(x) \|_{\mathrm{op}}$.

We start off by fixing $\mu\coloneqq\mathrm{Unif}([0,1]^d)$ throughout the entire proof, so that $d\mu(x)=dx$ on $[0,1]^d$. All change-of-variables and integration-by-parts identities below are stated relative to the Lebesgue measure.

\smallskip
\noindent\emph{Nonparametric lower bound for $T_0$ via balanced optimal transport.}

The class $\mathcal U_\alpha$ contains a balanced OT subclass in which $\lambda_0\equiv 1$, $M_\mu=M_\nu=1$, and the Monge problem \eqref{eq:gh_monge} reduces to the balanced OT problem on $[0,1]^d$. The minimax lower bound $n^{-2\alpha/(2\alpha-2+d)}$ for smooth transport map estimation in balanced OT was established in \cite{hutter2021minimax}. Therefore,
\begin{align}
    \label{eq:lb-T-nonpar}
    \inf_{\hat T} \mathfrak M_n^T(\hat T)  \gtrsim  n^{-2\alpha/(2\alpha-2+d)}.
\end{align}

\smallskip
\noindent\emph{Nonparametric lower bound for $\lambda_0$ by reduction to balanced optimal transport.}

We show that estimating the growth map $\lambda_0$ is at least as hard as estimating a balanced OT transport map of smoothness $\alpha$, yielding the same nonparametric rate. The key observation is that $\lambda_0=\exp(-z+\frac1{4}\|\nabla z\|^2)$ inherits $\alpha$-smoothness from $\|\nabla z\|^2$.

Let $\xi\in C^\infty(\R)$ be a non-zero bump with $\supp(\xi)\subset[0,1]$, $\xi(0)=\xi(1)=0$, and $\xi(x_*)\ne 0$, $\xi'(x_*)\ne 0$ for some $x_*\in(0,1)$; define
\begin{align}
    g(x)=\prod_{i=1}^d\xi(x_i), \qquad x=(x_1,\dots,x_d).
\end{align}

Fix a small constant $a>0$ and a smooth cutoff $\chi\in C^\infty_c((0,1)^d)$ with $\chi\equiv 1$ on $[\delta_0,1-\delta_0]^d$ for some $\delta_0\in(0,1/4)$. Define
\begin{align}
    \label{eq:base-potential}
    \varphi_0(x) \coloneqq a \chi(x) x_1.
\end{align}
Since $\varphi_0\in C^\infty_c((0,1)^d)$, the induced pair $(T_0,\lambda_0)$ via \eqref{eq:convenient} satisfies $T_0=\mathrm{id}$ and $\lambda_0=1$ on a neighborhood of $\partial[0,1]^d$. For $a$ small enough, the function $\Psi_0(x) \coloneqq \frac1{2}\|x\|^2-\varphi_0$ remains strictly convex (since $\|\nabla^2 \varphi_0\|_{\mathrm{op},\infty}\le Ca$), and $T_0:[0,1]^d\to[0,1]^d$ is a diffeomorphism. This construction yields $\nabla \varphi_0(x)=(a,0,\dots,0)$ for all $x \in [\delta_0,1-\delta_0]^d$.

Let $J=\lceil\theta  n^{1/(2\alpha-2+d)}\rceil$ for a constant $\theta>0$ to be determined later, and place a grid $\{x^{(j)}\}_{j\in\mJ}$ with $x^{(j)}_i=(j_i-1)/J$, with $\mJ\subset[J]^d$, satisfying $[x^{(j)},x^{(j)}+1/J]^d\subset[\delta_0,1-\delta_0]^d$; the cardinality $N_J\coloneqq|\mJ|\asymp J^d$ satisfies $c_{l} J^d\le N_J\le c_{u} J^d$ for some constants $c_{l},c_{u}>0$. Define
\begin{align}
    g_j(x) \coloneqq \frac{\varepsilon_{\mathrm b}}{J^{\alpha+1}} g(J(x-x^{(j)})),
\end{align}
so that $\{\supp (g_j)\}_{j\in\mJ}$ are pairwise disjoint and contained in $[\delta_0,1-\delta_0]^d$. Since $\partial^\flat g_j(\cdot)=\varepsilon_{\mathrm b}  J^{|\flat|-\alpha-1}\partial^\flat g(J(\cdot-x^{(j)}))$ for any multi-index $\flat$, and $\alpha>1$, one may fix $\varepsilon_{\mathrm b}$ small enough and $J\ge J_0$ large enough that uniformly in $j,J$:
\begin{align}
    \label{eq:bump-bounds}
    \|g_j\|_\infty\le J^{-\alpha-1}, \quad \|\nabla g_j\|_\infty\le J^{-\alpha}, \quad \|\nabla^2 g_j\|_{\mathrm{op},\infty}\le 1/2, \quad \|g_j\|_{C^{\alpha+1}}\le C_g,
\end{align}
for some constant $C_g>0$. Applying External result~\ref{lem:vg} with $N_J\ge 8$, we obtain binary strings $\omega^{(0)},\dots,\omega^{(K)}\in\{0,1\}^{\mJ}$ with $\omega^{(0)}=\mathbf{0}$, $K\ge 2^{N_J/8}$, and $\|\omega^{(k)}-\omega^{(k')}\|_2^2\ge N_J/8$ for $k\ne k'$. Define
\begin{align}
    \varphi_k(x) &\coloneqq \varphi_0(x) + \sum_{j\in\mJ}\omega^{(k)}_j g_j(x),\qquad k=0,\dots,K,\\
    T_k(x) &\coloneqq x - \nabla \varphi_k(x), \quad \lambda_k(x)\coloneqq\exp \Big( -\varphi_k(x)+\frac1{4}\|\nabla \varphi_k(x)\|^2\Big),\\
    \nu_k &\coloneqq (T_k)_\#(\lambda_k^2\mu).
\end{align}
Write $\Psi_k(x)\coloneqq\frac1{2}\|x\|^2-\varphi_k(x)$, so that $T_k=\nabla\Psi_k$. Since $\nabla^2 \varphi_k = \nabla^2 \varphi_0 + \nabla^2(\varphi_k-\varphi_0)$ and $\|\nabla^2 \varphi_0\|_{\mathrm{op},\infty}\le Ca$, $\|\nabla^2(\varphi_k-\varphi_0)\|_{\mathrm{op},\infty}\le 1/2$ by \eqref{eq:bump-bounds}, we have $\nabla^2\Psi_k\succeq(1/2-Ca)I_d\succ 0$ for $a$ small enough, so $\Psi_k$ is strictly convex and $\varphi_k$ is $c$-concave. It then follows from Theorem~\ref{prop:monge-potential} that $(T_k,\lambda_k)$ is the unique solution of the Monge problem \eqref{eq:gh_monge} associated with $\varphi_k$ for the pair $(\mu,\nu_k)$. Consequently, $(T_k,\lambda_k)$ is uniquely associated with $(\mu,\nu_k)$.

Moreover, $\varphi_k\in C^{\alpha+1}(\Omega)$ with $\|\varphi_k\|_{C^{\alpha+1}(\Omega)}\le\|\varphi_0\|_{C^{\alpha+1}}+C_g$, so $T_k\in C^\alpha(\Omega;\Omega)$ and $\lambda_k\in C^\alpha(\Omega)$ with uniform bounds. Because $\varphi_0\in C^\infty_c((0,1)^d)$ and $w_k\coloneqq \varphi_k-\varphi_0$ is supported strictly inside $[\delta_0,1-\delta_0]^d$, we have $\nabla \varphi_k\equiv 0$ on a neighborhood of $\partial[0,1]^d$, hence $T_k\equiv\mathrm{id}$ there and $T_k([0,1]^d)\subset[0,1]^d$. 


Because $\mu=\mathrm{Unif}([0,1]^d)$ and each perturbation is supported strictly inside $[0,1]^d$, the induced target densities remain uniformly positive, uniformly bounded, and equal to $1$ near the boundary. Together with the uniform $C^\alpha$ bounds on $T_k$ and $\lambda_k$, this shows that, after fixing $B$ and $\Lambda$ large enough, each pair $(\mu,\nu_k)$ belongs to $\mathcal U_\alpha$. In particular, the lower-bound construction lives in the same regime as the upper bound for the kernel-based method (Theorem~\ref{thm:main_cube_kernel_rates}).

To utilize the multiple hypotheses in External result~\ref{thm:tsy-fano}, we first calculate the $L^2(\mu)$-separation between $\lambda_k$'s. Denote $\ell_0 \coloneqq \log \lambda_0 = -\varphi_0+\frac1{4}\|\nabla \varphi_0\|^2$ and $\ell_k \coloneqq \log \lambda_k = -\varphi_k+\frac1{4}\|\nabla \varphi_k\|^2$. Since $\nabla \varphi_0(x)=(a,0,\dots,0)$ for all $x\in[\delta_0,1-\delta_0]^d$, we compute
\begin{align}
    \ell_k - \ell_0 &= \bigl(-\varphi_k+\frac1{4}\|\nabla \varphi_k\|^2\bigr) - \bigl(-\varphi_0+\frac1{4}\|\nabla \varphi_0\|^2\bigr) \notag\\
    &= -(\varphi_k-\varphi_0) + \frac1{2}\nabla \varphi_0\cdot\nabla(\varphi_k-\varphi_0) + \frac1{4}\|\nabla(\varphi_k-\varphi_0)\|^2 \\
    &= \frac{a}{2} \partial_1(\varphi_k-\varphi_0) - (\varphi_k-\varphi_0) + \frac1{4}\|\nabla(\varphi_k-\varphi_0)\|^2. \label{eq:u-expand}
\end{align}
The three terms have pointwise sizes $O(a J^{-\alpha})$, $O(J^{-\alpha-1})$, and $O(J^{-2\alpha})$; since $a$ is a fixed constant, the cross term ${ (a/2)}\partial_1(\varphi_k-\varphi_0)$ dominates for $J$ large. By $\lambda_k-\lambda_0=\lambda_0(e^{\ell_k-\ell_0}-1)$ and $|e^u-1-u|\le u^2 e^{|u|}$,
\begin{align}
    \lambda_k - \lambda_0 = \lambda_0\Big(\frac{a}{2} \partial_1(\varphi_k-\varphi_0) + R_k\Big), \label{eq:lambda-expand}
\end{align}
where $|R_k|\le |(\varphi_k-\varphi_0)|+C\|\nabla(\varphi_k-\varphi_0)\|^2+C'(\ell_k-\ell_0)^2\le C''J^{-\alpha-1}$ for $J\ge J_0$.

We now consider $\lambda_k - \lambda_{k'}$ for $k\ne k'$. Setting $\lambda_{\min}\coloneqq\inf_{x \in [0,1]^d}\lambda_0(x) >0$, we use the leading $\partial_1$-cross term in~\eqref{eq:lambda-expand} together with the disjointness of $\{\supp(\nabla g_j)\}_{j\in\mJ}$. We claim
\begin{align}
    \label{eq:cg-positive}
    c_g \coloneqq \int_{[0,1]^d}|\partial_1 g(y)|^2 dy  >  0.
\end{align}
Indeed, $\partial_1 g$ is the partial derivative of the bump $g(x)=\prod_i\xi(x_i)$ in the first coordinate. If $\partial_1 g\equiv 0$ on $[0,1]^d$, then $g$ would be constant in $x_1$; combined with $\supp(g)\subset[0,1]^d$ and $\xi(0)=\xi(1)=0$, this would force $g(x_*)=0$, contradicting $\xi'(x_*)\ne 0$. Hence $\partial_1 g$ is non-zero on a set of positive Lebesgue measure, giving~\eqref{eq:cg-positive}. By the rescaling $\partial_1 g_j(x)=\varepsilon_{\mathrm b} J^{-\alpha}(\partial_1 g)(J(x-x^{(j)}))$ and the substitution $y=J(x-x^{(j)})$,
\begin{align}
    \int|\partial_1 g_j|^2 d\mu = \varepsilon_{\mathrm b}^2 c_g J^{-2\alpha-d}.
\end{align}
By \eqref{eq:lambda-expand}, the disjoint supports of $\{\partial_1 g_j\}_{j\in\mJ}$, and $\|\omega^{(k)}-\omega^{(k')}\|_2^2\ge N_J/8\ge c_l J^d/8$, applying the inequality $\|A+B\|^2\ge\frac1{2}\|A\|^2-\|B\|^2$ with $A=\lambda_0(a/2)\partial_1(\varphi_k-\varphi_{k'})$ and $B=\lambda_0(R_k-R_{k'})$ yields
\begin{align}
    \label{eq:lambda-sep}
    \|\lambda_k-\lambda_{k'}\|_{L^2(\mu)}^2
    &\ge \frac{a^2\lambda_{\min}^2}{8}\sum_{j\in\mJ}(\omega_j^{(k)}-\omega_j^{(k')})^2 \int|\partial_1 g_j|^2 d\mu - \lambda_{\max}^2 \int(R_k-R_{k'})^2 d\mu \notag\\
    &\ge \frac{a^2\lambda_{\min}^2}{8}\cdot\frac{c_l J^d}{8}\cdot\varepsilon_{\mathrm b}^2 c_g J^{-2\alpha-d} - C J^{-2\alpha-2} \notag\\
    &\ge c_\lambda J^{-2\alpha},
\end{align}
for $J$ large enough that the remainder $C J^{-2\alpha-2}$ is absorbed; the constant $c_\lambda\coloneqq(1/128) a^2\lambda_{\min}^2\varepsilon_{\mathrm b}^2 c_g c_l>0$ is independent of $J,k,k'$.

Next, we bound the KL divergence $D_\KL(\bar\nu_k\|\bar\nu_0)$ between the normalized hypotheses $\bar\nu_k\coloneqq\nu_k/\nu_k(\Omega)$. Throughout, we denote $h_k\coloneqq d\nu_k/d\mu$.

Write $M_k\coloneqq\nu_k(\Omega)=\int h_k d\mu$, so that the density of $\bar\nu_k$ with respect to $\mu$ is $\bar h_k\coloneqq h_k/M_k$. Since $\nu_k=(T_k)_\#(\lambda_k^2\mu)$, $M_k=\int\lambda_k^2 d\mu$, and because $\|\varphi_k\|_\infty,\|\nabla \varphi_k\|_\infty$ are bounded uniformly in $k$, the weight $\lambda_k=\exp(-\varphi_k+\frac1{4}\|\nabla \varphi_k\|^2)$ satisfies $0<\lambda_{\min}\le\lambda_k\le\lambda_{\max}<\infty$ uniformly in $k$. Combined with $\mu(\Omega)=1$, this yields
\begin{align}
    \lambda_{\min}^2 \le  M_k \le \lambda_{\max}^2\qquad\text{for all }k=0,\dots,K,
\end{align}
We bound the KL-divergence by the $\chi^2$-divergence in two steps. First, the Cauchy--Schwarz inequality applied to $(M_k-M_0)=\int(h_k-h_0)d\mu=\int\frac{h_k-h_0}{\sqrt{h_0}}\cdot\sqrt{h_0} d\mu$ gives
\begin{align}
    \label{eq:CS-mass}
    (M_k-M_0)^2 \le \Bigl(\int \frac{(h_k-h_0)^2}{h_0} d\mu\Bigr)\Bigl(\int h_0 d\mu\Bigr) = M_0 \int \frac{(h_k-h_0)^2}{h_0} d\mu.
\end{align}
Splitting $\bar h_k-\bar h_0=\frac{h_k-h_0}{M_k}-\frac{h_0(M_k-M_0)}{M_0 M_k}$, we have
\begin{align}
    D_\KL(\bar\nu_k\|\bar\nu_0)
    \le \int \frac{(\bar h_k-\bar h_0)^2}{\bar h_0} d\mu
    &\le \frac{2M_0}{M_k^2} \int \frac{(h_k-h_0)^2}{h_0} d\mu + \frac{2(M_k-M_0)^2}{M_k^2 M_0} \int h_0 d\mu \notag\\
    &= \frac{2M_0}{M_k^2} \int \frac{(h_k-h_0)^2}{h_0} d\mu + \frac{2(M_k-M_0)^2}{M_k^2} \notag\\
    &\stackrel{\eqref{eq:CS-mass}}{\le}\frac{4M_0}{M_k^2} \int \frac{(h_k-h_0)^2}{h_0} d\mu  \le  C \int \frac{(h_k-h_0)^2}{h_0} d\mu,
\end{align}
where the last inequality uses $\lambda_{\min}^2\le M_k\le\lambda_{\max}^2$ to bound $4M_0/M_k^2\le 4\lambda_{\max}^2/\lambda_{\min}^4=\colon C$, a constant depending only on $\lambda_{\min}$ and $\lambda_{\max}$. To bound the integral, we consider $r_k(y)\coloneqq h_k(y)/h_0(y)$. If $\|\log r_k\|_\infty\le 1/2$ for a sufficiently large $J$, the inequality $|e^u-1-u|\le u^2$ for $|u|\le 1/2$ applied to $u=\log r_k(y)$ and multiplied by $h_0(y)$ gives:
\begin{align}
    \label{eq:hk-expand}
    h_k(y)-h_0(y)  =  h_0(y) \log r_k(y) + R_k(y),
    \qquad |R_k(y)|\le h_0(y) (\log r_k(y))^2.
\end{align}
Consequently,
\begin{align}
    \label{eq:I-reduction}
    \int \frac{(h_k-h_0)^2}{h_0} d\mu  \le  \|h_0\|_\infty \int(\log r_k)^2 d\mu.
\end{align}
It thus suffices to bound $\|\log r_k\|_\infty$. By the change-of-variables formula for $\nu_k=(T_k)_\#(\lambda_k^2\mu)$, $\lambda_k=\exp(-\varphi_k+\frac1{4}\|\nabla \varphi_k\|^2)$, and $\nabla T_k=I-\nabla^2\varphi_k$,
\begin{align}
    \log h_k(y)
    & =  \Bigl[ 2\log\lambda_k(x)-\log\det \nabla T_k(x) \Bigr]\Big\vert_{x=T_k^{-1}(y)} \\
    & =  \Bigl[ \underbrace{-2\varphi_k(x)+\frac1{2}\|\nabla \varphi_k(x)\|^2-\log\det(I-\nabla^2\varphi_k(x))}_{L_k(x)} \Bigr]\Big\vert_{x=T_k^{-1}(y)}.
\end{align}
Thus $\log r_k(y)=L_k(T_k^{-1}(y))-L_0(T_0^{-1}(y))$, which we split as
\begin{align}
    \label{eq:logr-split}
    \log r_k(y)  =  (L_k-L_0)(T_k^{-1}(y)) + \bigl[L_0(T_k^{-1}(y))-L_0(T_0^{-1}(y))\bigr].
\end{align}
We first bound $L_k(x) - L_0(x)$. By definition,
\begin{align}
    L_k(x)-L_0(x)
    &= -2(\varphi_k-\varphi_0)(x) + \frac1{2}\bigl(\|\nabla \varphi_k(x)\|^2-\|\nabla \varphi_0(x)\|^2\bigr) \notag\\
    &\quad + \bigl[-\log\det(I-\nabla^2 \varphi_k(x))+\log\det(I-\nabla^2 \varphi_0(x))\bigr].
\end{align}
We bound each term separately.
\begin{enumerate}
\item[(i)] \emph{Zeroth-order term.} From the bounds following \eqref{eq:u-expand}, $\varphi_k-\varphi_0=O_\infty(J^{-\alpha-1})$.

\item[(ii)] \emph{Gradient-squared term.} Since $\nabla \varphi_0=(a,0,\dots,0)$ on $[\delta_0,1-\delta_0]^d$,
\begin{align}
    \frac1{2}\|\nabla \varphi_k\|^2-\frac1{2}\|\nabla \varphi_0\|^2
    = a \partial_1(\varphi_k-\varphi_0)+\frac1{2}\|\nabla(\varphi_k-\varphi_0)\|^2
    = O_\infty(a J^{-\alpha}) + O_\infty(J^{-2\alpha}).
\end{align}

\item[(iii)] \emph{Log-determinant term.} We claim that pointwise on $\Omega$,
\begin{align}
    \label{eq:logdet-collapse}
    -\log\det(I-\nabla^2 \varphi_k(x)) + \log\det(I-\nabla^2 \varphi_0(x))  =  -\log\det(I-\nabla^2(\varphi_k-\varphi_0)(x)).
\end{align}
To see this, split $\Omega$ into the bump supports and their complement. On each bump support $\mathrm{supp}(g_j)\subset[\delta_0,1-\delta_0]^d$, $\chi\equiv 1$, so $\varphi_0(x)=ax_1$ is affine and $\nabla^2 \varphi_0(x)\equiv 0$; hence $\nabla^2 \varphi_k=\nabla^2(\varphi_k-\varphi_0)$ and $\log\det(I-\nabla^2 \varphi_0)=0$, which gives \eqref{eq:logdet-collapse}. Off the bump supports, $\varphi_k-\varphi_0=\sum_j\omega_j^{(k)}g_j=0$, so $\nabla^2 \varphi_k=\nabla^2 \varphi_0$ and both sides of \eqref{eq:logdet-collapse} vanish. The claim thus follows.

Applying the Taylor expansion $\log\det(I-A)=-\mathrm{tr}(A)+O(\|A\|_{\mathrm{op}}^2)$ with $A=\nabla^2(\varphi_k-\varphi_0)$, which is valid since $\|\nabla^2(\varphi_k-\varphi_0)\|_{\mathrm{op},\infty}\le 1/2$ by \eqref{eq:bump-bounds}, we have
\begin{align}
    -\log\det(I-\nabla^2(\varphi_k-\varphi_0))  =  \Delta(\varphi_k-\varphi_0) + O_\infty(J^{2-2\alpha}),
\end{align}
where $O_\infty(J^{2-2\alpha})$ follows from $\|\nabla^2(\varphi_k-\varphi_0)\|_{\mathrm{op},\infty}=O(J^{1-\alpha})$.
\end{enumerate}
Combining (i)--(iii) with $\varphi_k-\varphi_0=O_\infty(J^{-\alpha-1})$, and noting that for $\alpha\ge 1$ and $J\ge 1/a$ the bound $J^{1-\alpha}\ge a J^{-\alpha}\ge J^{-\alpha-1}\ge J^{-2\alpha}$ lets us absorb lower-order terms,
\begin{align}
    \label{eq:deltaL-expand}
    L_k(x)-L_0(x)  =  \Delta \varphi_k(x)-\Delta \varphi_0(x)+O_\infty(a J^{-\alpha})+O_\infty(J^{2-2\alpha}).
\end{align}
For the second term in \eqref{eq:logr-split}, we have $L_0(x)=-2\varphi_0(x)+\frac1{2}a^2=-2ax_1+\mathrm{const}$ on $[\delta_0, 1-\delta_0]^d$, so $\nabla L_0\equiv-2ae_1$ and:
\begin{align}\label{eq:L0-expand}
    L_0(T_k^{-1}(y))-L_0(T_0^{-1}(y))  =  -2a e_1 \cdot \bigl(T_k^{-1}(y)-T_0^{-1}(y)\bigr),
\end{align}
To bound the difference, we fix $y\in\Omega$ and set $x_k\coloneqq T_k^{-1}(y)$, $x_0\coloneqq T_0^{-1}(y)$. From $y=T_k(x_k)=x_k-\nabla \varphi_k(x_k)$ and $y=T_0(x_0)=x_0-\nabla \varphi_0(x_0)$, we have
\begin{align}
    (x_k-x_0)  -  \bigl(\nabla \varphi_0(x_k)-\nabla \varphi_0(x_0)\bigr)
     =  \nabla \varphi_k(x_k)-\nabla \varphi_0(x_k).
\end{align}
By the fundamental theorem of calculus, $\nabla \varphi_0(x_k)-\nabla \varphi_0(x_0)=A_y(x_k-x_0)$ with $A_y\coloneqq\int_0^1 \nabla^2 \varphi_0(x_0+t(x_k-x_0)) dt$ and $\|A_y\|_{\mathrm{op}}\le\|\nabla^2 \varphi_0\|_{\mathrm{op},\infty}\le Ca$. Hence
\begin{align}
    (I-A_y)(x_k-x_0) = \nabla \varphi_k(x_k)-\nabla \varphi_0(x_k).
\end{align}
For a sufficiently small $a>0$, $I-A_y$ is invertible with $\|(I-A_y)^{-1}\|_{\mathrm{op}}\le(1-Ca)^{-1}$. Consequently, using $\|\nabla(\varphi_k-\varphi_0)\|_\infty\le J^{-\alpha}$ from \eqref{eq:bump-bounds}, we obtain:
\begin{align}\label{eq:T-inv-diff}
    \|T_k^{-1}(y)-T_0^{-1}(y)\|  \le  \frac{\|\nabla(\varphi_k-\varphi_0)\|_\infty}{1-Ca}  =  O(J^{-\alpha}),
\end{align}
Substituting \eqref{eq:deltaL-expand}, \eqref{eq:L0-expand} and \eqref{eq:T-inv-diff} into~\eqref{eq:logr-split} yields:
\begin{align}
    \log r_k(y)  =  \Delta \varphi_k(T_k^{-1}(y))-\Delta \varphi_0(T_k^{-1}(y))+O_\infty(a J^{-\alpha})+O_\infty(J^{2-2\alpha}).
\end{align}
In particular $\|\log r_k\|_\infty=O(J^{1-\alpha})\le 1/2$ for a sufficiently large $J$, justifying the hypothesis used in~\eqref{eq:hk-expand}.

Plugging the bound for $\log r_k$ into~\eqref{eq:I-reduction}, yields:
\begin{align}\label{eq:KL-bound-master}
D_\KL(\bar\nu_k\|\bar\nu_0)\le C \int(\Delta \varphi_k\circ T_k^{-1}-\Delta \varphi_0\circ T_k^{-1})^2 d\mu + O(a^2 J^{-2\alpha}) + O(J^{4-4\alpha}).
\end{align}
We now bound the leading term $\int(\Delta(\varphi_k-\varphi_0)\circ T_k^{-1})^2 d\mu$. The change of variables $x=T_k^{-1}(y)$ has Jacobian $\det \nabla T_k(x)=\det(I-\nabla^2(\varphi_k-\varphi_0))=1+O(J^{1-\alpha})$, so $\|\Delta(\varphi_k-\varphi_0)\circ T_k^{-1}\|_{L^2(\mu)}^2\le C\|\Delta(\varphi_k-\varphi_0)\|_{L^2(\mu)}^2$. Using disjoint supports of $\{g_j\}_{j\in\mJ}$, $(\omega_j^{(k)})^2=\omega_j^{(k)}\in\{0,1\}$, and the rescaling $\Delta g_j(x)=\varepsilon_{\mathrm b} J^{1-\alpha}(\Delta g)(J(x-x^{(j)}))$,
\begin{align}
    \int(\Delta(\varphi_k-\varphi_0))^2 d\mu
     = \sum_{j\in\mJ}\omega_j^{(k)}\cdot\frac{\varepsilon_{\mathrm b}^2}{J^{2\alpha-2+d}} \int(\Delta g)^2 dy
     \le  C \varepsilon_{\mathrm b}^2 J^{2-2\alpha},
\end{align}
using $\sum_{j\in\mJ}\omega_j^{(k)}\le N_J\le c_u J^d$. The two remainders in~\eqref{eq:KL-bound-master} contribute $O(a^2 J^{-2\alpha})$ and $O(J^{4-4\alpha})$, both dominated by $J^{2-2\alpha}$ for $\alpha\ge 1$. Hence,
\begin{align}
    D_\KL(\bar\nu_k\|\bar\nu_0)  \le  C J^{2-2\alpha}.
\end{align}
By \eqref{eq:iid_kl} and $\mu$ fixed across hypotheses,
\begin{align}
    D_\KL(\mathbb P_{\mu,\nu_k}^n\|\mathbb P_{\mu,\nu_0}^n)
    = n D_\KL(\bar\nu_k\|\bar\nu_0)
    \le C_0 n J^{2-2\alpha}.
\end{align}
From External result~\ref{lem:vg} and $N_J\ge c_l J^d$, $\log K\ge (N_J/8)\log 2\ge (c_l J^d\log 2)/8$, so
\begin{align}
    (\log K)/9 \ge (c_l J^d\log 2)/72.
\end{align}
Using $J\ge\theta n^{1/(2\alpha-2+d)}$ from the definition $J=\lceil\theta n^{1/(2\alpha-2+d)}\rceil$,
\begin{align}
    C_0 n J^{2-2\alpha} \le (\log K)/9
    \quad\Longleftrightarrow\quad
    \theta^{2\alpha-2+d} \ge 72 C_0/(c_l\log 2).
\end{align}
Therefore, by choosing
\begin{align}
    \theta  >  \left(72 C_0/(c_l\log 2)\right)^{ 1/(2\alpha-2+d)},
\end{align}
we have $D_\KL(\mathbb P_{\mu,\nu_k}^n\|\mathbb P_{\mu,\nu_0}^n)\le\log K/9$, as required by External result~\ref{thm:tsy-fano}.

Applying External result~\ref{thm:tsy-fano} with $d_\lambda(\lambda,\lambda')^2\coloneqq\|\lambda-\lambda'\|_{L^2(\mu)}^2$ and separation $\gtrsim J^{-2\alpha}$ from~\eqref{eq:lambda-sep} yields
\begin{align}
    \label{eq:lb-lambda-nonpar}
    \inf_{\hat\lambda} \mathfrak M_n^\lambda(\hat\lambda)  \gtrsim  J^{-2\alpha}  \asymp  n^{-2\alpha/(2\alpha-2+d)}.
\end{align}

\smallskip
\noindent\emph{Parametric lower bound from two hypotheses.}

We show $\inf_{\hat T}\mathfrak M_n^T(\hat T)\gtrsim n^{-1}$ and $\inf_{\hat\lambda}\mathfrak M_n^\lambda(\hat\lambda)\gtrsim n^{-1}$ simultaneously. Retain $\varphi_0$ as in \eqref{eq:base-potential} and $\mu=\mathrm{Unif}([0,1]^d)$. For $\tilde\theta>0$ small, define
\begin{align}
    \varphi_1(x)\coloneqq \varphi_0(x)+(\tilde\theta/\sqrt n) g(x),
\end{align}
and let $(T_i,\lambda_i,\nu_i)$ be the corresponding pairs and targets via \eqref{eq:convenient}, $i=0,1$. For $\tilde\theta$ small and $n$ large, $\Psi_1=\frac1{2}\|\cdot\|^2-\varphi_1$ remains strictly convex, so $\varphi_1$ is $c$-concave. It then follows from Theorem~\ref{prop:monge-potential} that $(T_1,\lambda_1)$ is the unique solution of \eqref{eq:gh_monge} associated with $\varphi_1$ for the pair $(\mu,\nu_1)$.

For the separation between the two transport maps, we have
\begin{align}
    \|T_0-T_1\|_{L^2(\mu)}^2
    = (\tilde\theta^2/n) \int \|\nabla g\|^2 d\mu  \gtrsim  1/n.
\end{align}
For the growth maps, since $\nabla \varphi_0=(a,0,\dots,0)$ on $\supp(g)\subset[\delta_0,1-\delta_0]^d$, the same expansion as in~\eqref{eq:lambda-expand} with $w=(\tilde\theta/\sqrt n) g$ gives
\begin{align}
    \lambda_1(x)-\lambda_0(x) = \lambda_0(x) \frac{\tilde\theta}{\sqrt n}\bigl[(a/2)\partial_1 g(x) - g(x)\bigr] + R(x),
\end{align}
with $\|R\|_\infty=O(n^{-1})$. Squaring and integrating against $d\mu=dx$,
\begin{align}
    \|\lambda_0-\lambda_1\|_{L^2(\mu)}^2 \ge \lambda_{\min}^2 \frac{\tilde\theta^2}{n} \int_{[0,1]^d} \bigl[(a/2)\partial_1 g - g\bigr]^2 dx - O(n^{-2}).
\end{align}
We claim
\begin{align}
    \label{eq:cgpar-positive}
    c_g^{\mathrm{par}} \coloneqq \int_{[0,1]^d} \bigl[(a/2)\partial_1 g(y) - g(y)\bigr]^2 dy  >  0.
\end{align}
Indeed, suppose for contradiction $(a/2)\partial_1 g - g\equiv 0$ on $[0,1]^d$. Along any line $t\mapsto (t,x_2,\dots,x_d)$ with fixed $(x_2,\dots,x_d)\in[0,1]^{d-1}$, the function $\varphi(t)\coloneqq g(t,x_2,\dots,x_d)$ would satisfy the linear ODE $(a/2)\varphi'(t)=\varphi(t)$, so $\varphi(t)=\varphi(0)e^{2t/a}$ wherever the line lies in $[0,1]$. Since $\varphi$ vanishes outside $[0,1]$ (as $\supp(g)\subset[0,1]^d$ and $\xi(0)=\xi(1)=0$), continuity forces $\varphi(0)=0$, hence $\varphi\equiv 0$. Applying this to every such line and to $x_*$ gives $g(x_*)=0$, contradicting $g(x_*)\ne 0$. Therefore $(a/2)\partial_1 g - g$ is non-zero on a set of positive Lebesgue measure, giving~\eqref{eq:cgpar-positive}. Hence
\begin{align}
    \|\lambda_0-\lambda_1\|_{L^2(\mu)}^2 \ge \lambda_{\min}^2 \frac{\tilde\theta^2}{n} c_g^{\mathrm{par}} - O(n^{-2})  \gtrsim  1/n.
\end{align}
We now consider the KL divergence. The expansion \eqref{eq:KL-bound-master} with $J=1$ gives $D_\KL(\bar\nu_1\|\bar\nu_0)\lesssim \tilde\theta^2/n$, so $n D_\KL(\bar\nu_1\|\bar\nu_0)\lesssim\tilde\theta^2$. By External result~\ref{thm:tsy-twopoint},
\begin{align}
    \label{eq:lb-T-par}
    \inf_{\hat T} \mathfrak M_n^T(\hat T) & \gtrsim  1/n,\\
    \label{eq:lb-lambda-par}
    \inf_{\hat\lambda} \mathfrak M_n^\lambda(\hat\lambda) & \gtrsim  1/n.
\end{align}

Combining \eqref{eq:lb-T-nonpar}, \eqref{eq:lb-lambda-nonpar}, \eqref{eq:lb-T-par}, and \eqref{eq:lb-lambda-par} yields Theorem~\ref{thm:lb-ubot}.
\end{proof}




\section{Wasserstein convergence of the weighted empirical measure} \label{sec:empirical_wasserstein}


We record here that the Wasserstein convergence of $\tilde\mu_n$ is an immediate reduction to the probability-measure case. In particular, the idea of multiplying the transport plan in the proof of \cite{weed2019sharp} by $M_\mu$ is correct. The only point to keep in mind is that the factor $M_\mu$ appears at the level of $W_p^p$, while $W_p$ itself scales like $M_\mu^{1/p}$.

\begin{definition}
Let $p\in[1,\infty)$ and let $\alpha,\beta\in \mM_+(\Omega)$ satisfy $\alpha(\Omega)=\beta(\Omega)<\infty$.
We define the $p$-Wasserstein distance between $\alpha$ and $\beta$ by
\begin{align}
    W_p(\alpha,\beta) \coloneqq \left( \inf_{\pi\in\Pi(\alpha,\beta)} \int_{\Omega\times\Omega} \|x-y\|^p   d\pi(x,y) \right)^{1/p},
\end{align}
where $\Pi(\alpha,\beta)$ denotes the set of couplings of $\alpha$ and $\beta$.
When $\alpha$ and $\beta$ are probability measures, this is the usual Wasserstein distance.
\end{definition}

\begin{lemma}[Scaling in the total mass] \label{lem:mass_scaling_wasserstein}
Let $m>0$ and let $\alpha,\beta$ be probability measures on $\Omega$.
Then
\begin{align}
    W_p^p(m\alpha,m\beta) = m   W_p^p(\alpha,\beta),
    \qquad
    W_p(m\alpha,m\beta) = m^{1/p} W_p(\alpha,\beta).
\end{align}
\end{lemma}

\begin{proof}
The map $\pi \mapsto m\pi$ is a bijection from $\Pi(\alpha,\beta)$ onto $\Pi(m\alpha,m\beta)$.
Therefore,
\begin{align}
    \inf_{\gamma\in\Pi(m\alpha,m\beta)} \int \|x-y\|^p   d\gamma(x,y)
    = m \inf_{\pi\in\Pi(\alpha,\beta)} \int \|x-y\|^p   d\pi(x,y).
\end{align}
Taking the $p$th root yields the second identity.
\end{proof}

We introduce main ingredients for bounding the convergence rates of empirical measures in Wasserstein distances.

For a probability measure $\eta$ on $\Omega$ and $\tau\in[0,1)$, let $N_\epsilon(S)$ be the minimum number of closed balls of diameter $\epsilon$ needed to cover $S$, and define:
\begin{align}
    N_\epsilon(\eta,\tau) \coloneqq \inf\{N_\epsilon(S): S\subset \Omega,\ \eta(S)\ge 1-\tau\},
    \qquad
    d_\epsilon(\eta,\tau) \coloneqq \frac{\log N_\epsilon(\eta,\tau)}{-\log \epsilon}.
\end{align}
We also set
\begin{align}
    d_p^*(\eta)
    &\coloneqq \inf\left\{ s>2p : \limsup_{\epsilon\downarrow 0} d_\epsilon\left(\eta,\epsilon^{sp/(s-2p)}\right) \le s \right\}, \\
    d_*(\eta)
    &\coloneqq \lim_{\tau\downarrow 0} \liminf_{\epsilon\downarrow 0} d_\epsilon(\eta,\tau).
\end{align}
These are the upper and lower Wasserstein dimensions introduced by \cite{weed2019sharp}.


\begin{proposition}[Weighted empirical measure: convergence and rates]
\label{prop:weighted_empirical_wasserstein}
Let $M_\mu\in(0,\infty)$.
Let $X_1,X_2,\dots \stackrel{\mathrm{i.i.d.}}{\sim} \bar\mu$, and define
\begin{align}
    \bar\mu_n \coloneqq {1/n}\sum_{i=1}^n \delta_{X_i},
    \qquad
    \tilde\mu_n \coloneqq  {M_\mu/n}\sum_{i=1}^n \delta_{X_i}.
\end{align}
Then, the following hold:
\begin{enumerate}[label=(\roman*)]
    \item For every $n\in\N$,
    \begin{align}
        W_p(\mu,\tilde\mu_n) = (M_\mu)^{1/p} W_p(\bar\mu,\bar\mu_n).
    \end{align}
    In particular, $W_p(\mu,\tilde\mu_n)\to 0$ almost surely as $n\to\infty$.

    \item If $s>d_p^*(\bar\mu)$, then
    \begin{align}
    \label{eq:weighted_empirical_wp_power_rate}
        \Ep[W_p^p(\mu,\tilde\mu_n)] \lesssim M_\mu n^{-p/s}.
    \end{align}
    Consequently,
    \begin{align}
        \Ep[W_p(\mu,\tilde\mu_n)] \lesssim (M_\mu)^{1/p} n^{-1/s}.
    \end{align}

    \item If $t<d_*(\bar\mu)$ and $\nu_n\in\mM_+(\Omega)$ is any discrete measure with
    $\nu_n(\Omega)=M_\mu$ and $|\mathrm{supp}(\nu_n)|\le n$, then
    \begin{align}
        W_p(\mu,\nu_n) \gtrsim (M_\mu)^{1/p} n^{-1/t}.
    \end{align}
    In particular, the same lower bound holds for $\nu_n=\tilde\mu_n$.
\end{enumerate}
\end{proposition}




\begin{proof}
The identity in (i) is exactly Lemma~\ref{lem:mass_scaling_wasserstein} with $\alpha=\bar\mu$ and $\beta=\bar\mu_n$.
Because $\overline{\Omega}$ is compact, the empirical measures satisfy $\bar\mu_n \Rightarrow \bar\mu$ almost surely.
Fix $x_0\in\Omega$.
Then $x\mapsto \|x-x_0\|^p$ is bounded and continuous on $\overline{\Omega}$, so along the same event
\begin{align*}
    \int \|x-x_0\|^p d\bar\mu_n(x)
    \to
    \int \|x-x_0\|^p d\bar\mu(x).
\end{align*}
Weak convergence together with convergence of the $p$th moments is equivalent to convergence in $W_p$, hence $W_p(\bar\mu,\bar\mu_n)\to 0$ almost surely.
Applying Lemma~\ref{lem:mass_scaling_wasserstein} once more gives $W_p(\mu,\tilde\mu_n)\to 0$ almost surely.

For (ii), fix $s>d_p^*(\bar\mu)$.
By definition of $d_p^*(\bar\mu)$, the hypothesis of Proposition~5 in \cite{weed2019sharp} is satisfied for the probability measure $\bar\mu$ on the compact metric space $\overline{\Omega}$ after a harmless rescaling to $\mathrm{diam}(\Omega)\le 1$.
Hence
\begin{align*}
    \Ep[W_p^p(\bar\mu,\bar\mu_n)]
    \le
    C_1 n^{-p/s} + C_2 n^{-1/2}.
\end{align*}
Because $s>2p$, the second term is of smaller order and can be absorbed into the first for all sufficiently large $n$.
Using Lemma~\ref{lem:mass_scaling_wasserstein}, we obtain
\begin{align*}
    \Ep[W_p^p(\mu,\tilde\mu_n)]
    &=
    M_\mu \Ep[W_p^p(\bar\mu,\bar\mu_n)]
    \\
    &\lesssim
    M_\mu n^{-p/s},
\end{align*}
which proves \eqref{eq:weighted_empirical_wp_power_rate}.
The displayed bound for $\Ep[W_p(\mu,\tilde\mu_n)]$ follows from Jensen's inequality:
\begin{align*}
    \Ep[W_p(\mu,\tilde\mu_n)]
    \le
    \Ep[W_p^p(\mu,\tilde\mu_n)]^{1/p}
    \lesssim
    M_\mu^{1/p} n^{-1/s}.
\end{align*}

For (iii), let $\bar\nu_n \coloneqq \nu_n/(M_\mu)$.
Then $\bar\nu_n$ is a probability measure supported on at most $n$ points.
The lower-bound part of Theorem~1 of \cite{weed2019sharp}, together with the sentence immediately following that theorem, implies that for every $t<d_*(\bar\mu)$,
\begin{align*}
    W_p(\bar\mu,\bar\nu_n) \gtrsim n^{-1/t}.
\end{align*}
Applying Lemma~\ref{lem:mass_scaling_wasserstein} once again yields
\begin{align*}
    W_p(\mu,\nu_n)
    =
    (M_\mu)^{1/p} W_p(\bar\mu,\bar\nu_n)
    \gtrsim
    (M_\mu)^{1/p} n^{-1/t}.
\end{align*}
Taking $\nu_n=\tilde\mu_n$ proves the last claim.
\end{proof}



\begin{remark}
If one follows the constructive proof of \cite{weed2019sharp} more literally, every partial transport plan in their dyadic argument can indeed be multiplied by $M_\mu$.
Equivalently,
\begin{align}
    W_p^p(\mu,\tilde\mu_n) = (M_\mu) W_p^p(\bar\mu,\bar\mu_n),
\end{align}
so any bound proved at the level of $W_p^p$ for the normalized empirical measure transfers verbatim after multiplying the right-hand side by $M_\mu$.
\end{remark}

\section{Details of the simulation} 

\subsection{Computational resources}\label{sec:comp}
The experiments were run on a Linux computing node with an 8-core Intel(R) Xeon(R) Gold 5222 processor and 755 GB of system memory. The experiments were run on a single NVIDIA Tesla V100-SXM2 GPU, utilizing 32 GB of VRAM.

\subsection{Simulation and implementation details}\label{sec:sim_detail}
In our simulation, the source and target measures $\mu,\nu \in \mM_+([0,1]^d)$ have $1$-Hölder smooth densities $\propto \prod_{i=1}^d \lvert \sin \pi(x_i-c)\rvert$, where $c=0.3$ for $\mu$ and $c=0.7$ for $\nu$. We scale the measures so that $\mu([0,1]^d)=1$ and $\nu([0,1]^d)=2.5$. We generate source and target samples of size $n\in \{100,200,500,1000\}$ over 10 random seeds. For PB-1NN and PB-Kernel, we solve the discrete UOT problem by framing it as non-negative regularized linear regression as described in \cite{Chapel2021}. For PI-kernel and PI-wavelet, we solve the optimization problem in \eqref{eq:kernel_estimator} using the L-BFGS algorithm \cite{Liu1989} over the discretized grid with resolution $128, 64, 32, 32$ for $d=1,2,3,4$, respectively. For the plan-based estimators, we use the UOT solver from the Python Optimal Transport library \cite{flamary2021pot}. For the purpose of learning rate estimation, we first estimate the oracle number of kernel basis elements $L_n$ and wavelet resolution levels $J_n$ via cross-validation with $n=5000$, and then scale $L_n$ and $J_n$ for each $n\in \{100,200,500,1000\}$ according to Theorem~\ref{thm:main_plan_based_rates} and~\ref{thm:main_wavelet_rates}, respectively. The cone program in the SSUOT estimator was solved using the CVXPY library \cite{diamond2016cvxpy} with MOSEK solver \cite{mosek}.

\bibliography{main}
\bibliographystyle{alpha}


\end{document}


